% Limitation des conséquences liées aux échanges d'eau, de substances dissoutes et de particules entre la cellule et le milieu ambiant — SVT Tle D — Cours signé OnBuch+
% Programme officiel MINESEC. Compiler avec : tectonic N01-echanges-cellulaires.tex
\newcommand{\DOCMATIERE}{SVT}
\newcommand{\DOCNIVEAU}{Tle D}
\newcommand{\DOCMODULE}{Module I : Le Monde Vivant — Échanges cellulaires}
\newcommand{\DOCLECON}{N01}
\newcommand{\DOCTITRE}{Limitation des conséquences liées aux échanges d'eau, de substances dissoutes et de particules entre la cellule et le milieu ambiant}
% OnBuch+ — préambule « cours signés » ENRICHI (compilé avec tectonic / XeLaTeX)
% Système visuel « Pop » d'OnBuch+ : fond crème (jamais blanc), bordures encre
% épaisses, ombres « dures » décalées, titres Archivo Black en capitales.
% Version MATHÉMATIQUES : graphes (pgfplots/tikz), boîtes pédagogiques
% (définition, propriété, méthode, attention, le savais-tu, exercice résolu,
% exercice, TP, bilan), page de garde et signature OnBuch+ ancrée en bas.
%
% Macros attendues dans main.tex AVANT \input{preamble} :
%   \DOCMATIERE  \DOCNIVEAU  \DOCMODULE  \DOCLECON (numéro)  \DOCTITRE
\documentclass[11pt]{article}

\usepackage{fontspec}
\usepackage[french]{babel}
\usepackage[a4paper,margin=2cm,top=3.4cm,bottom=2.8cm,headheight=1.4cm,footskip=1.3cm]{geometry}
\usepackage{amsmath,amssymb,amsfonts}
\usepackage{mathtools} % \xmapsto, \coloneqq... souvent utilisés par les modèles
\usepackage{cancel} % simplifications barrées (bilans de forces)
% \abs{z} (module, valeur absolue) et \norm{u} (norme), utilisés par les modèles
\providecommand{\abs}{}\let\abs\relax\DeclarePairedDelimiter{\abs}{\lvert}{\rvert}
\providecommand{\norm}{}\let\norm\relax\DeclarePairedDelimiter{\norm}{\lVert}{\rVert}
\usepackage[table]{xcolor} % [table] : fournit aussi \rowcolors (alternance de lignes)
\usepackage{pagecolor}
\usepackage{tikz}
\usetikzlibrary{arrows.meta,positioning,calc,shapes.geometric,shapes.misc,shapes.arrows,decorations.pathmorphing,decorations.pathreplacing,decorations.markings,patterns,shadows,mindmap,trees,fit,backgrounds,matrix,intersections,angles,quotes}
\usepackage{pgfplots}
\usepackage[version=4]{mhchem} % équations nucléaires \ce{^{235}_{92}U}
\usepackage[european,siunitx]{circuitikz} % schémas électriques
\pgfplotsset{compat=1.18}
\usepgfplotslibrary{fillbetween,groupplots} % aires entre courbes (calcul intégral)
\usepackage{enumitem}
\usepackage{fancyhdr}
% [most] charge toutes les bibliothèques tcolorbox (listings, minted, raster,
% documentation...) et gonfle inutilement la mémoire TeX sur un document long
% (~300 boîtes) : on ne charge que ce qui est réellement utilisé.
\usepackage[skins,breakable]{tcolorbox}
\usepackage{graphicx}
\usepackage{colortbl}
\usepackage{booktabs}
\usepackage{tabularx}
\usepackage{diagbox} % case d'en-tête diagonale des tableaux à double entrée
% Colonnes extensibles centrée / gauche / droite (C, L, R), souvent utilisées par les modèles
\newcolumntype{C}{>{\centering\arraybackslash}X}
\newcolumntype{L}{>{\raggedright\arraybackslash}X}
\newcolumntype{R}{>{\raggedleft\arraybackslash}X}
\usepackage{array}
\usepackage{multicol}
\usepackage{siunitx}
% parse-numbers=false : accepte aussi \SI{2,5 \times 10^{-3}}{...}
\sisetup{locale=FR,output-decimal-marker={,},group-minimum-digits=4,parse-numbers=false}
\DeclareSIUnit\ohm{\ensuremath{\symup{\Omega}}} % U+2126 absent de la police
\DeclareSIUnit\division{div} % réglages d'oscilloscope (ms/div, V/div)
\DeclareSIUnit\tr{tr} % tours (vitesse de rotation, ex. \SI{1500}{\tr\per\minute})
\DeclareSIUnit\molar{M} % concentration molaire (ex. \SI{3}{\molar}, \SI{1}{\micro\molar})
\DeclareSIUnit\an{an} % année (le modèle confond parfois avec \year, un compteur TeX) — ex. \SI{5}{\kilo\meter\per\an}
\DeclareSIUnit\FCFA{FCFA} % franc CFA (ex. \SI{500}{\FCFA\per\kilo\gram})
\DeclareSIUnit\degreeBrix{\textdegree Brix} % teneur en sucre d'un sirop/jus (réfractométrie)
\DeclareSIUnit\month{mois} % le modèle utilise parfois \month, un compteur TeX, comme unité de durée
\usepackage{qrcode}
% \degree : commande fréquemment utilisée par les modèles pour un angle ou une
% température en dehors de \SI{}{} (non fournie par les paquets chargés).
\providecommand{\degree}{^{\circ}}
% \qed (fin de démonstration), utilisé par les modèles sans amsthm
\providecommand{\qed}{\hfill$\square$}
% \barre : double barre des valeurs interdites dans un tableau de variations (tkz-tab)
\providecommand{\barre}{\ensuremath{\|}}
% environnement proof (amsthm non chargé) utilisé par les modèles
\ifdefined\proof\else\newenvironment{proof}[1][Démonstration]{\par\noindent\textit{#1.}\ }{\qed\par}\fi
\usepackage{lastpage}
\usepackage{hyperref}
\hypersetup{hidelinks}

% ============================================================
% Palette « Pop » OnBuch+ — mêmes hex que la classe P (plus_ui.dart)
% ============================================================
\definecolor{poporange}{HTML}{F59321}
\definecolor{popdark}{HTML}{E07A0C}
\definecolor{popcream}{HTML}{FFF6E8}
\definecolor{popink}{HTML}{1C1714}
\definecolor{poppurple}{HTML}{7A5AE0}
\definecolor{popgreen}{HTML}{2E9E6A}
\definecolor{poppink}{HTML}{E0457B}
\definecolor{popblue}{HTML}{2F7DE1}
\definecolor{popgold}{HTML}{C98A12}
\definecolor{popmuted}{HTML}{8A7F76}
\definecolor{popline}{HTML}{EADFD2}
\definecolor{popgray}{HTML}{8A7F76}
\colorlet{popgrey}{popgray}
% Alias tolérants (le modèle nomme parfois les couleurs différemment)
\colorlet{popwhite}{white}
\colorlet{popblack}{popink}
\colorlet{popred}{poppink}
\colorlet{popyellow}{popgold}
% Teintes claires (fonds de tableaux/encadrés) — le modèle nomme parfois
% les couleurs avec un suffixe L (ex. popblueL) sans les avoir définies.
\colorlet{poporangeL}{poporange!20}
\colorlet{popdarkL}{popdark!20}
\colorlet{poppurpleL}{poppurple!20}
\colorlet{popgreenL}{popgreen!20}
\colorlet{poppinkL}{poppink!20}
\colorlet{popblueL}{popblue!20}
\colorlet{popgoldL}{popgold!20}
\colorlet{popredL}{popred!20}
\colorlet{popgrayL}{popgray!20}
% Teintes claires pour fonds de tableaux / zones
\colorlet{popblueL}{popblue!10!white}
\colorlet{poporangeL}{poporange!14!white}
\colorlet{popgreenL}{popgreen!12!white}
\colorlet{poppurpleL}{poppurple!12!white}
\colorlet{poppinkL}{poppink!12!white}
\colorlet{popgoldL}{popgold!14!white}

\pagecolor{popcream}

% ============================================================
% Typographie
% ============================================================
\defaultfontfeatures{Ligatures=TeX}
\setmainfont{PlusJakartaSans-Medium.ttf}[Path=../../../fonts/,BoldFont=PlusJakartaSans-Bold.ttf]
% Maths en sans-serif (Fira Math) avec chiffres et lettres droites de Plus
% Jakarta, pour que les formules se fondent dans le texte.
\usepackage[math-style=ISO,bold-style=ISO]{unicode-math}
\setmathfont{FiraMath-Regular.otf}[Path=../../../fonts/]
\setmathfont{PlusJakartaSans-Medium.ttf}[Path=../../../fonts/,range={up/num,up/{Latin,latin},\mathup}]
\usepackage{icomma} % virgule décimale sans espace en mode math
% Caractères mathématiques Unicode absents des polices de texte (Plus Jakarta,
% Archivo Black) : rendus par leur équivalent mathématique, en texte comme en
% formule — sinon ils s'affichent en carré vide (ex. « Arithmétique dans ℤ »).
\usepackage{newunicodechar}
\newunicodechar{ℝ}{\ensuremath{\mathbb{R}}}
\newunicodechar{ℤ}{\ensuremath{\mathbb{Z}}}
\newunicodechar{ℂ}{\ensuremath{\mathbb{C}}}
\newunicodechar{ℕ}{\ensuremath{\mathbb{N}}}
\newunicodechar{ℚ}{\ensuremath{\mathbb{Q}}}
\newunicodechar{𝔻}{\ensuremath{\mathbb{D}}}
\newunicodechar{α}{\ensuremath{\alpha}}
\newunicodechar{β}{\ensuremath{\beta}}
\newunicodechar{γ}{\ensuremath{\gamma}}
\newunicodechar{δ}{\ensuremath{\delta}}
\newunicodechar{ε}{\ensuremath{\varepsilon}}
\newunicodechar{θ}{\ensuremath{\theta}}
\newunicodechar{λ}{\ensuremath{\lambda}}
\newunicodechar{μ}{\ensuremath{\mu}}
\newunicodechar{σ}{\ensuremath{\sigma}}
\newunicodechar{τ}{\ensuremath{\tau}}
\newunicodechar{φ}{\ensuremath{\varphi}}
\newunicodechar{ψ}{\ensuremath{\psi}}
\newunicodechar{ω}{\ensuremath{\omega}}
\newunicodechar{Σ}{\ensuremath{\Sigma}}
% majuscules produites par \MakeUppercase dans les titres (α-aminés -> Α)
\newunicodechar{Α}{\ensuremath{\alpha}}
\newunicodechar{Β}{\ensuremath{\beta}}
\newunicodechar{Γ}{\ensuremath{\Gamma}}
\newunicodechar{Δ}{\ensuremath{\Delta}}
\newunicodechar{Ε}{\ensuremath{\varepsilon}}
\newunicodechar{Θ}{\ensuremath{\Theta}}
\newunicodechar{Λ}{\ensuremath{\Lambda}}
\newunicodechar{Μ}{\ensuremath{\mu}}
\newunicodechar{Τ}{\ensuremath{\tau}}
\newunicodechar{Φ}{\ensuremath{\Phi}}
\newunicodechar{Ψ}{\ensuremath{\Psi}}
\newunicodechar{Ω}{\ensuremath{\Omega}}
\newunicodechar{↦}{\ensuremath{\mapsto}}
\newunicodechar{∈}{\ensuremath{\in}}
\newunicodechar{∉}{\ensuremath{\notin}}
\newunicodechar{∧}{\ensuremath{\wedge}}
\newunicodechar{⇔}{\ensuremath{\Leftrightarrow}}
\newunicodechar{⟺}{\ensuremath{\Longleftrightarrow}}
\newunicodechar{⇒}{\ensuremath{\Rightarrow}}
\newunicodechar{⟹}{\ensuremath{\Longrightarrow}}
\newunicodechar{∘}{\ensuremath{\circ}}
\newunicodechar{∪}{\ensuremath{\cup}}
\newunicodechar{∩}{\ensuremath{\cap}}
\newunicodechar{≡}{\ensuremath{\equiv}}
\newunicodechar{⊂}{\ensuremath{\subset}}
\newunicodechar{⊥}{\ensuremath{\perp}}
\newunicodechar{∃}{\ensuremath{\exists}}
\newunicodechar{∀}{\ensuremath{\forall}}
\newunicodechar{∛}{\ensuremath{\sqrt[3]{\,}}}
\newunicodechar{∜}{\ensuremath{\sqrt[4]{\,}}}
\newunicodechar{⃗}{\ensuremath{{}^{\to}}}
\newunicodechar{ⁿ}{\ifmmode^{n}\else\textsuperscript{\ensuremath{n}}\fi}
\newunicodechar{ˣ}{\ifmmode^{x}\else\textsuperscript{\ensuremath{x}}\fi}
\newunicodechar{ᵇ}{\ifmmode^{b}\else\textsuperscript{\ensuremath{b}}\fi}
\newunicodechar{ʳ}{\ifmmode^{r}\else\textsuperscript{\ensuremath{r}}\fi}
\newunicodechar{ᵘ}{\ifmmode^{u}\else\textsuperscript{\ensuremath{u}}\fi}
\newunicodechar{ʸ}{\ifmmode^{y}\else\textsuperscript{\ensuremath{y}}\fi}
\newunicodechar{ᵃ}{\ifmmode^{a}\else\textsuperscript{\ensuremath{a}}\fi}
\newunicodechar{ᵉ}{\ifmmode^{e}\else\textsuperscript{\ensuremath{e}}\fi}
\newunicodechar{ᵏ}{\ifmmode^{k}\else\textsuperscript{\ensuremath{k}}\fi}
\newunicodechar{ᵖ}{\ifmmode^{p}\else\textsuperscript{\ensuremath{p}}\fi}
\newunicodechar{ᴺ}{\ifmmode^{N}\else\textsuperscript{\ensuremath{N}}\fi}
\newunicodechar{ᶜ}{\ifmmode^{c}\else\textsuperscript{\ensuremath{c}}\fi}
\newunicodechar{ʰ}{\ifmmode^{h}\else\textsuperscript{\ensuremath{h}}\fi}
\newunicodechar{ᵠ}{\ifmmode^{q}\else\textsuperscript{\ensuremath{q}}\fi}
\newunicodechar{ₙ}{\ifmmode_{n}\else\textsubscript{\ensuremath{n}}\fi}
\newunicodechar{ₐ}{\ifmmode_{a}\else\textsubscript{\ensuremath{a}}\fi}
\newunicodechar{ᵢ}{\ifmmode_{i}\else\textsubscript{\ensuremath{i}}\fi}
\newunicodechar{ᵣ}{\ifmmode_{r}\else\textsubscript{\ensuremath{r}}\fi}
\newunicodechar{ₖ}{\ifmmode_{k}\else\textsubscript{\ensuremath{k}}\fi}
\newunicodechar{ᵦ}{\ifmmode_{\beta}\else\textsubscript{\ensuremath{\beta}}\fi}
\newunicodechar{ₓ}{\ifmmode_{x}\else\textsubscript{\ensuremath{x}}\fi}
\newfontfamily\popdisplay{ArchivoBlack-Regular.ttf}[Path=../../../fonts/]
\newfontfamily\poplogo{SpaceGrotesk-Bold.ttf}[Path=../../../fonts/]
\newfontfamily\popsemi{PlusJakartaSans-SemiBold.ttf}[Path=../../../fonts/]
\newfontfamily\popxbold{PlusJakartaSans-ExtraBold.ttf}[Path=../../../fonts/]
\newfontfamily\popxboldls{PlusJakartaSans-ExtraBold.ttf}[Path=../../../fonts/,LetterSpace=3.0]

\setlist[enumerate]{leftmargin=1.7em,itemsep=5pt,topsep=4pt}
\setlist[itemize]{leftmargin=1.7em,itemsep=4pt,topsep=4pt,label={\color{poporange}\textbullet}}
\setlength{\parindent}{0pt}
\setlength{\parskip}{5pt}
\renewcommand{\arraystretch}{1.25}

% pgfplots : style Pop par défaut
\pgfplotsset{
  every axis/.append style={axis line style={popink,line width=1pt},tick style={popink},
    grid style={popline,line width=0.5pt},label style={font=\small},tick label style={font=\footnotesize}},
  popcurve/.style={poporange,line width=1.8pt,no markers,smooth},
}
% Même style disponible dans un \draw TikZ simple (hors pgfplots)
\tikzset{popcurve/.style={poporange,line width=1.8pt}}
\tikzset{popgrid/.style={popline,very thin}}
\colorlet{popcurve}{poporange} % le nom du style est parfois utilisé comme couleur

% ============================================================
% Pill / squiggle
% ============================================================
\newcommand{\poppill}[3]{%
  \tikz[baseline=-0.55ex]\node[draw=popink,line width=1.2pt,rounded corners=2.6pt,%
    fill=#2,inner xsep=6.5pt,inner ysep=3pt]{%
    \popxboldls\fontsize{7}{7}\selectfont\color{#3}\MakeUppercase{#1}};%
}
\newcommand{\popsquiggle}[1]{%
  \tikz[baseline=-0.4ex]{
    \draw[#1,line width=2.6pt,line cap=round]
      plot [smooth] coordinates {(0,0.09) (0.34,0.01) (0.68,0.21) (1.02,0.01) (1.36,0.10)};
  }%
}
% Mot-clé surligné (marqueur orange sous le texte)
\newcommand{\cle}[1]{{\popxbold\color{popdark}#1}}

% ============================================================
% En-tête / pied
% ============================================================
\pagestyle{fancy}
\fancyhf{}
\renewcommand{\headrulewidth}{0pt}
\renewcommand{\footrulewidth}{0pt}
\lhead{\raisebox{-0.15ex}{\poplogo\fontsize{13}{13}\selectfont\color{popink}OnBuch\color{poporange}+}}
\rhead{\poppill{\DOCMATIERE\ \textbullet\ \DOCNIVEAU\ \textbullet\ Leçon \DOCLECON}{white}{popink}}
\lfoot{\popxbold\fontsize{7.6}{7.6}\selectfont\color{popmuted}\MakeUppercase{Cours signé OnBuch+}\ \textbullet\ {\popsemi\DOCTITRE}}
\rfoot{\tikz[baseline=-0.6ex]\node[rounded corners=8.5pt,fill=popink,minimum height=17pt,minimum width=17pt,inner xsep=5pt,inner sep=0]{\popxbold\fontsize{8}{8}\selectfont\color{popcream}\thepage\,/\,\pageref*{LastPage}};}

% ============================================================
% Titres
% ============================================================
\newcommand{\chaptitre}[2]{%
  \begin{tcolorbox}[enhanced,colback=poporange,colframe=popink,boxrule=2.6pt,arc=11pt,
      drop shadow=popink,left=15pt,right=15pt,top=12pt,bottom=11pt,before skip=3pt,after skip=18pt]
    \poppill{#2}{popink}{popcream}\par\vspace{9pt}
    {\raggedright\hyphenpenalty=10000\exhyphenpenalty=10000\popdisplay\fontsize{22}{24}\selectfont\color{popcream}\MakeUppercase{#1}\par}
  \end{tcolorbox}
}
% Section numérotée : pastille orange avec le numéro + titre Archivo
\newcounter{popsec}
\newcommand{\coursec}[1]{%
  \stepcounter{popsec}\setcounter{popsub}{0}%
  \par\vspace{16pt}\noindent
  \begin{minipage}{\linewidth}
  \tikz[baseline=-0.7ex]\node[draw=popink,line width=1.6pt,fill=poporange,rounded corners=4pt,minimum size=22pt,inner sep=0]{\popdisplay\fontsize{12}{12}\selectfont\color{popcream}\thepopsec};%
  \hspace{8pt}{\popdisplay\fontsize{15}{17}\selectfont\color{popink}\MakeUppercase{#1}}\par
  \vspace{4pt}\noindent{\color{poporange}\rule{36pt}{3.6pt}}
  \end{minipage}\par\nopagebreak\vspace{8pt}
}
\newcounter{popsub}
\newcommand{\courssub}[1]{%
  \stepcounter{popsub}%
  \par\vspace{9pt}\noindent{\popxbold\fontsize{11.5}{13}\selectfont\color{popink}{\color{poporange}\thepopsec.\thepopsub}\hspace{6pt}#1}\par\nopagebreak\vspace{4pt}
}

% ============================================================
% Boîtes pédagogiques — même recette : fond blanc, bordure épaisse colorée,
% ombre dure encre, badge plein. Titre optionnel [#1].
% ============================================================
\tcbset{popbase/.style={breakable,enhanced,colback=white,boxrule=2.3pt,arc=9pt,
  shadow={3pt}{-3pt}{0pt}{popink},left=14pt,right=14pt,top=11pt,bottom=11pt,before skip=11pt,after skip=15pt,
  fontupper=\popsemi\fontsize{10.6}{14.5}\selectfont\color{popink},
  fontlower=\popsemi\fontsize{10.6}{14.5}\selectfont\color{popink}}}
% Coche dessinée (le glyphe ✓ n'existe pas dans les polices du document)
\newcommand{\popcheck}[1]{\tikz[baseline=-0.5ex]\draw[#1,line width=1.6pt,line cap=round,line join=round](-0.13,0.0)--(-0.04,-0.09)--(0.13,0.1);}
\providecommand{\coche}{\popcheck{popgreen}}
\providecommand{\resultat}[1]{\par\smallskip\noindent\fcolorbox{popgreen}{popgreenL}{\parbox{\dimexpr\linewidth-2\fboxsep-2\fboxrule\relax}{\textbf{Résultat.} #1}}\par\smallskip}
\newcommand{\popboxhead}[3]{\poppill{#1}{#2}{white}\ifx\relax#3\relax\else\hspace{6pt}{\popxbold\fontsize{10.6}{12}\selectfont\color{popink}#3}\fi\par\vspace{7pt}}

\newtcolorbox{aretenir}[1][]{popbase,colframe=popgreen,before upper={\popboxhead{À retenir}{popgreen}{#1}}}
\newtcolorbox{exemplebox}[1][]{popbase,colframe=poporange,before upper={\popboxhead{Exemple}{poporange}{#1}}}
\newtcolorbox{definition}[1][]{popbase,colframe=popblue,before upper={\popboxhead{Définition}{popblue}{#1}}}
\newtcolorbox{propriete}[1][]{popbase,colframe=poppurple,before upper={\popboxhead{Propriété}{poppurple}{#1}}}
\newtcolorbox{methode}[1][]{popbase,colframe=popgold,before upper={\popboxhead{Méthode}{popgold}{#1}}}
\newtcolorbox{attention}[1][]{popbase,colframe=poppink,before upper={\popboxhead{Attention}{poppink}{#1}}}
\newtcolorbox{experience}[1][]{popbase,colframe=popblue,colback=popblueL,before upper={\popboxhead{Expérience}{popblue}{#1}}}
% « Le savais-tu ? » : carte violette pleine (contexte, culture, Cameroun)
\newtcolorbox{savaistu}[1][]{popbase,colback=poppurple,colframe=popink,
  fontupper=\popsemi\fontsize{10.4}{14.2}\selectfont\color{white},
  before upper={\poppill{Le savais-tu\,?}{white}{poppurple}\ifx\relax#1\relax\else\hspace{6pt}{\popxbold\color{white}#1}\fi\par\vspace{7pt}}}
% Exercice résolu : énoncé en haut, \tcblower puis la solution sur fond crème
\newcounter{popexo}\newcounter{popexr}
\newtcolorbox{exoresolu}[1][]{popbase,colframe=poporange,colbacklower=popcream,
  segmentation style={popink,line width=1.2pt,dashed},
  before upper={\stepcounter{popexr}\popboxhead{Exercice résolu \thepopexr}{poporange}{#1}},
  before lower={\poppill{Solution}{popink}{popcream}\par\vspace{6pt}}}
% Exercice (non résolu en place) — la correction est dans la partie Corrigés
\newtcolorbox{exercice}[1][]{popbase,colframe=popink,
  before upper={\stepcounter{popexo}\popboxhead{Exercice \thepopexo}{popink}{#1}}}
% Corrigé
\newtcolorbox{corrige}[1][]{popbase,colframe=popgreen,colback=popgreenL,
  before upper={\popboxhead{Corrigé}{popgreen}{#1}}}
% Objectifs / prérequis / bilan
\newtcolorbox{objectifs}{popbase,colframe=popink,colback=poporangeL,
  before upper={\popboxhead{Objectifs de la leçon}{popink}{}}}
\newtcolorbox{prerequis}{popbase,colframe=popink,colback=popblueL,
  before upper={\popboxhead{Prérequis}{popblue}{}}}
\newtcolorbox{bilan}{popbase,colframe=popink,colback=white,boxrule=2.8pt,
  before upper={\popboxhead{Fiche bilan}{poporange}{L'essentiel à connaître pour l'examen}}}
% Figure encadrée avec légende
\newtcolorbox{popfigure}[1][]{enhanced,breakable=false,colback=white,colframe=popline,boxrule=1.4pt,arc=8pt,
  left=10pt,right=10pt,top=10pt,bottom=8pt,before skip=10pt,after skip=12pt,halign=center,
  lower separated=false,halign lower=center,
  fontlower=\popsemi\fontsize{9}{11.5}\selectfont\color{popmuted},
  #1}
\newcommand{\legende}[1]{\par\vspace{6pt}{\popsemi\fontsize{9}{11.5}\selectfont\color{popmuted}#1}}

% ============================================================
% Signature OnBuch+
% ============================================================
\newcommand{\poplogoinline}[1]{{\poplogo\fontsize{#1}{#1}\selectfont\color{popink}OnBuch\color{poporange}+}}
\newcommand{\signatureonbuch}{%
  \begin{tcolorbox}[enhanced,colback=popink,colframe=popink,boxrule=2.2pt,arc=10pt,
      drop shadow=poporange,left=14pt,right=14pt,top=10pt,bottom=10pt,before skip=0pt,after skip=0pt]
    \tikz[baseline=-0.7ex]\node[circle,fill=popgreen,draw=popcream,line width=1.3pt,minimum size=22pt,inner sep=0]{\popcheck{white}};%
    \hspace{9pt}\begin{minipage}[c]{0.62\linewidth}
      {\popxbold\fontsize{12}{13}\selectfont\color{popcream}Signé\ \textbullet\ {\poplogo OnBuch\color{poporange}+}}\par\vspace{2pt}
      {\raggedright\popsemi\fontsize{8}{10.5}\selectfont\color{popline}\MakeUppercase{Équipe pédagogique OnBuch+}\\\MakeUppercase{Conforme au programme officiel MINESEC}\par}
    \end{minipage}\hfill
    {\poplogo\fontsize{22}{22}\selectfont\color{popcream}OnBuch\color{poporange}+}
  \end{tcolorbox}%
}

% ============================================================
% Page de garde
% #1 titre · #2 matière · #3 niveau · #4 module · #5 n° leçon · #6 durée ·
% #7 description / objectifs (texte ou itemize)
% ============================================================
\newcommand{\pagegarde}[7]{%
  \thispagestyle{empty}
  \newgeometry{a4paper,margin=1.7cm,top=1.6cm,bottom=1.4cm}%
  \begin{tikzpicture}[remember picture,overlay]
    \fill[poporange!18] ([xshift=-3.2cm,yshift=-3.6cm]current page.north east) circle (4.6cm);
    \fill[poppurple!14] ([xshift=2.4cm,yshift=7.6cm]current page.south west) circle (3.8cm);
  \end{tikzpicture}%
  {\poplogo\fontsize{30}{30}\selectfont\color{popink}OnBuch\color{poporange}+}\hfill
  \raisebox{0.6ex}{\poppill{Cours signé \textbullet\ Premium}{popink}{popcream}}\par
  \vspace{1pt}\popsquiggle{poppurple}\par
  \vspace{8pt}
  \begin{tcolorbox}[enhanced,colback=poporange,colframe=popink,boxrule=2.8pt,arc=13pt,
      drop shadow=popink,left=18pt,right=18pt,top=12pt,bottom=13pt,after skip=10pt]
    \poppill{#2 \textbullet\ #3}{popink}{popcream}\hspace{5pt}\poppill{#4}{white}{popink}\par\vspace{8pt}
    {\popdisplay\fontsize{14}{14}\selectfont\color{popink}LEÇON #5}\par\vspace{4pt}
    {\raggedright\hyphenpenalty=10000\exhyphenpenalty=10000\popdisplay\fontsize{25}{27}\selectfont\color{popcream}\MakeUppercase{#1}\par}
    \vspace{8pt}
    {\popxbold\fontsize{9.5}{9.5}\selectfont\color{popink}\MakeUppercase{Durée conseillée : #6}}
  \end{tcolorbox}
  \begin{tcolorbox}[enhanced,colback=white,colframe=popink,boxrule=2.2pt,arc=10pt,
      drop shadow=popink,left=16pt,right=16pt,top=10pt,bottom=10pt,after skip=8pt]
    \poppill{Dans cette leçon}{poppurple}{white}\par\vspace{5pt}
    \popsemi\fontsize{9.3}{12.6}\selectfont\color{popink}#7
  \end{tcolorbox}
  \noindent\begin{minipage}[t]{0.487\linewidth}
    \begin{tcolorbox}[enhanced,colback=popgreen,colframe=popink,boxrule=2.2pt,arc=10pt,
        drop shadow=popink,left=11pt,right=11pt,top=10pt,bottom=10pt,height=3.1cm,valign=center]
      \tcbox[colback=white,colframe=popink,boxrule=1.3pt,arc=5pt,boxsep=0pt,nobeforeafter,
        left=3pt,right=3pt,top=3pt,bottom=3pt,valign=center]{\qrcode[height=1.9cm]{https://play.google.com/store/apps/details?id=com.luvvix.onbuch&hl=fr}}%
      \hspace{8pt}\begin{minipage}[c]{0.5\linewidth}
        {\popdisplay\fontsize{11}{12.5}\selectfont\color{white}\MakeUppercase{Télécharge OnBuch}}\par\vspace{4pt}
        {\raggedright\popsemi\fontsize{8.4}{11}\selectfont\color{white}Gratuit \textbullet\ Google Play\\Tuteur IA, annales, concours, résultats\par}
      \end{minipage}
    \end{tcolorbox}
  \end{minipage}\hfill
  \begin{minipage}[t]{0.487\linewidth}
    \begin{tcolorbox}[enhanced,colback=poppurple,colframe=popink,boxrule=2.2pt,arc=10pt,
        drop shadow=popink,left=11pt,right=11pt,top=10pt,bottom=10pt,height=3.1cm,valign=center]
      \tikz[baseline=-0.7ex]\node[circle,fill=white,draw=popink,line width=1.3pt,minimum size=1.6cm,inner sep=0]{\popdisplay\fontsize{24}{24}\selectfont\color{poppurple}+};%
      \hspace{8pt}\begin{minipage}[c]{0.6\linewidth}
        {\popdisplay\fontsize{11}{12.5}\selectfont\color{white}\MakeUppercase{Passe à OnBuch+}}\par\vspace{4pt}
        {\raggedright\popsemi\fontsize{8.4}{11}\selectfont\color{white}Tous les cours signés, Tuteur IA illimité, fiches PDF — onglet OnBuch+ de l'app\par}
      \end{minipage}
    \end{tcolorbox}
  \end{minipage}
  \vspace{8pt}
  \popcontenu
  \vfill
  \signatureonbuch
  \restoregeometry
  \newpage
}

% Rangée « Ce cours contient » (page de garde)
\newcommand{\poptile}[3]{% #1 couleur #2 gros texte #3 légende
  \begin{tcolorbox}[enhanced,colback=#1,colframe=popink,boxrule=2pt,arc=9pt,shadow={2.5pt}{-2.5pt}{0pt}{popink},
      left=6pt,right=6pt,top=9pt,bottom=9pt,halign=center,valign=center,height=1.75cm,width=\linewidth,nobeforeafter]
    {\popdisplay\fontsize{12.5}{13}\selectfont\color{white}\MakeUppercase{#2}}\par\vspace{3pt}
    {\centering\popsemi\fontsize{7.8}{9.5}\selectfont\color{white}#3\par}
  \end{tcolorbox}}
\newcommand{\popcontenu}{%
  \par\noindent{\popxboldls\fontsize{7.5}{7.5}\selectfont\color{popmuted}CE COURS SIGNÉ CONTIENT}\par\vspace{7pt}
  \noindent\begin{minipage}[t]{0.235\linewidth}\poptile{popblue}{Cours}{complet, illustré, pas à pas}\end{minipage}\hfill
  \begin{minipage}[t]{0.235\linewidth}\poptile{poporange}{Méthodes}{et exercices résolus}\end{minipage}\hfill
  \begin{minipage}[t]{0.235\linewidth}\poptile{poppink}{Exercices}{type examen + corrigés}\end{minipage}\hfill
  \begin{minipage}[t]{0.235\linewidth}\poptile{popgreen}{Fiche bilan}{l'essentiel à retenir}\end{minipage}\par}

% Sommaire
\newtcolorbox{sommaire}{enhanced,colback=white,colframe=popink,boxrule=2.2pt,arc=10pt,
  drop shadow=popink,left=18pt,right=18pt,top=13pt,bottom=13pt,before skip=6pt,after skip=20pt,
  before upper={\poppill{Sommaire}{poppurple}{white}\par\vspace{9pt}\popsemi\fontsize{10.6}{14.5}\selectfont\color{popink}}}

% Signature de fin de cours — poussée tout en bas de la dernière page
\newcommand{\findecours}{%
  \par\vfill\nopagebreak
  \begin{center}\popsquiggle{poporange}\end{center}\vspace{4pt}
  \signatureonbuch
}
\AtBeginDocument{\def\llbracket{\mathopen{[\mkern-3.5mu[}}\def\rrbracket{\mathclose{]\mkern-3.5mu]}}} % intervalles d'entiers (glyphe ⟦ absent de la police)
% Glyphes absents de Fira Math : redéfinis à partir de glyphes disponibles
\AtBeginDocument{%
  \def\setminus{\mathbin{\backslash}}\let\smallsetminus\setminus
  \def\vdots{\mathord{\vbox{\baselineskip3.2pt\lineskiplimit0pt\kern4pt\hbox{.}\hbox{.}\hbox{.}}}}%
  \def\ddots{\mathinner{\mkern1mu\raise6.5pt\hbox{.}\mkern2mu\raise3.6pt\hbox{.}\mkern2mu\raise0.7pt\hbox{.}\mkern1mu}}%
  \def\triangle{\mathord{\scalebox{0.9}{$\symup{\Delta}$}}}%
  \def\nabla{\mathord{\raisebox{\depth}{\scalebox{1}[-1]{$\symup{\Delta}$}}}}%
  \def\checkmark{\popcheck{popdark}}%
  \def\star{\mathbin{\ast}}\def\bigstar{\mathord{\ast}}%
  \def\diamond{\mathbin{\scalebox{0.6}{$\Diamond$}}}%
  \def\bigcup{\mathop{\mathchoice{\raisebox{-0.3em}{\scalebox{1.7}{$\cup$}}}{\scalebox{1.25}{$\cup$}}{\cup}{\cup}}\displaylimits}%
  \def\bigcap{\mathop{\mathchoice{\raisebox{-0.3em}{\scalebox{1.7}{$\cap$}}}{\scalebox{1.25}{$\cap$}}{\cap}{\cap}}\displaylimits}%
  \def\lesseqqgtr{\mathrel{\lessgtr}}\def\bigoplus{\mathop{\oplus}\displaylimits}%
}
\providecommand{\ssi}{si et seulement si} % abréviation fréquente
\AtBeginDocument{\providecommand{\wideparen}{\overparen}} % arc de cercle

%% FIGFIX-BEGIN v5 — correctifs globaux des figures (docs/onbuch-generation/tools/figfix)
\usepackage{adjustbox}\usepackage{etoolbox}\tcbuselibrary{raster}
% toute figure TikZ/pgfplots plus large que la ligne est réduite à la largeur disponible
\BeforeBeginEnvironment{tikzpicture}{\begin{adjustbox}{max width=\linewidth}}
\AfterEndEnvironment{tikzpicture}{\end{adjustbox}}
% légendes pgfplots lisibles par-dessus les courbes
\pgfplotsset{every axis/.append style={legend style={fill=white,fill opacity=0.92,text opacity=1,draw=black!15,inner sep=3pt}}}
% étiquettes d'arêtes (syntaxe edge["..."]) avec fond blanc
\tikzset{every edge quotes/.append style={fill=white,inner sep=1.2pt}}
%% FIGFIX-END
\begin{document}

\pagegarde{Limitation des conséquences liées aux échanges d'eau, de substances dissoutes et de particules entre la cellule et le milieu ambiant}{SVT}{Tle D}{Module I : Le Monde Vivant — Échanges cellulaires}{N01}{12 h}{Cette leçon étudie les mécanismes d'échanges entre la cellule et le milieu ambiant : osmose, diffusion simple et facilitée, transport actif, endocytose et exocytose. À partir d'expériences historiques (Dutrochet, Pfeiffer) et d'observations microscopiques, elle explique la turgescence, la plasmolyse, l'hémolyse et la crénelure, ainsi que leurs applications dans la vie courante et la médecine.
\vspace{4pt}
\begin{multicols}{2}\small
\begin{itemize}
\item À la fin de cette leçon, je sais définir l'osmose, la diffusion, le transport actif, l'endocytose et l'exocytose
\item À la fin de cette leçon, je sais identifier au microscope ou sur des électronographies les états de turgescence, de plasmolyse, d'hémolyse et de crénelure des cellules
\item À la fin de cette leçon, je sais interpréter les résultats des expériences de Dutrochet et de Pfeiffer
\item À la fin de cette leçon, je sais concevoir un protocole expérimental mettant en évidence le phénomène d'osmose
\item À la fin de cette leçon, je sais réaliser et interpréter le graphe de la vitesse de transport en fonction de la concentration pour distinguer diffusion simple et diffusion facilitée
\item À la fin de cette leçon, je sais distinguer transport passif et transport actif à partir de critères expérimentaux (sens du gradient, énergie, saturation)
\end{itemize}\end{multicols}}

\begin{sommaire}
\begin{multicols}{2}
\begin{enumerate}
\item Échanges d'eau : observations expérimentales
\item Interprétation des échanges d'eau : osmose et dialyse
\item Échanges de substances dissoutes : transport passif
\item Échanges de substances dissoutes : transport actif
\item Échanges de particules : endocytose et exocytose
\item Importance et applications des échanges cellulaires
\item Activité d'intégration
\item Méthodes et exercices résolus
\item Exercices
\item Corrigés des exercices
\item Fiche bilan
\end{enumerate}
\end{multicols}
\end{sommaire}

\begin{objectifs}
\begin{itemize}
\item À la fin de cette leçon, je sais définir l'osmose, la diffusion, le transport actif, l'endocytose et l'exocytose
\item À la fin de cette leçon, je sais identifier au microscope ou sur des électronographies les états de turgescence, de plasmolyse, d'hémolyse et de crénelure des cellules
\item À la fin de cette leçon, je sais interpréter les résultats des expériences de Dutrochet et de Pfeiffer
\item À la fin de cette leçon, je sais concevoir un protocole expérimental mettant en évidence le phénomène d'osmose
\item À la fin de cette leçon, je sais réaliser et interpréter le graphe de la vitesse de transport en fonction de la concentration pour distinguer diffusion simple et diffusion facilitée
\item À la fin de cette leçon, je sais distinguer transport passif et transport actif à partir de critères expérimentaux (sens du gradient, énergie, saturation)
\item À la fin de cette leçon, je sais décrire les mécanismes de la phagocytose, de la pinocytose et de l'exocytose
\item À la fin de cette leçon, je sais expliquer l'importance des échanges cellulaires dans la vie des végétaux et des animaux et citer des applications concrètes (sérum physiologique, conservation des aliments, réhydratation)
\end{itemize}
\end{objectifs}

\begin{prerequis}
\begin{itemize}
\item Structure de la membrane plasmique (modèle mosaïque fluide) et de la paroi végétale
\item Organisation de la cellule végétale (vacuole, paroi, chloroplastes) et de la cellule animale
\item Notions de solution, de concentration et de gradient de concentration
\item Notion de pression et de mouvement brownien (diffusion vue en physique-chimie)
\item Rôle de l'ATP comme source d'énergie cellulaire (respiration cellulaire, Première)
\end{itemize}
\end{prerequis}

\newpage

\coursec{Échanges d'eau : observations expérimentales}

Au marché de Mokolo à Yaoundé, une vendeuse asperge régulièrement d'eau ses feuilles de ndolé flétries : en quelques minutes, elles redeviennent fermes et croquantes. À l'hôpital de district, un infirmier refuse d'injecter au malade de l'eau pure en perfusion : il utilise toujours du sérum physiologique. Pourquoi l'eau fait-elle revivre les légumes mais tuerait-elle les globules rouges ? Et comment les racines du manioc absorbent-elles les sels minéraux du sol même quand ceux-ci y sont très dilués ? Cette leçon explique les mécanismes par lesquels la cellule échange eau, substances dissoutes et particules avec son milieu, et comment elle en limite les conséquences parfois dramatiques.

Pour répondre à ces questions, commençons par le plus simple : observons ce qui arrive à des cellules végétales et animales placées dans des solutions de concentrations différentes. C'est l'objet de cette première section.

\courssub{La membrane plasmique : frontière vivante de la cellule}

Toute cellule est délimitée par une \cle{membrane plasmique} (ou membrane cytoplasmique), une enveloppe vivante, très fine (environ \SI{7,5}{nm} d'épaisseur, invisible au microscope optique) et \cle{élastique}. Elle est constituée d'une double couche de phospholipides dans laquelle sont insérées des protéines. Cette membrane contrôle tous les échanges entre le milieu intracellulaire et le milieu extracellulaire : c'est une membrane dite \cle{semi-perméable} (ou hémiperméable), car elle laisse passer librement certaines molécules (l'eau) et en bloque d'autres (la plupart des solutés).

Chez les \textbf{végétaux}, la membrane plasmique est doublée, à l'extérieur, d'une \cle{paroi pecto-cellulosique} rigide et épaisse, constituée de cellulose et de pectine. Cette paroi est \textbf{perméable} (elle laisse tout traverser) mais elle est \textbf{inextensible} : elle ne peut pas se dilater. Cette différence structurale entre cellule animale et cellule végétale est la clé de toute la leçon.

\begin{attention}[Deux enveloppes à ne pas confondre]
Chez la cellule végétale, il faut distinguer la \textbf{paroi pecto-cellulosique} (externe, rigide, perméable, non vivante) de la \textbf{membrane plasmique} (interne, fine, élastique, semi-perméable, vivante). Au microscope optique, on voit la paroi ; la membrane n'est visible que lorsqu'elle se décolle de la paroi (plasmolyse). La cellule animale ne possède que la membrane plasmique.
\end{attention}

\courssub{Comportement des cellules végétales dans des solutions de concentrations différentes}

\begin{experience}[Cellules végétales en milieux de concentrations variées]
\textbf{Matériel :} lamelles d'épiderme de bulbe d'oignon (ou feuilles d'Élodée, plante aquatique), microscope optique, lames et lamelles, eau distillée, solution de NaCl à \SI{9}{\per mille}, solution concentrée de NaCl (par exemple \SI{100}{g.L^{-1}}) ou solution concentrée de saccharose.

\textbf{Protocole :} on dépose trois lamelles d'épiderme d'oignon sur trois lames différentes ; on ajoute sur la première une goutte d'eau distillée, sur la deuxième une goutte de solution de NaCl à \SI{9}{\per mille}, sur la troisième une goutte de solution concentrée de NaCl. On observe au microscope optique (grossissement $\times 400$) au bout de quelques minutes.

\textbf{Observations :}
\begin{itemize}
\item \textbf{Dans l'eau distillée :} les cellules sont gonflées, la vacuole est très volumineuse et occupe presque toute la cellule ; le cytoplasme et le noyau sont plaqués contre la paroi. On dit que la cellule est \cle{turgescente}.
\item \textbf{Dans la solution de NaCl à \SI{9}{\per mille} :} les cellules ont leur aspect habituel : vacuole de taille moyenne, protoplasme appliqué contre la paroi sans pression excessive. C'est l'\textbf{état normal}.
\item \textbf{Dans la solution concentrée :} le protoplasme (membrane plasmique + cytoplasme + noyau + vacuole) se rétracte et se décolle de la paroi, formant une masse arrondie au centre de la cellule ; un espace rempli de la solution extérieure apparaît entre la paroi et la membrane. C'est la \cle{plasmolyse}.
\end{itemize}

\textbf{Interprétation :} l'eau est entrée dans la cellule placée en eau distillée (la vacuole gonfle), et elle est sortie de la cellule placée en solution concentrée (la vacuole perd son eau, le protoplasme rétrécit). La paroi rigide, elle, ne bouge pas : c'est pourquoi la cellule végétale ne peut ni éclater en milieu dilué, ni se ratatiner complètement en milieu concentré.
\end{experience}

\begin{definition}[Turgescence et plasmolyse]
La \cle{turgescence} est l'état d'une cellule végétale gonflée par entrée d'eau : la vacuole est volumineuse et le protoplasme est fortement plaqué contre la paroi pecto-cellulosique, qui s'oppose à l'éclatement. La pression exercée par le contenu cellulaire sur la paroi est appelée \cle{pression de turgescence}.

La \cle{plasmolyse} est l'état d'une cellule végétale qui a perdu de l'eau : le protoplasme se rétracte et se décolle de la paroi. La cellule plasmolysée reste \textbf{vivante} : replacée dans l'eau distillée, elle redevient turgescente (déplasmolyse).
\end{definition}

\begin{savaistu}[La turgescence, moteur de la fermeté des plantes]
C'est la turgescence qui maintient dressées les plantes non ligneuses : une salade ou une feuille de ndolé flétrie est un tissu dont les cellules ont perdu leur eau ; trempée dans l'eau, elle redevient ferme car ses cellules se turgescent. À l'inverse, saler des légumes (ou un escargot !) provoque la sortie d'eau des cellules : c'est le principe de la conservation des aliments par le sel ou le sucre, qui déshydrate aussi les micro-organismes et bloque leur développement.
\end{savaistu}

\courssub{Comportement des globules rouges dans des solutions de concentrations différentes}

\begin{experience}[Hématies en milieux de concentrations variées]
\textbf{Matériel :} sang frais défibriné ou citraté, microscope optique, eau distillée, sérum physiologique (solution de NaCl à \SI{9}{\per mille}), solution concentrée de NaCl.

\textbf{Protocole :} on réalise trois gouttes de sang dilué, chacune additionnée de l'une des trois solutions, et on observe rapidement au microscope optique.

\textbf{Observations :}
\begin{itemize}
\item \textbf{Dans l'eau distillée :} les globules rouges gonflent, deviennent sphériques puis \textbf{éclatent} : ils libèrent leur hémoglobine qui colore le milieu en rouge. On ne voit plus que des « fantômes » (membranes vides). C'est l'\cle{hémolyse}.
\item \textbf{Dans le sérum physiologique :} les hématies conservent leur forme normale de disque biconcave (environ \SI{7,5}{\micro m} de diamètre). C'est l'\textbf{état normal}.
\item \textbf{Dans la solution concentrée :} les globules rouges se ratatinent, leur surface devient irrégulière, hérissée de spicules : ce sont des hématies \cle{crénelées}.
\end{itemize}

\textbf{Interprétation :} comme pour la cellule végétale, l'eau entre dans l'hématie placée en eau distillée et sort de l'hématie placée en solution concentrée. Mais la cellule animale, dépourvue de paroi rigide, n'a aucune protection : sa membrane élastique seule ne résiste pas à la pression, et la cellule éclate.
\end{experience}

\begin{definition}[Hémolyse et crénelure]
L'\cle{hémolyse} est l'éclatement des globules rouges (hématies) placés dans un milieu très dilué, avec libération de l'hémoglobine dans le milieu. La cellule est \textbf{détruite} : le phénomène est irréversible.

La \cle{crénelure} est la déformation des hématies placées dans un milieu très concentré : elles perdent leur eau, se ratatinent et présentent des spicules à leur surface.
\end{definition}

\courssub{Le vocabulaire des solutions : hypotonique, isotonique, hypertonique}

Pour décrire ces phénomènes, on compare toujours la concentration de la solution extérieure à celle du milieu intracellulaire.

\begin{definition}[Tonicité d'une solution]
Une solution est dite :
\begin{itemize}
\item \cle{hypotonique} si sa concentration en substances dissoutes est \textbf{inférieure} à celle du contenu cellulaire : l'eau \textbf{entre} dans la cellule (turgescence ou hémolyse) ;
\item \cle{isotonique} si sa concentration est \textbf{égale} à celle du contenu cellulaire : les échanges d'eau se compensent, la cellule garde son état normal ;
\item \cle{hypertonique} si sa concentration est \textbf{supérieure} à celle du contenu cellulaire : l'eau \textbf{sort} de la cellule (plasmolyse ou crénelure).
\end{itemize}
\end{definition}

\begin{attention}[La tonicité est toujours relative]
On ne dit jamais « une solution hypertonique » tout court : une solution est hypertonique \textbf{par rapport à} un contenu cellulaire de référence. L'eau distillée est hypotonique par rapport à toute cellule ; le sérum physiologique à \SI{9}{\per mille} est isotonique par rapport au plasma sanguin, mais hypotonique par rapport à une cellule végétale très concentrée.
\end{attention}

\courssub{Bilan comparatif : pourquoi la cellule végétale résiste-t-elle ?}

\begin{center}
\begin{tabularx}{\linewidth}{|l|X|X|X|}
\hline
\rowcolor{popblueL} \textbf{Milieu} & \textbf{Hypotonique} & \textbf{Isotonique} & \textbf{Hypertonique} \\
\hline
\textbf{Mouvement d'eau} & Entrée d'eau & Équilibre & Sortie d'eau \\
\hline
\textbf{Cellule végétale} & Turgescence (paroi rigide empêche l'éclatement) & État normal & Plasmolyse (cellule vivante, réversible) \\
\hline
\textbf{Globule rouge} & Hémolyse (éclatement, destruction) & État normal & Crénelure (ratatinement) \\
\hline
\end{tabularx}
\end{center}

La \textbf{paroi pecto-cellulosique} explique la différence fondamentale : inextensible, elle s'oppose au gonflement et la cellule végétale développe une pression de turgescence qui finit par équilibrer l'entrée d'eau. La cellule animale, entourée de sa seule membrane élastique, éclate dès que l'entrée d'eau est excessive. La cellule végétale est ainsi protégée contre les variations de concentration du milieu ; la cellule animale ne survit que si le milieu extracellulaire (le plasma, chez nous) est maintenu à concentration constante.

\begin{popfigure}
\begin{tikzpicture}[scale=0.92, every node/.style={font=\small}]
% Titres colonnes
\node[font=\bfseries] at (2.2,5.6) {Milieu hypotonique};
\node[font=\bfseries] at (6.6,5.6) {Milieu isotonique};
\node[font=\bfseries] at (11,5.6) {Milieu hypertonique};
% Ligne 1 : cellule végétale
\node[font=\bfseries, rotate=90] at (-1.6,3.6) {Cellule végétale};
% Hypotonique : turgescente
\draw[thick, popgreen] (0.4,2.4) rectangle (4,4.8);
\draw[thick, popgreen!60!black, fill=popgreenL] (0.55,2.55) rectangle (3.85,4.65);
\draw[fill=popblue!25] (0.7,2.7) rectangle (3.7,4.5);
\node at (2.2,3.6) {vacuole};
\node[align=center, font=\footnotesize] at (2.2,1.9) {Turgescence : membrane\\ plaquée contre la paroi};
% Isotonique : normale
\draw[thick, popgreen] (4.8,2.4) rectangle (8.4,4.8);
\draw[thick, popgreen!60!black, fill=popgreenL] (5.0,2.6) rectangle (8.2,4.6);
\draw[fill=popblue!15] (5.3,2.9) rectangle (7.9,4.3);
\node at (6.6,3.6) {vacuole};
\node[align=center, font=\footnotesize] at (6.6,1.9) {État normal};
% Hypertonique : plasmolysée
\draw[thick, popgreen] (9.2,2.4) rectangle (12.8,4.8);
\draw[thick, poporange, fill=poporangeL] (10.0,2.9) .. controls (9.7,3.6) and (10.2,4.4) .. (11,4.4) .. controls (12.1,4.4) and (12.3,3.6) .. (12.1,3.1) .. controls (11.9,2.7) and (10.5,2.6) .. (10.0,2.9) -- cycle;
\draw[fill=popblue!20] (10.3,3.1) ellipse (0.75 and 0.5);
\node[font=\footnotesize] at (10.3,3.1) {vacuole};
\node[align=center, font=\footnotesize] at (11,1.9) {Plasmolyse : protoplasme\\ décollé de la paroi};
% Ligne 2 : globule rouge
\node[font=\bfseries, rotate=90] at (-1.6,-0.4) {Globule rouge};
% Hypotonique : hémolyse
\draw[thick, poppink, dashed] (1.6,-1.1) circle (0.55);
\draw[thick, poppink, dashed] (2.9,-0.2) circle (0.5);
\draw[poppink, fill=poppink!20] (0.6,0.4) .. controls (1.5,0.9) and (3,0.8) .. (3.9,0.3) .. controls (3.2,0.1) and (1.4,0.1) .. (0.6,0.4) -- cycle;
\node[align=center, font=\footnotesize] at (2.2,-1.9) {Hémolyse : éclatement,\\ libération d'hémoglobine};
% Isotonique : normal
\draw[thick, poppink, fill=poppinkL] (6.6,-0.4) circle (0.7);
\draw[thick, poppink, fill=white] (6.6,-0.4) circle (0.35);
\node[align=center, font=\footnotesize] at (6.6,-1.9) {État normal\\ (disque biconcave)};
% Hypertonique : crénelée
\draw[thick, poppink, fill=poppinkL] (11,-0.4)
  \foreach \a in {0,30,...,330}{ -- ++(\a:0.55) -- ++(\a+15:0.75) } -- cycle;
\node[align=center, font=\footnotesize] at (11,-1.9) {Crénelure : hématie\\ ratatinée, spiculée};
\end{tikzpicture}
\legende{Comportement comparé d'une cellule végétale (ligne du haut) et d'un globule rouge (ligne du bas) placés dans des milieux hypotonique, isotonique et hypertonique. En vert : paroi pecto-cellulosique ; en bleu : vacuole ; en orange : protoplasme rétracté ; en rouge : hématies et hémoglobine libérée.}
\end{popfigure}

\begin{methode}[Identifier au microscope l'état d'une cellule placée dans une solution inconnue]
\begin{enumerate}
\item \textbf{Identifier le type cellulaire} : présence d'une paroi et d'une vacuole $\Rightarrow$ cellule végétale ; disques biconcaves sans noyau $\Rightarrow$ hématies.
\item \textbf{Pour une cellule végétale}, observer la vacuole et la position du protoplasme :
\begin{itemize}
\item vacuole très volumineuse, protoplasme plaqué contre la paroi $\Rightarrow$ \textbf{turgescence} $\Rightarrow$ solution \textbf{hypotonique} ;
\item aspect habituel $\Rightarrow$ état normal $\Rightarrow$ solution \textbf{isotonique} ;
\item protoplasme rétracté, décollé de la paroi $\Rightarrow$ \textbf{plasmolyse} $\Rightarrow$ solution \textbf{hypertonique}.
\end{itemize}
\item \textbf{Pour des hématies}, observer leur forme :
\begin{itemize}
\item cellules éclatées, milieu rosé, « fantômes » $\Rightarrow$ \textbf{hémolyse} $\Rightarrow$ solution \textbf{hypotonique} ;
\item disques réguliers $\Rightarrow$ solution \textbf{isotonique} ;
\item cellules ratatinées, spiculées $\Rightarrow$ \textbf{crénelure} $\Rightarrow$ solution \textbf{hypertonique}.
\end{itemize}
\item \textbf{Conclure} en précisant le sens du mouvement d'eau (entrée ou sortie) qui explique l'état observé.
\end{enumerate}
\end{methode}

\begin{attention}[Plasmolyse $\neq$ hémolyse]
Erreur classique au Baccalauréat : confondre ces deux termes. La \textbf{plasmolyse} concerne la cellule \textbf{végétale} : le protoplasme se rétracte mais la cellule reste \textbf{vivante} (phénomène réversible). L'\textbf{hémolyse} concerne le \textbf{globule rouge} : la cellule \textbf{éclate} et est \textbf{détruite} (phénomène irréversible). De même, ne dites jamais qu'une cellule végétale « éclate » en eau distillée : sa paroi rigide le lui interdit.
\end{attention}

\begin{savaistu}[Pourquoi le sérum physiologique à \SI{9}{\per mille} ?]
Le plasma sanguin a une concentration en sels équivalente à celle d'une solution de NaCl à \SI{9}{g.L^{-1}} (soit \SI{9}{\per mille}). Le \cle{sérum physiologique} est donc \textbf{isotonique} au plasma : perfusé à un malade, il ne provoque ni entrée ni sortie nette d'eau dans les globules rouges, qui restent intacts. Une perfusion d'eau pure, hypotonique, provoquerait une hémolyse massive, potentiellement mortelle : c'est pourquoi l'infirmier de l'hôpital de district a raison de la refuser. À l'inverse, les feuilles de ndolé de la vendeuse de Mokolo, aspergées d'eau (milieu hypotonique), voient leurs cellules se turgescer : d'où leur fermeté retrouvée.
\end{savaistu}

\begin{exoresolu}[Exploiter des observations microscopiques]
\textbf{Énoncé.} Trois lots de lamelles d'épiderme d'oignon sont placés respectivement dans trois solutions A, B et C de concentrations inconnues. Après 10 minutes, on observe : dans A, des cellules au protoplasme décollé de la paroi ; dans B, des cellules à la vacuole très volumineuse ; dans C, des cellules d'aspect normal. (1) Nommer l'état des cellules dans chaque solution. (2) Classer les solutions A, B, C par concentration croissante en justifiant. (3) Que se passe-t-il si l'on replace les cellules du lot A dans l'eau distillée ?
\tcblower
\textbf{Solution.} (1) Dans A, le protoplasme est décollé de la paroi : c'est la \textbf{plasmolyse}. Dans B, la vacuole est très volumineuse et le protoplasme est plaqué contre la paroi : c'est la \textbf{turgescence}. Dans C, les cellules sont à l'\textbf{état normal}.

(2) La plasmolyse traduit une sortie d'eau : A est \textbf{hypertonique} par rapport au contenu cellulaire. La turgescence traduit une entrée d'eau : B est \textbf{hypotonique}. L'état normal correspond à une solution \textbf{isotonique}. D'où le classement par concentration croissante : B $<$ C $<$ A.

(3) Replacées dans l'eau distillée (milieu hypotonique), les cellules du lot A, restées vivantes, regagnent de l'eau : la vacuole regonfle et le protoplasme se réapplique contre la paroi. C'est la \textbf{déplasmolyse}, qui prouve que la plasmolyse est réversible, contrairement à l'hémolyse.
\end{exoresolu}

\begin{exercice}[À toi de jouer]
On place du sang dans trois tubes contenant respectivement de l'eau distillée (tube 1), du sérum physiologique (tube 2) et une solution de NaCl à \SI{50}{g.L^{-1}} (tube 3). Après centrifugation, le liquide surnageant du tube 1 est rouge, celui des tubes 2 et 3 est incolore. (1) Interpréter la couleur du surnageant du tube 1. (2) Décrire l'aspect des hématies des tubes 2 et 3 au microscope. (3) Expliquer pourquoi le tube 2 correspond au sérum physiologique.
\end{exercice}

\begin{corrige}[Exercice]
(1) Le surnageant rouge du tube 1 traduit la présence d'\textbf{hémoglobine libre} : les hématies ont éclaté (\textbf{hémolyse}) car l'eau distillée, hypotonique, est entrée massivement dans les cellules dépourvues de paroi.

(2) Dans le tube 2 (isotonique), les hématies ont leur forme normale de disque biconcave. Dans le tube 3 (hypertonique), elles ont perdu de l'eau : elles sont \textbf{crénelées}, ratatinées et spiculées.

(3) Dans le tube 2, il n'y a ni hémolyse ni crénelure : les échanges d'eau sont équilibrés, ce qui signifie que la solution est \textbf{isotonique} au plasma, comme le sérum physiologique à \SI{9}{\per mille}.
\end{corrige}

\begin{aretenir}[L'essentiel de la section]
\begin{itemize}
\item La membrane plasmique est fine, élastique et semi-perméable ; chez le végétal, elle est protégée par une paroi pecto-cellulosique rigide et inextensible.
\item Milieu \textbf{hypotonique} $\Rightarrow$ entrée d'eau $\Rightarrow$ \textbf{turgescence} (végétal) ou \textbf{hémolyse} (animal).
\item Milieu \textbf{isotonique} $\Rightarrow$ équilibre des échanges $\Rightarrow$ état normal.
\item Milieu \textbf{hypertonique} $\Rightarrow$ sortie d'eau $\Rightarrow$ \textbf{plasmolyse} (végétal, réversible) ou \textbf{crénelure} (animal).
\item La paroi rigide explique que la cellule végétale résiste (pression de turgescence) là où la cellule animale éclate.
\end{itemize}
\end{aretenir}

Ces observations posent une question fondamentale : par quel mécanisme l'eau traverse-t-elle la membrane, toujours dans le sens qui dilue le milieu le plus concentré ? C'est ce que nous allons découvrir dans la section suivante, en interprétant les expériences historiques de Dutrochet et de Pfeffer : le phénomène d'\textbf{osmose}.

\coursec{Interprétation des échanges d'eau : osmose et dialyse}

Dans la section précédente, nous avons \emph{observé} les conséquences visibles des échanges d'eau : la cellule végétale devient turgescente ou se plasmolyse, le globule rouge éclate (hémolyse) ou se crénelle. Mais \emph{pourquoi} l'eau se déplace-t-elle dans un sens ou dans l autre ? Quelle est la force motrice ? C'est à ces questions que répondent les expériences fondatrices de Dutrochet et de Pfeffer, qui nous mèneront aux concepts rigoureux d'\textbf{osmose} et de \textbf{dialyse}.

\courssub{L'expérience de Dutrochet : la première mise en évidence de l'osmose}

\begin{experience}[Expérience de Dutrochet (1826) — Naissance du concept d'osmose]
\textbf{Matériel :} Un cæcum de porc (membrane animale naturelle, poreuse), une solution concentrée de sucre (ou d'alcool), de l'eau pure, un récipient.
\textbf{Protocole :}
\begin{enumerate}
    \item On remplit le cæcum de porc avec la solution concentrée de sucre.
    \item On ferme hermétiquement l'orifice du cæcum (par ligature).
    \item On plonge l'ensemble dans un grand volume d'eau pure.
    \item On observe l'évolution du système au cours du temps.
\end{enumerate}
\textbf{Observations :}
\begin{itemize}
    \item Au début : le cæcum est mou, la solution est à l'intérieur.
    \item Après un certain temps : le cæcum se gonfle, devient dur et tendu. Il peut même éclater sous la pression.
    \item L'eau pure du récipient a traversé la membrane pour entrer dans la poche.
\end{itemize}
\textbf{Interprétation :} L'eau traverse la membrane du milieu le moins concentré en solutés (eau pure) vers le milieu le plus concentré (solution de sucre). La membrane laisse passer l'eau mais retient les molécules de sucre : elle est \textbf{semi-perméable} (ou hémiperméable). Ce mouvement d'eau crée une pression à l'intérieur de la poche : la \textbf{pression osmotique}.
\tcblower
\textbf{Schéma récapitulatif :}
\begin{popfigure}
\begin{tikzpicture}[scale=0.9, every node/.style={font=\small}]
    % Avant
    \begin{scope}[xshift=-4cm]
        \draw[popdark, thick] (-1.5,-2) -- (-1.5,2) -- (1.5,2) -- (1.5,-2) -- cycle; % Becher
        \draw[popblue, thick] (-1.5,0.5) -- (1.5,0.5); % Niveau eau
        \node[popblue] at (0,1.5) {Eau pure};
        
        \draw[poporange, thick, rounded corners=3pt] (-0.6,-1.5) .. controls (-0.6,-1.8) and (0.6,-1.8) .. (0.6,-1.5) .. controls (0.8,-1.5) and (0.8,1.5) .. (0.6,1.5) .. controls (0.6,1.8) and (-0.6,1.8) .. (-0.6,1.5) -- cycle; % Cecum
        \fill[poporange!30] (-0.55,-1.4) .. controls (-0.55,-1.7) and (0.55,-1.7) .. (0.55,-1.4) .. controls (0.55,1.4) and (-0.55,1.4) .. (-0.55,-1.4); % Solution
        \node[poporange, align=center] at (0,-0.5){Solution\\concentrée};
        \draw[->, thick, popgreen] (0,1.6) -- (0,1.0) node[midway, right] {?};
        \node at (0,-2.8) {\textbf{Avant}};
    \end{scope}
    
    % Flèche temps
    \draw[->, thick, popdark] (1.5,0) -- (2.5,0) node[midway, above] {Temps};
    
    % Après
    \begin{scope}[xshift=6.5cm]
        \draw[popdark, thick] (-1.5,-2) -- (-1.5,2) -- (1.5,2) -- (1.5,-2) -- cycle;
        \draw[popblue, thick] (-1.5,0.2) -- (1.5,0.2); % Niveau eau baissé légèrement
        \node[popblue] at (0,1.5) {Eau pure};
        
        \draw[poporange, thick, rounded corners=3pt] (-0.8,-1.8) .. controls (-0.8,-2.1) and (0.8,-2.1) .. (0.8,-1.8) .. controls (1.0,-1.8) and (1.0,1.8) .. (0.8,1.8) .. controls (0.8,2.1) and (-0.8,2.1) .. (-0.8,1.8) -- cycle; % Cecum gonflé
        \fill[poporange!30] (-0.75,-1.7) .. controls (-0.75,-2.0) and (0.75,-2.0) .. (0.75,-1.7) .. controls (0.75,1.7) and (-0.75,1.7) .. (-0.75,-1.7);
        \node[poporange, align=center] at (0,-0.2){Solution\\concentrée + Eau};
        \draw[->, very thick, popblue] (0,0.5) -- (0,1.2) node[midway, right] {\textbf{Entrée d'eau}};
        \draw[->, thick, popgreen] (-1.2,1.8) -- (-1.2,1.2) node[midway, left, align=center]{Pression\\osmotique};
        \node at (0,-2.8) {\textbf{Après équilibre / Éclatement}};
    \end{scope}
\end{tikzpicture}
\legende{Expérience de Dutrochet : l'eau traverse la membrane semi-perméable (cæcum) vers le milieu hypertonique, gonflant la poche jusqu'à l'éclatement.}
\end{popfigure}
\end{experience}

\begin{attention}[Piège classique : « L'eau va vers l'eau »]
\emph{Erreur fréquente :} « L'eau passe du milieu riche en eau vers le milieu pauvre en eau. »
\emph{Correction rigoureuse :} En chimie physique, on raisonne en \textbf{concentration en solutés} (potentiel hydrique). L'eau diffuse du milieu \textbf{hypotonique} (faible concentration en solutés, potentiel hydrique élevé) vers le milieu \textbf{hypertonique} (forte concentration en solutés, potentiel hydrique bas). Dire « l'eau va vers le milieu concentré \emph{en solutés} » est correct ; dire « l'eau va vers le milieu concentré \emph{en eau} » est une contradiction.
\end{attention}

\courssub{L'expérience de Pfeffer : quantification de la pression osmotique}

Si Dutrochet a \emph{vu} le phénomène, Wilhelm Pfeffer (1877) l'a \emph{mesuré} en inventant l'osmomètre à membrane artificielle, permettant de quantifier la pression osmotique.

\begin{experience}[Osomètre de Pfeffer — Mesure de la pression osmotique]
\textbf{Matériel :} Un récipient en porcelaine poreuse enduit de \textbf{ferrocyanure de cuivre} ($\ce{Cu2[Fe(CN)6]}$), membrane artificielle \emph{perm sélective} : perméable à l'eau, imperméable au saccharose. Un tube capillaire gradué vertical. Une solution de saccharose de concentration connue ($C$). De l'eau pure.
\textbf{Protocole :}
\begin{enumerate}
    \item On remplit l'osmomètre (porcelaine + tube) avec la solution de saccharose.
    \item On plonge l'appareil dans un bain d'eau pure.
    \item On observe la montée du liquide dans le tube capillaire.
    \item La montée s'arrête lorsque la pression hydrostatique de la colonne de liquide équilibre l'entrée d'eau par osmose.
\end{enumerate}
\textbf{Observations et résultats :}
\begin{itemize}
    \item Le niveau dans le tube monte : l'eau entre par osmose.
    \item L'équilibre est atteint pour une hauteur $h$. La pression osmotique $\Pi$ est égale à la pression hydrostatique : $\Pi = \rho g h$ (avec $\rho$ masse volumique de la solution, $g \approx \SI{9,81}{N/kg}$).
    \item Pfeffer a montré que $\Pi$ est proportionnelle à la concentration $C$ (loi de van't Hoff : $\Pi = i C R T$, analogue à la loi des gaz parfaits).
\end{itemize}
\tcblower
\textbf{Schéma de l'osmomètre à l'équilibre :}
\begin{popfigure}
\begin{tikzpicture}[scale=0.85, every node/.style={font=\small}]
    % Bain eau pure
    \draw[popdark, thick] (-5,-3) rectangle (5,0);
    \fill[popblue!10] (-5,-3) rectangle (5,0);
    \node[popblue] at (0,-1.5) {Eau pure (Milieu hypotonique)};
    
    % Osmomètre (porcelaine poreuse + membrane Cu2[Fe(CN)6])
    \draw[poporange, very thick] (-1.5,0) -- (-1.5,4.5); % Paroi gauche
    \draw[poporange, very thick] (1.5,0) -- (1.5,4.5);   % Paroi droite
    \draw[poporange, very thick] (-1.5,4.5) -- (1.5,4.5); % Fond haut (bouchon)
    % Membrane au fond (z=0) : trait double pour montrer semi-perméabilité
    \draw[poppurple, ultra thick] (-1.5,0) -- (1.5,0);
    \draw[poppurple, thick, densely dashed] (-1.5,-0.1) -- (1.5,-0.1);
    \node[poppurple, rotate=90] at (-1.8,0) {Membrane artificielle (Ferrocyanure de Cu)};
    
    % Solution saccharose dans l'osmomètre
    \fill[poporange!20] (-1.4,0.1) rectangle (1.4,3.2);
    \node[poporange] at (0,1.6) {Solution de saccharose $C$ (Milieu hypertonique)};
    
    % Tube capillaire
    \draw[popdark, thick] (-0.2,4.5) -- (-0.2,7);
    \draw[popdark, thick] (0.2,4.5) -- (0.2,7);
    \fill[poporange!20] (-0.2,4.5) rectangle (0.2,6.2); % Montée liquide
    \node[poporange] at (0,5.3) {Colonne $h$};
    
    % Flèches mouvement eau
    \draw[->, very thick, popblue] (0,-0.5) -- (0,0.5) node[midway, right] {\textbf{Flux d'eau (Osmose)}};
    \draw[->, thick, popgreen] (2.5,3) -- (2.5,1) node[midway, right, align=center]{Pression\\hydrostatique $\rho g h$};
    \draw[<->, popdark] (1.8,4.5) -- (1.8,6.2) node[midway, right] {$h$};
    
    % Légende équilibre
    \node[align=center, popdark] at (4,3) {Équilibre :\\Flux entrant (osmose)\\= Flux sortant (pression)};
\end{tikzpicture}
\legende{Osomètre de Pfeffer : la montée du liquide dans le tube capillaire traduit la pression osmotique $\Pi = \rho g h$. La membrane au ferrocyanure de cuivre est sélective : elle laisse passer $\ce{H2O}$ mais retient le saccharose.}
\end{popfigure}
\end{experience}

\begin{definition}[Osmose et Pression Osmotique]
\emph{\textbf{Osmose :}} Diffusion passive du \textbf{solvant} (généralement l'eau) à travers une \textbf{membrane semi-perméable} (hémiperméable), du milieu \textbf{hypotonique} (moins concentré en solutés, potentiel hydrique $\Psi_w$ plus élevé) vers le milieu \textbf{hypertonique} (plus concentré en solutés, $\Psi_w$ plus bas).

\emph{\textbf{Membrane semi-perméable :}} Membrane dont la porosité ou la structure chimique permet le passage des molécules d'eau (petites, polaires) mais bloque le passage des molécules de solutés (ions, sucres, protéines).

\emph{\textbf{Pression osmotique ($\Pi$) :}} Pression qu'il faut appliquer sur la solution hypertonique pour arrêter le flux net d'eau entrant par osmose. C'est une pression de \textbf{colloïde} (dépend du nombre de particules solutées par unité de volume, pas de leur nature). Unité : \textbf{Pascal (Pa)}. $\Pi = i \cdot C \cdot R \cdot T$ (Loi de van't Hoff), avec $i$ coefficient de van't Hoff (nombre de particules par formule), $C$ concentration molaire ($\si{mol/L}$), $R = \SI{8,314}{J/(mol.K)}$, $T$ température (K).
\end{definition}

\courssub{La dialyse : diffusion des solutés à travers une membrane perméable}

L'osmose concerne le \textbf{solvant} (eau). La dialyse concerne les \textbf{solutés dissous}.

\begin{definition}[Dialyse — Cristalloïdes vs Colloïdes]
\emph{\textbf{Dialyse :}} Diffusion passive des \textbf{substances dissoutes} (solutés) à travers une \textbf{membrane perméable} (poreuse), du milieu où elles sont \textbf{plus concentrées} vers le milieu où elles sont \textbf{moins concentrées}, jusqu'à équilibre des concentrations.

\emph{\textbf{Classification des solutés selon leur taille (Thomas Graham, 1861) :}}
\begin{itemize}
    \item \textbf{Cristalloïdes :} Petites molécules ou ions (\ce{Na+}, \ce{Cl-}, \ce{Glucose}, \ce{Uree}, \ce{H2O}). Traversent \textbf{facilement} les membranes poreuses (pores $\approx \SI{1}{\nano\meter}$ à $\SI{10}{\nano\meter}$). Diffusent vite.
    \item \textbf{Colloïdes :} Grosses molécules (protéines : albumine, hémoglobine ; polysaccharides : amidon). Ne traversent \textbf{PAS} les membranes poreuses standards (rétention colloïdale). Restent du côté où elles ont été placées.
\end{itemize}
\textbf{Conséquence physiologique majeure :} Les protéines plasmatiques (albumine surtout), restées dans le sang par dialyse (elles ne traversent pas la paroi capillaire), créent une \textbf{pression oncotique} (pression osmotique colloïdale) qui retient l'eau dans les vaisseaux sanguins (rôle anti-œdème).
\end{definition}

\begin{popfigure}
\begin{tikzpicture}[scale=0.9, every node/.style={font=\small}]
    % Compartiment Gauche (Sang/Intérieur)
    \draw[popdark, thick] (-4,-3) rectangle (0,3);
    \fill[poppink!10] (-4,-3) rectangle (0,3);
    \node[poppink] at (-2,2.5) {Compartiment 1 : Sang / Intérieur cellule};
    
    % Membrane au centre
    \draw[poppurple, ultra thick] (0,-3) -- (0,3);
    \foreach \y in {-2.5,-1.5,...,2.5} {
        \draw[poppurple, thick] (-0.1,\y) -- (0.1,\y); % Pores symboliques
    }
    \node[poppurple, rotate=90] at (0.4,0) {Membrane poreuse (Dialyseur)};
    
    % Compartiment Droit (Dialysat/Extérieur)
    \draw[popdark, thick] (0,-3) rectangle (4,3);
    \fill[popblue!10] (0,-3) rectangle (4,3);
    \node[popblue] at (2,2.5) {Compartiment 2 : Dialysat / Milieu extérieur};
    
    % Cristalloïdes (petits cercles) - traversent
    \foreach \x/\y in {-3/2, -2/1, -1/0, -2.5/-1, -1.5/-2, 1/2, 2/1, 3/0, 1.5/-1, 2.5/-2} {
        \fill[popgreen] (\x,\y) circle (0.08);
    }
    \node[popgreen] at (-2,-2.7) {Cristalloïdes (\ce{Na+}, \ce{K+}, \ce{Uree}, \ce{Glucose})};
    \draw[->, thick, popgreen] (-0.5,0.5) -- (0.5,0.5);
    \draw[->, thick, popgreen] (0.5,-1) -- (-0.5,-1);
    \node[popgreen, font=\tiny] at (0,0.8) {Diffusion bidirectionnelle nette : $[C]_1 \to [C]_2$};
    
    % Colloïdes (grosses formes) - retenus à gauche
    \foreach \x/\y in {-3.5/1.5, -2.5/0.5, -1.5/-0.5, -3/-1.5} {
        \draw[poporange, thick] (\x,\y) ellipse (0.15 and 0.08);
    }
    \node[poporange] at (-2, 2.0) {Colloïdes (Protéines, Amidon)};
    \draw[poporange, decorate, decoration={snake, amplitude=1pt, segment length=4pt}] (-0.3,1.5) -- (0.3,1.5) node[midway, above, font=\tiny] {BLOQUÉS};
    
    % Flèches nettes
    \draw[->, very thick, popdark] (-3.5,-3.5) -- (-0.5,-3.5) node[midway, below] {Épuration (Élimination déchets)};
    \draw[->, very thick, popdark] (0.5,-3.5) -- (3.5,-3.5) node[midway, below] {Perte substances utiles (à compenser)};
\end{tikzpicture}
\legende{Principes de la dialyse : les cristalloïdes (verts) traversent la membrane poreuse selon leur gradient de concentration ; les colloïdes (orange) sont retenus. C'est le principe du rein artificiel.}
\end{popfigure}

\courssub{Interprétation unifiée des observations de la Section 1}

Grâce à l'osmose et à la dialyse, nous expliquons maintenant \emph{quantitativement} les faits observés :

\begin{propriete}[Bilan interprétatif des états cellulaires]
\begin{center}
\begin{tabularx}{\linewidth}{|l|X|X|X|}\hline
\rowcolor{popblueL}
\textbf{Type cellulaire} & \textbf{Milieu extérieur} & \textbf{Sens net de l'eau (Osmose)} & \textbf{État cellulaire observé} \\ \hline
\cellcolor{popgreenL}\textbf{Végétale} & \textbf{Hypotonique} (eau pure, sol dilué) & Entrée massive ($\Psi_{ext} > \Psi_{int}$) & \textbf{Turgescence} : membrane plaquée contre paroi, pression de turgesse $P_t > 0$. Cellule rigide. \\ \hline
\cellcolor{popgreenL}\textbf{Végétale} & \textbf{Isotonique} ($C_{ext} \approx C_{int}$) & Flux égaux (équilibre dynamique) & \textbf{État normal} : membrane légèrement décollée par endroits. \\ \hline
\cellcolor{popgreenL}\textbf{Végétale} & \textbf{Hypertonique} (sirop, $\ce{NaCl}$ concentré) & Sortie nette ($\Psi_{ext} < \Psi_{int}$) & \textbf{Plasmolyse} : membrane se décolle de la paroi, cytoplasme se rétracte. Réversible si retour isotonie. \\ \hline
\cellcolor{poppinkL}\textbf{Animale (Globule rouge)} & \textbf{Hypotonique} (ex. $\ce{NaCl}$ 5 g/L) & Entrée massive (pas de paroi) & \textbf{Hémolyse} : gonflement $\to$ éclatement (lyse). Disparition du globule, libération \ce{Hb}. \\ \hline
\cellcolor{poppinkL}\textbf{Animale} & \textbf{Isotonique} (ex. $\ce{NaCl}$ 9 g/L = sérum physiologique) & Équilibre & \textbf{État normal} : disque biconcave, diamètre \SI{7,5}{\micro\meter}. \\ \hline
\cellcolor{poppinkL}\textbf{Animale} & \textbf{Hypertonique} (ex. $\ce{NaCl}$ 15 g/L) & Sortie nette & \textbf{Crénelure} : rétrécissement, bords dentelés (échinocytes). Réversible. \\ \hline
\end{tabularx}
\end{center}
\textbf{Note :} La dialyse intervient pour l'élimination des déchets métaboliques (urée, $\ce{CO2}$) et l'apport en nutriments ($\ce{O2}$, glucose) qui suivent leurs gradients de concentration à travers la membrane plasmique (perm sélective + protéines de transport).
\end{propriete}

\begin{exemplebox}[Exemple chiffré : Le globule rouge face à la tonicité du milieu]
On place trois gouttes de sang frais dans trois tubes contenant des solutions de $\ce{NaCl}$ différentes. On observe au microscope optique ($\times 400$) après \SI{2}{\minute} :
\begin{itemize}
    \item \textbf{Tube A : $\ce{NaCl}$ 5 g/L (hypotonique).} Les globules gonflent, deviennent sphériques, puis \textbf{éclatent} (hémolyse). Le surnageant devient rouge transparent (hémoglobine libre). \emph{Interprétation :} $\Pi_{int} > \Pi_{ext}$, entrée d'eau $\Delta V > 0$, rupture membrane ($P_{\text{critique}} \approx \SI{10}{atm}$).
    \item \textbf{Tube B : $\ce{NaCl}$ 9 g/L (isotonique = sérum physiologique).} Les globules conservent leur forme \textbf{disque biconcave} caractéristique. Pas de variation de volume. \emph{Interprétation :} $\Pi_{int} = \Pi_{ext}$, flux d'eau nuls nets.
    \item \textbf{Tube C : $\ce{NaCl}$ 15 g/L (hypertonique).} Les globules \textbf{se crénellent} (forme d'étoile, échinocytes), volume diminué. \emph{Interprétation :} $\Pi_{int} < \Pi_{ext}$, sortie d'eau $\Delta V < 0$, concentration cytoplasmique augmente.
\end{itemize}
\textbf{Application médicale :} Toute perfusion intraveineuse \textbf{doit être isotonique} (sérum physiologique \SI{0,9}{\%} $\ce{NaCl}$, ou glucose \SI{5}{\%}) pour éviter l'hémolyse (si hypotonique) ou la crénelure/déshydratation cellulaire (si hypertonique).
\end{exemplebox}

\courssub{Concevoir un protocole pour mettre en évidence l'osmose (Savoir-faire Bac)}

\begin{methode}[Protocole type : Mise en évidence de l'osmose par variation de masse (Méthode des tubercules de pomme de terre / carotte)]
\textbf{Question scientifique :} Comment la concentration du milieu extérieur influence-t-elle le sens et l'amplitude des échanges d'eau (osmose) dans une cellule végétale ?
\textbf{Variables :}
\begin{itemize}
    \item Variable indépendante (cause) : Concentration molaire de la solution de saccharose ($C$ en $\si{mol/L}$) : \num{0,0} (eau), \num{0,1}, \num{0,2}, \num{0,3}, \num{0,4}, \num{0,5}.
    \item Variable dépendante (effet) : Variation relative de masse $\frac{\Delta m}{m_i} = \frac{m_f - m_i}{m_i}$ (en \%) ou variation de longueur/turgescence.
    \item Variables contrôlées : Température (\SI{25}{\celsius}), durée (\SI{30}{\minute} à \SI{1}{\hour}), volume solution (\SI{20}{mL}), géométrie des cylindres (tronçonneuse à légumes : $\varnothing \SI{1}{cm}$, hauteur \SI{3}{cm}), même pomme de terre (même potentiel hydrique initial), essuyage doux papier buvard avant pesée finale.
\end{itemize}
\textbf{Matériel :} Pomme de terre (ou carotte, courgette), tronçonneuse à légumes / emporte-pièce, balance de précision (\SI{0,01}{g}), béchers \SI{50}{mL}, solutions de saccharose étalonnées, papier buvard, étiquettes.
\textbf{Protocole :}
\begin{enumerate}
    \item Préparer \num{6} cylindres identiques de pomme de terre. Les essuyer, les peser ($m_i$), noter la masse.
    \item Placer chaque cylindre dans un bécher distinct contenant \SI{20}{mL} d'une des \num{6} solutions. Démarrer le chronomètre.
    \item Après \SI{45}{\minute} (temps suffisant pour approche équilibre), retirer les cylindres.
    \item Les essuyer \textbf{très délicatement} et \textbf{complètement} (erreur majeure si eau de surface non enlevée) avec papier buvard.
    \item Peser immédiatement chaque cylindre ($m_f$).
    \item Calculer $\frac{\Delta m}{m_i} \times 100$. Tracer le graphique $\frac{\Delta m}{m_i} = f(C)$.
\end{enumerate}
\textbf{Résultats attendus et Interprétation :}
\begin{itemize}
    \item Courbe décroissante. Point d'intersection avec l'axe des abscisses ($\Delta m = 0$) $\to$ concentration isotonique $C_{iso}$ (équivalent à la concentration cellulaire interne).
    \item $C < C_{iso}$ (hypotonique) : $\Delta m > 0$ (gain masse = entrée eau = turgescence).
    \item $C > C_{iso}$ (hypertonique) : $\Delta m < 0$ (perte masse = sortie eau = plasmolyse).
\end{itemize}
\textbf{Témoin :} Le tube à \num{0,0} $\si{mol/L}$ (eau pure) montre le gain de masse maximal (turgescence maximale). L'absence de témoin « temps zéro » (pesée initiale) invaliderait la quantification.
\end{methode}

\begin{exoresolu}[Exercice type Bac : Détermination de la concentration isotonique]
\textbf{Énoncé :} On réalise le protocole ci-dessus sur des cylindres de carotte. Résultats :
\begin{center}
\begin{tabular}{|c|c|c|c|c|c|c|}\hline
$C$ (\si{mol/L}) & 0,00 & 0,10 & 0,20 & 0,30 & 0,40 & 0,50 \\ \hline
$\Delta m / m_i$ (\% ) & +18,2 & +10,5 & +2,1 & -5,8 & -13,4 & -20,1 \\ \hline
\end{tabular}
\end{center}
\begin{enumerate}
    \item Tracer le graphique $\frac{\Delta m}{m_i} = f(C)$.
    \item Déterminer graphiquement la concentration isotonique $C_{iso}$ de la cellule de carotte.
    \item Interpréter le signe de $\Delta m$ pour $C = \SI{0,10}{mol/L}$ et $C = \SI{0,40}{mol/L}$.
    \item Calculer la pression osmotique $\Pi$ du cytoplasme de la carotte à \SI{25}{\celsius} (on suppose $i=1$ pour le saccharose non dissocié).
\end{enumerate}
\tcblower
\textbf{Solution détaillée :}
\begin{enumerate}
    \item \textbf{Graphique :} Axe des abscisses $C$ (\si{mol/L}) de 0 à 0,5. Axe des ordonnées $\Delta m/m_i$ (\%) de -25 à +20. Points alignés grossièrement. Droite de régression tracée.
    \item \textbf{Concentration isotonique :} Lecture de l'abscisse pour $\Delta m = 0$. Interpolation entre 0,20 (+2,1) et 0,30 (-5,8).
    \[
    C_{iso} \approx \num{0,20} + \frac{\num{2,1}}{\num{2,1} + \num{5,8}} \times (\num{0,30} - \num{0,20}) = \num{0,20} + \frac{\num{2,1}}{\num{7,9}} \times \num{0,10} \approx \mathbf{\num{0,227}\ mol/L}
    \]
    (On retient \textbf{\SI{0,23}{mol/L}}).
    \item \textbf{Interprétation :}
    \begin{itemize}
        \item $C = \SI{0,10}{mol/L} < C_{iso}$ : Milieu \textbf{hypotonique}. $\Psi_{ext} > \Psi_{int}$. L'eau entre par osmose $\to$ \textbf{gain de masse} ($+10,5\%$). La cellule devient turgescente.
        \item $C = \SI{0,40}{mol/L} > C_{iso}$ : Milieu \textbf{hypertonique}. $\Psi_{ext} < \Psi_{int}$. L'eau sort par osmose $\to$ \textbf{perte de masse} ($-13,4\%$). La cellule se plasmolyse.
    \end{itemize}
    \item \textbf{Pression osmotique :} $\Pi = i C R T$.
    $i=1$, $C = \SI{0,227}{mol/L} = \SI{227}{mol/m^3}$, $R = \SI{8,314}{J/(mol.K)}$, $T = \SI{298,15}{K}$.
    \[
    \Pi = 1 \times 227 \times \num{8,314} \times \num{298,15} \approx \SI{562000}{Pa} = \mathbf{\num{5,62} \times 10^5\ Pa} \ (\approx \mathbf{5,6\ bar} \text{ ou } \mathbf{5,5\ atm}).
    \]
    \emph{Remarque :} Cette pression est considérable ! C'est elle qui assure la rigidité (turgescence) des tissus végétaux non lignifiés (jeunes pousses, feuilles, pétioles).
\end{enumerate}
\end{exoresolu}

\begin{savaistu}[La dialyse rénale artificielle : un rein de poche]
L'insuffisance rénale terminale (stade 5) empêche l'élimination de l'urée, de la créatinine, du potassium en excès et de l'eau. La \textbf{hémodialyse} (rein artificiel) utilise exactement le principe de la dialyse découvert par Graham et quantifié par Pfeffer :
\begin{itemize}
    \item Le sang du patient circule d'un côté d'une \textbf{membrane semi-perméable synthétique} (polysulfone, pores $\sim \SI{10}{\nano\meter}$).
    \item De l'autre côté circule le \textbf{dialysat} (solution stérile, composition ionique contrôlée : $\ce{Na+}$ \SI{140}{mmol/L}, $\ce{K+}$ \SI{2}{mmol/L}, $\ce{Ca^{2+}}$ \SI{1,5}{mmol/L}, $\ce{HCO3-}$ \SI{32}{mmol/L}, glucose \SI{5,5}{mmol/L}, \textbf{urée = 0}).
    \item \textbf{Gradient de concentration} : Urée, créatinine, $\ce{K+}$ excédentaire diffusent du sang (concentré) vers le dialysat (dilué) $\to$ \textbf{épuration}.
    \item \textbf{Ultrafiltration} : En appliquant une pression hydrostatique négative côté dialysat (ou positive côté sang), on force le passage d'eau (osmose forcée) $\to$ \textbf{déshydratation} contrôlée du patient.
    \item Les \textbf{protéines (albumine)} et \textbf{cellules sanguines} (colloïdes) sont retenues par la membrane.
\end{itemize}
Au Cameroun, les centres d'hémodialyse (Hôpital Général de Yaoundé, Laquintinie Douala, centres privés) traitent des centaines de patients, 3 séances de \SI{4}{\hour} par semaine. Le coût reste un défi majeur de santé publique. La \textbf{dialyse péritonéale} (utilisant le péritoine comme membrane naturelle) est une alternative domiciliaire en développement.
\end{savaistu}

\begin{aretenir}[Synthèse : Osmose vs Dialyse]
\begin{center}
\begin{tabularx}{\linewidth}{|l|X|X|}\hline
\rowcolor{popblueL}
\textbf{Critère} & \textbf{Osmose} & \textbf{Dialyse} \\ \hline
\textbf{Substance qui bouge} & Solvant (\textbf{Eau}) & Solutés dissous (\textbf{Cristalloïdes}) \\ \hline
\textbf{Membrane} & Semi-perméable (hémiperméable) & Poreuse (perm sélective par la taille) \\ \hline
\textbf{Moteur} & Gradient de potentiel hydrique ($\Delta \Psi_w$) / Concentration en solutés & Gradient de concentration du soluté ($\Delta C$) \\ \hline
\textbf{Sens} & Hypotonique $\to$ Hypertonique & Concentré $\to$ Dilué \\ \hline
\textbf{Pression d'équilibre} & Pression osmotique $\Pi$ (loi van't Hoff) & Égalité des concentrations ($C_1 = C_2$) \\ \hline
\textbf{Exemple cellulaire} & Turgescence / Plasmolyse / Hémolyse / Crénelure & Nutrition cellulaire (glucose, AA), Éxcrétion (urée, $\ce{CO2}$) \\ \hline
\end{tabularx}
\end{center}
\textbf{Distinction cruciale :} L'osmose modifie le \textbf{volume} cellulaire (entrée/sortie d'eau). La dialyse modifie la \textbf{composition chimique} du cytoplasme (entrée/sortie de solutés). Les deux sont des processus \textbf{passifs} (ne coûtent pas d'ATP directement), obéissant à la seconde loi de la thermodynamique (augmentation de l'entropie).
\end{aretenir}

\coursec{Échanges de substances dissoutes : transport passif}

Après avoir analysé les mouvements d'eau à travers la membrane plasmique (osmose) et séparé les solutés selon leur taille (dialyse), nous nous posons maintenant une question fondamentale : \textbf{comment les substances dissoutes (ions, glucose, acides aminés, gaz) franchissent-elles la barrière lipidique de la membrane pour entrer ou sortir de la cellule ?} Contrairement à l'eau qui traverse la bicouche par simple diffusion (et via les aquaporines), la plupart des solutés polaires ou chargés ne peuvent pas diffuser librement à travers le cœur hydrophobe de la membrane. Leur passage nécessite des voies protéiques spécifiques. Étudions les deux grands modes de \textbf{transport passif} : la diffusion simple et la diffusion facilitée.

\courssub{La diffusion simple : passage direct à travers la bicouche lipidique}

\textbf{Mécanisme et molécules concernées}\par

Seules les molécules de \textbf{petite taille} et \textbf{non polaires (liposolubles)}, ainsi que quelques très petites molécules polaires non chargées, peuvent traverser la bicouche phospholipidique sans aide protéique. Elles se dissolvent dans le cœur hydrophobe de la membrane et diffusent selon leur gradient de concentration.

\begin{definition}[Diffusion simple]
Mouvement net de molécules ou d'ions à travers la membrane plasmique, \textbf{directement à travers la bicouche lipidique}, du compartiment où leur concentration (ou potentiel électrochimique) est la plus élevée vers celui où elle est la plus faible, \textbf{sans intervention de protéine transporteuse} et \textbf{sans consommation d'énergie métabolique (ATP)}.
\end{definition}

\noindent
\textbf{Principaux substrats :} dioxygène ($\ce{O2}$), dioxyde de carbone ($\ce{CO2}$), azote ($\ce{N2}$), éthanol, stéroïdes, petites molécules liposolubles (vitamines D, A, E, K).

\begin{popfigure}
\begin{tikzpicture}[scale=1.1, every node/.style={font=\small}]
% Bicouche lipidique : deux feuillets
\draw[fill=poporangeL!50, draw=poporange!80!black, thick] (-3.5, 0.5) rectangle (3.5, 1.2) node[midway, text=popdark, align=center] {Têtes polaires\\ (face extracellulaire)};
\draw[fill=poporangeL!20, draw=poporange!80!black, thick] (-3.5, -0.5) rectangle (3.5, 0.5) node[midway, text=popdark, align=center] {Queues hydrophobes\\ (acides gras)};
\draw[fill=poporangeL!50, draw=poporange!80!black, thick] (-3.5, -1.2) rectangle (3.5, -0.5) node[midway, text=popdark, align=center] {Têtes polaires\\ (face cytoplasmique)};

% Molécules O2 et CO2 traversant
\foreach \y in {-0.8, 0, 0.8} {
    \draw[popblue, thick, ->] (-4.5, \y) -- (-3.5, \y);
    \node[popblue, left] at (-4.5, \y) {$\ce{O2}$ / $\ce{CO2}$};
    \draw[popblue, thick, ->] (3.5, \y) -- (4.5, \y);
}
% Légendes milieux
\node[align=center, text width=3.5cm] at (-5.5, -1.5) {Milieu extracellulaire\\ (forte [$\ce{O2}$])};
\node[align=center, text width=3.5cm] at (5.5, -1.5) {Cytoplasme\\ (faible [$\ce{O2}$])};
\draw[decorate, decoration={brace, amplitude=5pt}, popdark] (-4, 1.2) -- (4, 1.2) node[midway, above=5pt] {Membrane plasmique (bicouche lipidique)};
\end{tikzpicture}
\legende{Schéma de la \cle{diffusion simple} : les gaz respiratoires ($\ce{O2}$, $\ce{CO2}$) traversent librement la bicouche lipidique selon leur gradient de concentration. Aucune protéine n'est requise.}
\end{popfigure}

\textbf{Cinétique de la diffusion simple : proportionnalité à la concentration}\par

Pour caractériser ce transport, on mesure la \textbf{vitesse d'entrée ($v$)} d'une substance dans la cellule (ex. hématies) en fonction de sa \textbf{concentration extracellulaire ($C_e$)}.

\begin{experience}[Cinétique de la diffusion simple]
\textbf{Matériel :} Hématies humaines, solutions de $\ce{O2}$ ou d'éthanol à concentrations croissantes ($C_e$), spectrophotomètre ou électrode à $\ce{O2}$.
\textbf{Protocole :} On incube des hématies dans des solutions de concentrations connues. On mesure la vitesse initiale d'entrée de la substance (variation de concentration intracellulaire par unité de temps).
\textbf{Résultats attendus :} La vitesse $v$ est proportionnelle à $C_e$. Le graphe $v = f(C_e)$ est une \textbf{droite passant par l'origine}.
\textbf{Interprétation :} Plus le gradient de concentration est fort, plus les chocs efficaces entre molécules et membrane sont fréquents. Il n'y a pas de limite de nombre de « portes » d'entrée.
\end{experience}

\begin{popfigure}
\begin{tikzpicture}
\begin{axis}[
    width=\linewidth, height=8cm,
    axis lines=middle,
    xlabel={Concentration extracellulaire $C_e$ (mmol/L)},
    ylabel={Vitesse d'entrée $v$ (mmol$\cdot$L$^{-1}\cdot$s$^{-1}$)},
    xmin=0, xmax=50,
    ymin=0, ymax=12,
    xtick={0,10,20,30,40,50},
    ytick={0,2,4,6,8,10,12},
    grid=major, grid style={popline},
    legend pos=north west,
    legend style={font=\small, fill=popcream, draw=popdark},
    domain=0:50, samples=100,
]
% Diffusion simple : droite
\addplot[popblue, very thick, samples=2] {0.2*x};
\addlegendentry{\cle{Diffusion simple} ($v = k \cdot C_e$)}

% Diffusion facilitée : saturation (Michaelis-Menten)
\addplot[poporange, very thick] {10*x/(15+x)};
\addlegendentry{\cle{Diffusion facilitée} ($v = \frac{V_{max} \cdot C_e}{K_m + C_e}$)}

% Annotations
\node[popblue, anchor=south west] at (axis cs:45,9.5) {Pente = $k$ (perméabilité)};
\node[poporange, anchor=south east] at (axis cs:45,9.5) {Plateau = $V_{max}$};
\draw[poporange, dashed] (axis cs:0,10) -- (axis cs:50,10) node[right, pos=0.5, poporange] {$V_{max}$};
\draw[poporange, dashed] (axis cs:15,0) -- (axis cs:15,10) node[above, pos=0.5, poporange] {$K_m$};
\end{axis}
\end{tikzpicture}
\legende{Cinétique du transport passif. \textbf{Courbe bleue (\cle{diffusion simple})} : vitesse proportionnelle à la concentration (droite passant par l'origine). \textbf{Courbe orange (\cle{diffusion facilitée})} : vitesse croissante puis plateau de saturation ($V_{max}$) traduisant le nombre fini de transporteurs. $K_m$ = concentration pour $v = V_{max}/2$ (affinité du transporteur).}
\end{popfigure}

\begin{propriete}[Cinétique de la diffusion simple]
La vitesse de transport $v$ est proportionnelle à la concentration extracellulaire $C_e$ (tant que la concentration intracellulaire reste négligeable) :
\[ v = P \cdot C_e \]
où $P$ est le coefficient de perméabilité (dépend de la taille, de la lipophilicité, de l'épaisseur de la membrane). Le graphe $v = f(C_e)$ est une \textbf{droite passant par l'origine}. Il n'y a \textbf{pas de saturation}.
\end{propriete}

\courssub{La diffusion facilitée : transport par protéines spécifiques, sans énergie}

\textbf{Preuve expérimentale : la saturation du transport}\par

Si l'on réalise la même expérience cinétique avec du \textbf{glucose} sur des hématies (globules rouges), on obtient une courbe radicalement différente : la vitesse augmente avec la concentration, mais \textbf{plafonne} à une valeur maximale $V_{max}$.

\begin{experience}[Mise en évidence de la diffusion facilitée (Glucose / Hématies)]
\textbf{Matériel :} Hématies fraîches, solutions de glucose marquées ($^{14}\text{C}$-glucose ou $^{3}\text{H}$-glucose) à concentrations variées ($0$ à $50$ mmol/L), filtration rapide sur filtre à nitrocellulose, scintillation liquide.
\textbf{Protocole :} Incubation courte (secondes) à $37\,^\circ\text{C}$ pour mesurer la vitesse \emph{initiale} d'influx ($v_0$). Arrêt brutal par filtration et lavage à froid. Mesure de la radioactivité intracellulaire.
\textbf{Résultats (données classiques) :}
\begin{center}
\begin{tabularx}{\linewidth}{|c|X|X|}\hline
$C_e$ (mmol/L) & $v_0$ (mmol$\cdot$L$^{-1}\cdot$s$^{-1}$) & Observation \\ \hline
1 & 0,6 & Croissance rapide \\
5 & 2,5 & \\
10 & 4,2 & \\
15 & 5,3 & \\
20 & 6,0 & Ralentissement \\
25 & 6,4 & \textbf{Plateau} \\
30 & 6,5 & \textbf{Plateau} \\
50 & 6,6 & \textbf{Plateau ($V_{max} \approx 6,6$)} \\ \hline
\end{tabularx}
\end{center}
\textbf{Interprétation :} Le plateau ($V_{max}$) prouve que le nombre de sites de passage est \textbf{fini}. Le glucose ne traverse pas la bicouche lipidique (imperméable au glucose), mais emprunte des \textbf{protéines transporteuses (GLUT1 sur hématies)} en nombre limité. Lorsque tous les transporteurs sont occupés (saturés), la vitesse ne peut plus augmenter.
\end{experience}

\begin{definition}[Diffusion facilitée]
Transport passif d'une substance à travers la membrane plasmique \textbf{via une protéine transporteuse spécifique (canal ou transporteur)}, \textbf{dans le sens du gradient électrochimique} (de la concentration forte vers la faible), \textbf{sans consommation d'ATP}. Elle présente une \textbf{spécificité} (un transporteur pour un ou quelques substrats) et une \textbf{saturation} (vitesse maximale $V_{max}$).
\end{definition}

\textbf{Deux types de protéines pour la diffusion facilitée}\par

\begin{itemize}
    \item \textbf{Les canaux ioniques (\cle{Channel proteins}) :} Pores aqueux traversant la membrane. Ouverture/fermeture commandée (voltage, ligand, mécanique). Passage très rapide ($\sim 10^7$ ions/s). Ex. : canal $\ce{K+}$ voltage-dépendant, canal $\ce{Na+}$ voltage-dépendant, \cle{aquaporines} (eau).
    \item \textbf{Les transporteurs / perméases (\cle{Carrier proteins}) :} Protéines à conformation variable. Fixation du substrat $\rightarrow$ changement de conformation $\rightarrow$ libération de l'autre côté. Plus lents ($\sim 10^2$--$10^4$ molécules/s). Ex. : \textbf{GLUT1} (glucose érythrocytaire), GLUT4 (insulino-dépendant, muscle/adipose), transporteurs d'acides aminés.
\end{itemize}

\begin{popfigure}
\begin{tikzpicture}[scale=1.0, every node/.style={font=\small}]
% Membrane
\draw[fill=poporangeL!30, draw=poporange!80!black, thick] (-4.5, -1.2) rectangle (4.5, 1.2);
\draw[dashed, popdark] (-4.5, 0) -- (4.5, 0) node[midway, above=2pt] {Milieu extracellulaire};
\node[align=center] at (0, -1.6) {Cytoplasme};

% Canal ionique (gauche)
\draw[fill=popblue!30, draw=popblue, thick, rounded corners=2pt] (-3, -1.2) rectangle (-2, 1.2);
\draw[fill=popblue!60] (-2.5, 0) circle (0.35);
\draw[popblue, thick, ->] (-2.5, 1.5) -- (-2.5, 0.6);
\draw[popblue, thick, ->] (-2.5, -0.6) -- (-2.5, -1.5);
\node[popblue, align=center] at (-2.5, 1.8) {Canal $\ce{K+}$};
\node[popblue, align=center] at (-2.5, -1.9) {$\ce{K+}$};

% Transporteur GLUT (droite) - conformation ouverte vers l'extérieur
\draw[fill=popgreen!30, draw=popgreen, thick, rounded corners=3pt] (2, 0.2) -- (3, 0.6) -- (3.5, 0.2) -- (3, -0.2) -- (2, 0.2);
\draw[fill=popgreen!30, draw=popgreen, thick, rounded corners=3pt] (2, -0.2) -- (3, -0.6) -- (3.5, -0.2) -- (3, 0.2) -- (2, -0.2);
\draw[fill=poppink, draw=poppink!80!black] (2.7, 0) circle (0.15) node[white, scale=0.7] {Glc};
\draw[popgreen, thick, ->] (2.7, 1.0) -- (2.7, 0.5);
\draw[popgreen, thick, ->] (2.7, -0.5) -- (2.7, -1.0);
\node[popgreen, align=center] at (2.7, 1.3) {GLUT1};
\node[popgreen, align=center] at (2.7, -1.3) {Glucose};

% Légendes
\node[align=left, text width=4.5cm] at (-5.5, 0.5) {\textbf{Canal} : pore fixe, passage rapide, sélectivité par taille/charge.};
\node[align=left, text width=4.5cm] at (-5.5, -0.8) {\textbf{Transporteur} : site de fixation, changement conformationnel, spécificité forte, saturation.};
\end{tikzpicture}
\legende{Deux modes de \cle{diffusion facilitée}. À gauche : un \textbf{canal ionique} (ex. $\ce{K+}$) forme un pore aqueux permanent ou commandé. À droite : un \textbf{transporteur (GLUT1)} lie le glucose, change de conformation pour le libérer de l'autre côté. Les deux sont passifs (gradient électrochimique).}
\end{popfigure}

\textbf{Critères distinctifs du transport passif (simple ou facilité)}\par

\begin{aretenir}[Critères du transport passif]
Pour affirmer qu'un transport est \textbf{passif}, il faut vérifier \textbf{trois critères} simultanés :
\begin{enumerate}
    \item \textbf{Sens du gradient électrochimique :} du potentiel électrochimique élevé vers le faible (concentration forte $\rightarrow$ faible pour un neutre ; selon gradient concentration + charge pour un ion).
    \item \textbf{Absence de consommation d'ATP :} le transport se poursuit en présence d'inhibiteurs de la respiration mitochondriale (cyanure, DNP, azide) ou en l'absence de glucose (pas de glycolyse). L'énergie provient uniquement de l'entropie (désordre) du système.
    \item \textbf{Cinétique :} proportionnelle à $C_e$ (simple) ou hyperbolique avec saturation (facilité).
\end{enumerate}
\end{aretenir}

\begin{methode}[Exploiter un graphe $v = f(C_e)$ pour identifier le mode de transport]
\begin{enumerate}
    \item \textbf{Observer la forme de la courbe} : droite passant par l'origine $\rightarrow$ \cle{diffusion simple}. Courbe hyperbolique avec plateau $\rightarrow$ \cle{diffusion facilitée}.
    \item \textbf{Si plateau (saturation) :} lire la vitesse maximale $V_{max}$ (ordonnée du plateau) et la concentration $K_m$ (abscisse pour $v = V_{max}/2$). $K_m$ bas = forte affinité du transporteur.
    \item \textbf{Conclure :} La saturation prouve l'intervention d'un nombre fini de protéines transporteuses $\rightarrow$ diffusion facilitée. L'absence de saturation $\rightarrow$ diffusion simple (ou nombre de transporteurs non limitant sur l'intervalle testé).
\end{enumerate}
\end{methode}

\begin{attention}[Piège classique : « Protéine = Actif »]
\textbf{FAUX.} La \cle{diffusion facilitée} fait intervenir des protéines (canaux, transporteurs) mais \textbf{ne consomme pas d'ATP}. L'énergie est fournie par le gradient de concentration (énergie libre de Gibbs $\Delta G < 0$). Seuls les \textbf{transports actifs} (primaire : ATPase ; secondaire : symport/antiport couplé à un gradient ionique) consomment de l'énergie métabolique directement ou indirectement.
\begin{itemize}
    \item \textbf{Indice au Bac :} Si le sujet précise « en présence de cyanure (inhibiteur cytochrome oxydase), le transport persiste » $\rightarrow$ c'est passif. S'il s'arrête $\rightarrow$ c'est actif (ou secondaire couplé à un gradient entretenu par une ATPase).
\end{itemize}
\end{attention}

\courssub{Exemple concret : l'entrée du glucose dans l'hématie (globule rouge)}

L'hématie (érythrocyte) est le modèle historique d'étude du transport du glucose. Elle exprime massivement le transporteur \textbf{GLUT1} (environ $2 \times 10^5$ copies/cellule).

\begin{exemplebox}[Paramètres cinétiques du GLUT1 chez l'homme]
\begin{itemize}
    \item $V_{max} \approx \SI{100}{nmol \cdot mg^{-1} \cdot min^{-1}}$ (vitesse maximale d'influx par mg de protéine membranaire).
    \item $K_m \approx \SI{5}{mmol/L}$ (concentration extracellulaire pour demi-vitesse maximale).
    \item Glycémie normale à jeun : $\approx \SI{5}{mmol/L}$ ($\sim \SI{0,90}{g/L}$).
    \item Glycémie post-prandiale : $\approx \SI{7}{--}\SI{10}{mmol/L}$.
\end{itemize}
\textbf{Interprétation physiologique :} À la glycémie normale ($\SI{5}{mmol/L}$), $v \approx \frac{100 \times 5}{5 + 5} = \SI{50}{nmol \cdot mg^{-1} \cdot min^{-1}}$ (seulement $50\%$ de $V_{max}$). Le transporteur est \textbf{sous-saturé} : l'entrée du glucose est proportionnelle à la glycémie (régulation fine). En cas d'hyperglycémie sévère ($> \SI{25}{mmol/L}$), le transporteur sature : l'entrée ne suit plus la glycémie, ce qui limite l'hyperglycémie intracellulaire (protection relative contre la glycation des protéines).
\end{exemplebox}

\courssub{Exercice résolu : Analyse d'une courbe de saturation}

\begin{exoresolu}[Identification du mode de transport d'un acide aminé]
\textbf{Énoncé :} On étudie l'influx d'un acide aminé (AA) dans des vésicules membranaires purifiées. La vitesse initiale $v_0$ (en nmol$\cdot$mg prot$^{-1}\cdot$min$^{-1}$) est mesurée en fonction de la concentration extracellulaire $[AA]_e$ (en mmol/L). Résultats :

\begin{center}
\begin{tabular}{|c|c|c|c|c|c|c|c|}\hline
$[AA]_e$ & 0,5 & 1 & 2 & 5 & 10 & 20 & 50 \\ \hline
$v_0$ & 12 & 22 & 38 & 68 & 88 & 96 & 98 \\ \hline
\end{tabular}
\end{center}

\textbf{1. Représentez graphiquement $v_0 = f([AA]_e)$.}
\textbf{2. Déterminez le mode de transport (simple ou facilité). Justifiez.}
\textbf{3. Estimez $V_{max}$ et $K_m$.}
\textbf{4. Précisez si ce transport consomme de l'ATP.}

\tcblower

\textbf{Solution détaillée :}

\textbf{1. Graphe :} (Voir figure ci-dessous). Les points s'alignent sur une courbe hyperbolique croissante tendant vers un plateau.

\begin{popfigure}
\begin{tikzpicture}
\begin{axis}[
    width=\linewidth, height=7cm,
    axis lines=middle,
    xlabel={$[AA]_e$ (mmol/L)},
    ylabel={$v_0$ (nmol$\cdot$mg$^{-1}\cdot$min$^{-1}$)},
    xmin=0, xmax=55,
    ymin=0, ymax=110,
    xtick={0,10,20,30,40,50},
    ytick={0,20,40,60,80,100},
    grid=major, grid style={popline},
    legend pos=south east,
]
\addplot[only marks, mark=*, popblue, thick] coordinates {
    (0.5,12) (1,22) (2,38) (5,68) (10,88) (20,96) (50,98)
};
\addlegendentry{Données exp.}
% Fit Michaelis-Menten approximatif pour le tracé guide
\addplot[poporange, thick, dashed, domain=0:55, samples=100] {100*x/(2.5+x)};
\addlegendentry{Modèle $V_{max}=100$, $K_m=2,5$}
\end{axis}
\end{tikzpicture}
\legende{Graphe $v_0 = f([AA]_e)$ pour l'exercice résolu. La courbe orange en pointillés est un ajustement de type Michaelis-Menten ($V_{max} \approx 100$, $K_m \approx 2,5$).}
\end{popfigure}

\textbf{2. Mode de transport :} \textbf{Diffusion facilitée.}
\textbf{Justification :} La courbe n'est pas une droite. Elle présente un \textbf{plateau de saturation} vers $98$--$100$ pour les fortes concentrations. Cela traduit l'existence d'un nombre fini de sites de transport (protéines transporteuses), caractéristique de la diffusion facilitée.

\textbf{3. Estimation des paramètres :}
\begin{itemize}
    \item $V_{max} \approx \SI{100}{nmol \cdot mg^{-1} \cdot min^{-1}}$ (valeur du plateau).
    \item $K_m$ : concentration pour $v_0 = V_{max}/2 = 50$. Par lecture graphique ou interpolation entre $[AA]_e=1$ ($v=22$) et $[AA]_e=2$ ($v=38$), $K_m \approx \SI{2,5}{mmol/L}$.
\end{itemize}

\textbf{4. Consommation d'ATP :} \textbf{Non.} C'est un transport passif (diffusion facilitée). Il se fait dans le sens du gradient de concentration (milieu extracellulaire $\rightarrow$ vésicule interne) sans apport d'énergie métabolique. La présence d'inhibiteurs respiratoires (cyanure, oligomycine) n'altérerait pas cette cinétique.
\end{exoresolu}

\courssub{Synthèse comparative : Diffusion simple vs Diffusion facilitée}

\begin{center}
\begin{tabularx}{\linewidth}{|l|X|X|}\hline
\rowcolor{popblueL}
\textbf{Critère} & \textbf{Diffusion simple} & \textbf{Diffusion facilitée} \\ \hline
\textbf{Voie de passage} & Bicouche lipidique (interstices) & Protéines transmembranaires (canaux ou transporteurs) \\ \hline
\textbf{Substrats typiques} & $\ce{O2}$, $\ce{CO2}$, $\ce{N2}$, éthanol, stéroïdes, molécules liposolubles & Ions ($\ce{Na+}$, $\ce{K+}$, $\ce{Cl-}$, $\ce{Ca^{2+}}$), glucose, acides aminés, eau (aquaporines), nucléosides \\ \hline
\textbf{Spécificité} & Faible (taille, lipophilicité) & \textbf{Élevée} (site de fixation stéréospécifique) \\ \hline
\textbf{Saturation} & \textbf{Non} (pas de limite protéique) & \textbf{Oui} ($V_{max}$ atteint quand tous transporteurs occupés) \\ \hline
\textbf{Cinétique $v = f(C_e)$} & Droite passant par l'origine ($v = P \cdot C_e$) & Hyperbole de Michaelis-Menten ($v = \frac{V_{max} \cdot C_e}{K_m + C_e}$) \\ \hline
\textbf{Vitesse maximale} & Illimitée (théoriquement) & Limitée par $V_{max}$ (nb transporteurs $\times$ taux de turnover) \\ \hline
\textbf{Régulation} & Physico-chimique (température, fluidité membrane) & \textbf{Biologique} : nombre de protéines (expression génique, trafic vésiculaire ex. GLUT4/insuline), modulation allostérique, phosphorylation \\ \hline
\textbf{Énergie (ATP)} & Aucune (passif) & Aucune (passif) \\ \hline
\textbf{Sens} & Gradient de concentration (fort $\rightarrow$ faible) & Gradient électrochimique (fort $\rightarrow$ faible) \\ \hline
\textbf{Inhibition} & Quasi impossible (sauf changement physique membrane) & \textbf{Oui} : inhibiteurs compétitifs (ressemblent substrat) ou non compétitifs (autre site) \\ \hline
\end{tabularx}
\end{center}

\begin{savaistu}[Le glucose et le paludisme au Cameroun]
Le parasite du paludisme (\emph{Plasmodium falciparum}) se développe à l'intérieur de l'hématie. Il a \textbf{besoin de glucose} pour sa glycolyse (seule source d'ATP). Il induit l'apparition de \textbf{nouveaux canaux de perméabilité (NPP - New Permeability Pathways)} dans la membrane de l'hématie infectée, augmentant massivement l'entrée de glucose (et d'autres solutés) par diffusion facilitée. Comprendre la cinétique de transport (saturation, inhibiteurs) a permis de tester des molécules bloquant ces canaux comme stratégie antipaludique. Au Cameroun, où le paludisme reste la première cause de consultation, la recherche sur les transporteurs de l'hématie infectée est un axe majeur de l'Institut de Recherche Médicale et d'Études des Plantes Médicinales (IMPM) et de l'Université de Yaoundé I.
\end{savaistu}

\begin{exercice}[Entraînement : Interprétation de courbes cinétiques]
On compare l'entrée de trois substances A, B, C dans un même type cellulaire. Les courbes $v = f(C_e)$ sont tracées ci-dessous.
\begin{center}
\begin{tikzpicture}
\begin{axis}[
    width=0.8\linewidth, height=6cm,
    axis lines=middle,
    xlabel={$C_e$}, ylabel={$v$},
    xmin=0, xmax=10, ymin=0, ymax=10,
    xtick=\empty, ytick=\empty,
    domain=0:10, samples=100,
]
\addplot[popblue, very thick] {x}; \node[popblue, right] at (axis cs:9,9) {A};
\addplot[poporange, very thick] {8*x/(2+x)}; \node[poporange, right] at (axis cs:9,7.5) {B};
\addplot[popgreen, very thick] {5}; \node[popgreen, right] at (axis cs:9,5) {C};
\end{axis}
\end{tikzpicture}
\end{center}
\textbf{1.} Identifiez le mode de transport pour chaque substance (simple, facilité, ou autre ?).
\textbf{2.} Pour la substance B, estimez $V_{max}$ et $K_m$.
\textbf{3.} La substance C pourrait-elle correspondre à un transport actif ? Justifiez.
\textbf{4.} Proposez une identité possible pour A, B, C (ex. $\ce{O2}$, glucose, $\ce{Na+}$ entrant par symport $\ce{Na+/glucose}$).
\end{exercice}

\begin{corrige}[Exercice : Interprétation de courbes cinétiques]
\textbf{1.}
\begin{itemize}
    \item \textbf{A : Droite passant par l'origine} $\rightarrow$ \textbf{Diffusion simple}. Vitesse proportionnelle à $C_e$, pas de saturation.
    \item \textbf{B : Courbe hyperbolique avec plateau} $\rightarrow$ \textbf{Diffusion facilitée}. Saturation = nombre fini de transporteurs.
    \item \textbf{C : Droite horizontale ($v$ constant, indépendant de $C_e$)} $\rightarrow$ \textbf{Transport actif (ou facilité saturé dès le début)}. Vitesse constante = $V_{max}$ atteinte même à basse concentration. Cela suggère un transporteur à \textbf{très forte affinité} ($K_m$ très bas) saturé dès les plus faibles $C_e$ testées, \textbf{OU} un transport actif (vitesse imposée par le turnover de la pompe, indépendante du gradient tant que substrat présent).
\end{itemize}
\textbf{2.} Pour B : $V_{max} \approx 8$ (unité de $v$). $K_m \approx 2$ (abscisse où $v = 4 = V_{max}/2$).
\textbf{3.} Oui, C est compatible avec un \textbf{transport actif} (vitesse constante, indépendante de $C_e$, énergie fournie par ATP). Mais une diffusion facilitée avec $K_m \ll C_{e,min}$ donnerait le même aspect. Il faut tester la sensibilité aux inhibiteurs métaboliques (cyanure, ouate) pour trancher : si $v$ chute $\rightarrow$ actif ; si $v$ inchangé $\rightarrow$ facilité à forte affinité.
\textbf{4.} Propositions : \textbf{A} = $\ce{O2}$ ou $\ce{CO2}$ (gaz, liposolubles). \textbf{B} = Glucose (via GLUT) ou acide aminé. \textbf{C} = $\ce{Na+}$ entrant par symport $\ce{Na+/glucose}$ (SGLT1, actif secondaire) ou $\ce{K+}$ entrant par pompe $\ce{Na+/K+}$-ATPase (actif primaire) ; ou encore un transporteur à très forte affinité (ex. transporteur sérotonine SERT, $K_m \sim \SI{0,1}{\mu M}$).
\end{corrige}

\begin{aretenir}[Bilan : Le transport passif des solutés]
\begin{itemize}
    \item La membrane est \textbf{sélectivement perméable} : lipides $\rightarrow$ diffusion simple ; ions/polaires $\rightarrow$ protéines (diffusion facilitée ou transport actif).
    \item \textbf{Diffusion simple} : traversant la bicouche, non spécifique, non saturable, cinétique linéaire. Ex. : gaz respiratoires.
    \item \textbf{Diffusion facilitée} : via canaux ou transporteurs (GLUT, canaux $\ce{K+}$), spécifique, saturable ($V_{max}$, $K_m$), cinétique hyperbolique. \textbf{Pas d'ATP}.
    \item L'analyse du graphe $v = f(C_e)$ est l'outil diagnostique principal (savoir-faire Bac).
    \item Ces échanges assurent l'approvisionnement en nutriments (glucose), l'évacuation des déchets ($\ce{CO2}$, urée), la signalisation (neurotransmetteurs) et l'excitabilité (canaux ioniques).
\end{itemize}
\end{aretenir}

\coursec{Échanges de substances dissoutes : transport actif}

Dans la section précédente, nous avons vu que la cellule peut laisser passer certaines substances \emph{dans le sens du gradient}, sans dépenser d'énergie : c'est le transport passif (diffusion simple ou facilitée). Mais une difficulté subsiste, et elle est considérable : comment une racine de plante peut-elle absorber les ions minéraux d'un sol dont la solution est \emph{beaucoup plus diluée} que le suc cellulaire ? Comment l'intestin peut-il capter les derniers ions d'un bol alimentaire presque épuisé ? Dans ces situations, la substance doit voyager « à contre-courant », du milieu pauvre vers le milieu riche. C'est exactement le rôle du \cle{transport actif}, que nous allons maintenant étudier.

\courssub{Situation-problème : l'absorption des ions minéraux par les racines}

\textbf{Observation.} Un plant de maïs cultivé sur un sol du haut plateau de l'Ouest camerounais prélève dans la solution du sol des ions comme \ce{K+}, \ce{NO3-} ou \ce{PO4^3-}. Or les mesures montrent que la concentration en ions \ce{K+} dans les cellules des poils absorbants peut atteindre environ \num{100} mmol/L alors que la solution du sol n'en contient parfois que \num{0,1} mmol/L : la cellule racinaire \emph{accumule} le potassium contre un gradient de concentration de 1 à \num{1000}.

\textbf{Question.} La diffusion, même facilitée, ne peut faire passer une substance que du milieu le plus concentré vers le milieu le moins concentré. Par quel mécanisme la cellule racinaire réussit-elle donc à faire entrer des ions \emph{contre} leur gradient ?

\textbf{Hypothèse.} Ce transport « à contre-courant » nécessite un travail, donc une dépense d'énergie fournie par la cellule. Vérifions cette hypothèse par l'expérience.

\courssub{Mise en évidence expérimentale : le transport contre le gradient exige de l'énergie}

\begin{experience}[Absorption d'ions par des racines en présence ou non d'inhibiteurs du métabolisme]
\textbf{Matériel.} Jeunes plants de la même espèce (par exemple des plants de maïs), solution nutritive diluée contenant un ion marqué et dosable (par exemple \ce{K+}), enceintes thermostatées, glace, gaz azote (\ce{N2}) pour chasser le dioxygène, cyanure de potassium (KCN, inhibiteur de la respiration cellulaire — manipulé uniquement en milieu encadré).

\textbf{Protocole.} On place des lots identiques de racines dans la même solution diluée de \ce{K+} et on mesure la quantité d'ions absorbés par unité de temps dans quatre conditions :
\begin{enumerate}
  \item lot témoin : \SI{25}{\celsius}, milieu aéré ;
  \item lot A : \SI{4}{\celsius} (froid), milieu aéré ;
  \item lot B : \SI{25}{\celsius}, milieu maintenu sans \ce{O2} (barbotage d'azote) ;
  \item lot C : \SI{25}{\celsius}, milieu aéré additionné de cyanure.
\end{enumerate}

\textbf{Résultats.} La vitesse d'absorption de \ce{K+} est forte dans le lot témoin, mais elle s'effondre (proche de zéro) dans les trois autres lots.

\begin{center}
\begin{tabularx}{\linewidth}{|l|X|c|}
\hline
\rowcolor{popblueL} \textbf{Lot} & \textbf{Condition expérimentale} & \textbf{Absorption de \ce{K+}} \\
\hline
Témoin & \SI{25}{\celsius}, présence de \ce{O2} & forte (100 \%) \\
\hline
A & froid (\SI{4}{\celsius}) & très faible ($\approx$ 10 \%) \\
\hline
B & absence de \ce{O2} & quasi nulle ($\approx$ 5 \%) \\
\hline
C & cyanure (inhibiteur respiratoire) & quasi nulle ($\approx$ 3 \%) \\
\hline
\end{tabularx}
\end{center}

\textbf{Interprétation.} Le froid ralentit les réactions enzymatiques ; l'absence de \ce{O2} et le cyanure bloquent la respiration cellulaire, c'est-à-dire la production d'\cle{ATP}. Dans les trois cas, ce qui est supprimé, c'est la \emph{production d'énergie} par la cellule. Or l'absorption des ions s'arrête alors, bien que le gradient de concentration soit inchangé.

\textbf{Conclusion.} L'absorption des ions minéraux contre leur gradient de concentration nécessite de l'énergie fournie par la respiration cellulaire sous forme d'ATP. Il s'agit donc d'un mécanisme \emph{actif}.
\end{experience}

\begin{attention}[Le froid ne « gèle » pas la membrane]
Une erreur fréquente consiste à dire que le froid « bloque la membrane ». Non : le froid ralentit l'activité enzymatique, donc la respiration cellulaire, donc la production d'ATP. C'est le manque d'énergie, et non une modification mécanique de la membrane, qui arrête le transport actif. Rédige toujours la chaîne causale complète : froid $\rightarrow$ ralentissement de la respiration $\rightarrow$ baisse de la production d'ATP $\rightarrow$ arrêt du transport actif.
\end{attention}

\courssub{Définition du transport actif}

\begin{definition}[Transport actif]
Le \cle{transport actif} est le déplacement d'une substance dissoute (ion ou molécule) à travers la membrane plasmique, \textbf{contre son gradient électrochimique} (c'est-à-dire du milieu où elle est la moins concentrée vers le milieu où elle est la plus concentrée), grâce à des \textbf{protéines membranaires spécifiques appelées pompes}, avec \textbf{consommation d'énergie} fournie par l'hydrolyse de l'ATP :
\[ \ce{ATP + H2O -> ADP + Pi} + \text{énergie} \]
\end{definition}

Trois conditions cumulatives définissent donc le transport actif : un déplacement \emph{contre le gradient}, l'intervention d'une \emph{protéine transporteuse} (la pompe), et une \emph{dépense d'ATP}. Si l'une de ces trois conditions manque, ce n'est pas du transport actif.

\courssub{Un exemple fondamental : la pompe \ce{Na+}/\ce{K+}}

La pompe \ce{Na+}/\ce{K+} (ou \ce{Na+}/\ce{K+}-ATPase) est une protéine transmembranaire présente dans la membrane de toutes les cellules animales. Elle fonctionne en continu et constitue l'exemple canonique de transport actif au programme.

\begin{propriete}[Fonctionnement de la pompe \ce{Na+}/\ce{K+} (admis)]
Pour \textbf{chaque molécule d'ATP hydrolysée}, la pompe \ce{Na+}/\ce{K+} :
\begin{itemize}
  \item \textbf{expulse 3 ions \ce{Na+}} de la cellule vers le milieu extracellulaire ;
  \item \textbf{fait entrer 2 ions \ce{K+}} du milieu extracellulaire vers le cytosol.
\end{itemize}
Ce bilan $3\ \ce{Na+}$ sortants contre $2\ \ce{K+}$ entrants est \emph{admis} : aucune démonstration n'est exigible au Baccalauréat, mais tu dois savoir l'utiliser.
\end{propriete}

\begin{popfigure}
\begin{center}
\begin{tikzpicture}[scale=1.0]
% Milieu extracellulaire / cytosol
\fill[popblueL] (-5,1.0) rectangle (5,3.2);
\fill[popcream] (-5,-3.2) rectangle (5,-1.0);
% Bicouche membranaire
\fill[popgold!45] (-5,-1.0) rectangle (5,-0.55);
\fill[popgold!45] (-5,0.55) rectangle (5,1.0);
\foreach \x in {-4.8,-4.4,...,4.8}{
  \fill[popgold] (\x,0.775) circle (0.09);
  \fill[popgold] (\x,-0.775) circle (0.09);
}
\node at (-4.6,2.6) {\small\textbf{Milieu extracellulaire}};
\node at (-4.6,2.1) {\small $[\ce{Na+}]$ forte, $[\ce{K+}]$ faible};
\node at (-4.6,-2.6) {\small\textbf{Cytosol}};
\node at (-4.6,-2.1) {\small $[\ce{Na+}]$ faible, $[\ce{K+}]$ forte};
% Pompe
\draw[fill=popgreenL,draw=popgreen!70!black,thick] (-0.9,-1.1) .. controls (-1.3,-0.3) and (-1.3,0.3) .. (-0.9,1.1) -- (0.9,1.1) .. controls (1.3,0.3) and (1.3,-0.3) .. (0.9,-1.1) -- cycle;
\node at (0,0) {\small\textbf{Pompe}};
\node at (0,-0.35) {\small\textbf{\ce{Na+}/\ce{K+}}};
% Flèches Na+ sortantes (3)
\foreach \x in {1.9,2.7,3.5}{
  \node[circle,fill=poporange,draw=poporange!60!black,inner sep=1.5pt] at (\x,-2.0) {\tiny\ce{Na+}};
  \draw[-{Stealth[length=2.5mm]},very thick,poporange!70!black] (\x,-1.7) -- (\x,2.0);
  \node[circle,fill=poporange,draw=poporange!60!black,inner sep=1.5pt] at (\x,2.3) {\tiny\ce{Na+}};
}
\node[poporange!70!black] at (2.7,2.85) {\small\textbf{3 \ce{Na+} expulsés}};
% Flèches K+ entrantes (2)
\foreach \x in {-2.2,-3.1}{
  \node[circle,fill=poppurple,draw=poppurple!60!black,inner sep=1.5pt,text=white] at (\x,2.0) {\tiny\ce{K+}};
  \draw[-{Stealth[length=2.5mm]},very thick,poppurple!70!black] (\x,1.7) -- (\x,-2.0);
  \node[circle,fill=poppurple,draw=poppurple!60!black,inner sep=1.5pt,text=white] at (\x,-2.3) {\tiny\ce{K+}};
}
\node[poppurple!70!black] at (-2.65,-2.85) {\small\textbf{2 \ce{K+} captés}};
% Hydrolyse ATP
\node[draw=popink,fill=poppinkL,rounded corners,align=center] at (4.0,-0.9) {\small \ce{ATP -> ADP + Pi}\\ \small + énergie};
\draw[-{Stealth[length=2mm]},thick,popink] (3.2,-0.9) -- (1.35,-0.6);
\end{tikzpicture}
\end{center}
\legende{Schéma fonctionnel de la pompe \ce{Na+}/\ce{K+}. La protéine transmembranaire (vert) hydrolyse une molécule d'ATP pour expulser 3 ions \ce{Na+} (orange) et faire entrer 2 ions \ce{K+} (violet), tous deux contre leur gradient de concentration.}
\end{popfigure}

\textbf{Rôle physiologique.} En expulsant 3 charges positives pour seulement 2 charges positives entrantes, la pompe rend l'intérieur de la cellule \emph{légèrement négatif} par rapport à l'extérieur : elle contribue ainsi au \cle{potentiel de membrane} (environ \SI{-70}{mV} dans un neurone au repos). Elle maintient aussi une forte concentration interne en \ce{K+} et une forte concentration externe en \ce{Na+}. Ces deux conséquences sont indispensables à la \emph{naissance et à la propagation du message nerveux}, que tu étudieras plus tard dans le programme : retiens dès maintenant que la pompe \ce{Na+}/\ce{K+} prépare les conditions électriques du fonctionnement du neurone.

\begin{savaistu}[Une pompe très gourmande en énergie]
Chez un animal au repos, environ \textbf{un tiers de l'ATP} produit par les cellules est consommé par la seule pompe \ce{Na+}/\ce{K+} ! Chez les neurones, cette proportion peut dépasser les deux tiers. C'est l'une des raisons pour lesquelles le cerveau, qui ne représente qu'environ 2 \% de la masse corporelle, consomme près de 20 \% du dioxygène de l'organisme. Une intoxication au cyanure, qui bloque la production d'ATP, est donc si dangereuse précisément parce qu'elle paralyse ces pompes et effondre les gradients ioniques vitaux.
\end{savaistu}

\courssub{Caractéristiques du transport actif}

Le transport actif possède quatre caractéristiques qu'il faut savoir énoncer et justifier :

\begin{enumerate}
  \item \textbf{Spécificité.} Chaque pompe ne transporte qu'une substance (ou un couple de substances) déterminée : la pompe \ce{Na+}/\ce{K+} ne transporte ni glucose ni \ce{Ca^2+}. Cette spécificité tient à la structure tridimensionnelle du site de fixation de la protéine, comme une serrure et sa clé.
  \item \textbf{Saturation.} Le nombre de pompes dans une membrane est limité. Quand toutes les pompes fonctionnent à leur vitesse maximale, augmenter encore la concentration de la substance n'augmente plus la vitesse de transport : la courbe vitesse $= f(\text{concentration})$ présente un \emph{plateau}, comme en diffusion facilitée.
  \item \textbf{Sens inverse du gradient.} La substance va du milieu le moins concentré vers le milieu le plus concentré : c'est la signature du transport actif.
  \item \textbf{Dépendance énergétique.} Le transport s'arrête en l'absence d'ATP (froid intense, anoxie, cyanure), ce que ne fait jamais un transport passif.
\end{enumerate}

\begin{attention}[Ne confonds pas « contre le gradient » et « vers le milieu concentré »]
L'erreur classique du candidat est d'écrire : « le transport actif fait passer la substance du milieu concentré vers le milieu dilué ». C'est l'inverse ! \textbf{Contre le gradient} signifie : du milieu \emph{moins} concentré vers le milieu \emph{plus} concentré — comme remonter une rivière à la nage, ce qui exige un effort (l'ATP). Astuce mnémotechnique : la diffusion « descend » la pente (gratuit), le transport actif « monte » la pente (payant).
\end{attention}

\courssub{Comment distinguer expérimentalement transport passif et transport actif ?}

\begin{methode}[Le test de l'inhibiteur énergétique]
Pour déterminer si le transport d'une substance à travers une membrane est passif ou actif :
\begin{enumerate}
  \item Mesurer la vitesse de transport de la substance dans des conditions normales (témoin : température optimale, présence de \ce{O2}).
  \item Refaire la mesure en supprimant la production d'ATP : absence de \ce{O2} (anoxie), ou ajout d'un inhibiteur respiratoire (cyanure), ou abaissement fort de la température.
  \item Comparer les vitesses :
  \begin{itemize}
    \item si le transport \textbf{s'arrête ou chute fortement} $\Rightarrow$ il dépend de l'ATP : c'est un \textbf{transport actif} ;
    \item si le transport \textbf{se poursuit normalement} $\Rightarrow$ il ne consomme pas d'énergie : c'est un \textbf{transport passif} (diffusion simple ou facilitée).
  \end{itemize}
  \item Pour trancher ensuite entre diffusion simple et diffusion facilitée, tracer la vitesse en fonction de la concentration : un plateau de saturation indique une diffusion facilitée (protéine en nombre limité).
\end{enumerate}
\end{methode}

\courssub{Bilan comparatif des trois modes de transport membranaire}

\begin{popfigure}
\begin{center}
\begin{tikzpicture}
\node[draw=popdark,thick,rounded corners,inner sep=6pt, align=center]{
\begin{tabular}{|p{2.6cm}|p{2.9cm}|p{2.9cm}|p{2.9cm}|}
\hline
\rowcolor{popblueL} \textbf{Critère} & \textbf{Diffusion simple} & \textbf{Diffusion facilitée} & \textbf{Transport actif} \\
\hline
Sens du transport & Dans le sens du gradient & Dans le sens du gradient & \textbf{Contre} le gradient \\
\hline
Protéine membranaire & Aucune & Oui (canal ou transporteur) & Oui (pompe ATPase) \\
\hline
Énergie (ATP) & Non & Non & \textbf{Oui} \\
\hline
Saturation & Non (vitesse proportionnelle) & Oui (plateau) & Oui (plateau) \\
\hline
Spécificité & Faible (taille, liposolubilité) & Forte & Forte \\
\hline
Exemples & \ce{O2}, \ce{CO2}, éthanol & Glucose dans les hématies, ions via canaux & Pompe \ce{Na+}/\ce{K+}, absorption des ions minéraux par les racines \\
\hline
\end{tabular}
};
\end{tikzpicture}
\end{center}
\legende{Tableau comparatif des trois modes de passage des substances dissoutes à travers la membrane plasmique. Seul le transport actif s'oppose au gradient et consomme de l'ATP.}
\end{popfigure}

\begin{aretenir}[L'essentiel]
\begin{itemize}
  \item Le \cle{transport actif} déplace une substance \textbf{contre son gradient électrochimique}, via une \textbf{pompe} membranaire, au prix d'une hydrolyse d'\textbf{ATP}.
  \item La pompe \ce{Na+}/\ce{K+} expulse \textbf{3 \ce{Na+}} et fait entrer \textbf{2 \ce{K+}} par ATP hydrolysé ; elle contribue au potentiel de membrane.
  \item Le test décisif : un transport stoppé par l'anoxie, le froid ou le cyanure est un transport actif.
  \item Grâce au transport actif, la cellule maintient des concentrations internes différentes du milieu extérieur : c'est une \textbf{condition de la vie cellulaire}.
\end{itemize}
\end{aretenir}

\courssub{Application et exercice résolu}

\begin{exoresolu}[Absorption du potassium par des racines de tomate]
Des racines de tomate sont placées dans une solution contenant \num{0,5} mmol/L de \ce{K+}. On mesure la concentration interne en \ce{K+} des cellules racinaires : elle atteint \num{80} mmol/L après quelques heures, et l'absorption se poursuit. En ajoutant du cyanure à la solution, l'absorption cesse presque totalement, alors que la concentration interne reste supérieure à celle du milieu.

\textbf{Questions.} 1) Montrer que l'absorption de \ce{K+} ne peut pas être une diffusion. 2) Quel est le mode de transport mis en jeu ? Justifier par deux arguments. 3) Expliquer le résultat obtenu avec le cyanure.
\tcblower
\textbf{Solution.}

1) La concentration interne (\num{80} mmol/L) devient très supérieure à celle du milieu externe (\num{0,5} mmol/L), et pourtant \ce{K+} continue d'entrer dans la racine. Or une diffusion, simple ou facilitée, ne transporte une substance que du milieu le plus concentré vers le milieu le moins concentré. L'absorption se faisant ici \emph{contre le gradient}, elle ne peut pas être une diffusion.

2) Il s'agit d'un \textbf{transport actif}. Argument 1 : le transport s'effectue du milieu le moins concentré (\num{0,5} mmol/L) vers le milieu le plus concentré (\num{80} mmol/L), donc contre le gradient. Argument 2 : le cyanure, inhibiteur de la respiration cellulaire, arrête l'absorption, ce qui prouve que le transport dépend de l'ATP produit par la respiration.

3) Le cyanure bloque la chaîne respiratoire mitochondriale : les cellules racinaires ne produisent presque plus d'ATP. Or les pompes à \ce{K+} ont besoin de l'énergie libérée par l'hydrolyse de l'ATP pour fonctionner. Privées d'ATP, elles s'arrêtent, et l'absorption de \ce{K+} cesse, même si le gradient de concentration demeure favorable à la sortie des ions.
\end{exoresolu}

\courssub{Conséquence fondamentale : la cellule construit et entretient sa différence}

Le transport actif explique à lui seul un fait capital : le milieu intracellulaire n'est \emph{jamais} en équilibre avec le milieu extérieur. Le cytosol est riche en \ce{K+} et pauvre en \ce{Na+} ; le liquide extracellulaire est riche en \ce{Na+} et pauvre en \ce{K+} ; les cellules racinaires concentrent les ions minéraux ; les cellules intestinales concentrent le glucose. Ces \emph{différences de composition}, créées et entretenues en permanence au prix d'une dépense continue d'ATP, sont la condition de nombreuses fonctions vitales : excitabilité des neurones et des muscles, absorption des nutriments, maintien du volume cellulaire. Une cellule qui cesse de produire de l'ATP voit ses gradients s'effondrer — et meurt. Le transport actif n'est donc pas un luxe : c'est l'un des piliers de la vie cellulaire.

\begin{exercice}[Pour t'entraîner]
On mesure la vitesse d'absorption d'un ion X par des cellules en fonction de la concentration externe en X, en présence de \ce{O2} (courbe 1) puis en absence de \ce{O2} (courbe 2). La courbe 1 montre une absorption qui se poursuit même lorsque la concentration interne dépasse la concentration externe ; la courbe 2 montre une absorption quasi nulle dans les mêmes conditions. 1) Quel est le mode de transport de X ? Justifier. 2) Quelle expérience complémentaire permettrait de vérifier que des protéines membranaires interviennent ?
\end{exercice}

\begin{corrige}[Exercice]
1) Il s'agit d'un transport actif, pour deux raisons : d'une part l'ion X est absorbé contre son gradient de concentration (la concentration interne devient supérieure à la concentration externe et l'absorption continue) ; d'autre part l'absence de \ce{O2}, qui supprime la respiration cellulaire donc la production d'ATP, abolit l'absorption : le transport dépend de l'énergie.

2) On peut tracer la vitesse d'absorption en fonction de la concentration externe en X : si la courbe présente un plateau de saturation (la vitesse plafonne malgré l'augmentation de la concentration), cela prouve l'intervention de protéines transporteuses en nombre limité. On peut aussi utiliser un inhibiteur spécifique des protéines (par exemple un agent bloquant les ATPases) et montrer que l'absorption cesse.
\end{corrige}

\coursec{Échanges de particules : endocytose et exocytose}

Nous venons de voir que les substances dissoutes — ions, glucose, acides aminés — traversent la membrane plasmique par diffusion simple, diffusion facilitée ou transport actif. Mais ces mécanismes ont une limite de taille : ils concernent des molécules ou des ions de petites dimensions. Or la cellule doit parfois absorber des particules volumineuses (une bactérie entière, des débris cellulaires, des gouttelettes de liquide) ou expulser des macromolécules (enzymes digestives, hormones, neuromédiateurs). Comment la cellule résout-elle ce problème ? C'est l'objet de cette section : les \cle{échanges de particules} par \cle{endocytose} et \cle{exocytose}.

\courssub{Situation-problème : le casse-tête des grosses particules}

\textbf{Observation.} Au CHU de Yaoundé, un technicien de laboratoire observe au microscope optique un frottis sanguin : on y voit des leucocytes (globules blancs) au contact de bactéries. Quelques minutes plus tard, les bactéries ont disparu du milieu extracellulaire et se retrouvent \emph{à l'intérieur} des leucocytes, où elles sont digérées.

\textbf{Problème.} Une bactérie mesure environ \SI{1}{\micro\meter} de diamètre, soit mille fois plus qu'une molécule de glucose. Aucun canal membranaire, aucune pompe ne peut la faire passer. Comment entre-t-elle dans le leucocyte ?

\textbf{Hypothèse.} La membrane plasmique, souple et fluide, peut se déformer pour \emph{englober} la particule dans une poche, sans que la particule ne « traverse » jamais la membrane elle-même.

C'est exactement ce que révèlent les électronographies : la membrane s'invagine, entoure la particule, puis se referme en une \cle{vésicule} intracellulaire. La particule reste donc toujours séparée du cytoplasme par une membrane. Ce mécanisme, qui consomme de l'énergie (ATP), est l'\cle{endocytose}.

\courssub{L'endocytose : entrée de particules dans la cellule}

\begin{definition}[Endocytose]
L'\cle{endocytose} est le mécanisme par lequel une cellule fait pénétrer dans son cytoplasme des particules volumineuses (solides ou liquides) : la membrane plasmique s'\cle{invagine} autour de la particule, puis se referme et se détache en formant une \cle{vésicule} intracellulaire. C'est un processus \cle{actif} : il consomme de l'ATP et repose sur la fluidité de la membrane.
\end{definition}

On distingue deux formes d'endocytose selon la nature de la particule absorbée.

\begin{definition}[Phagocytose]
La \cle{phagocytose} (du grec \emph{phagein}, manger : la « cellule mangeuse ») est l'absorption de \cle{particules solides} (bactéries, débris cellulaires, poussières) par la cellule. La membrane émet des expansions appelées \cle{pseudopodes} qui entourent la particule et fusionnent : la particule est alors enfermée dans une vésicule appelée \cle{phagosome}. Le phagosome fusionne ensuite avec un \cle{lysosome} (vésicule riche en enzymes digestives) pour former une vacuole digestive (ou phagolysosome) dans laquelle la particule est détruite.
\end{definition}

\begin{definition}[Pinocytose]
La \cle{pinocytose} (du grec \emph{pinein}, boire : la « cellule buveuse ») est l'absorption de \cle{gouttelettes de liquide extracellulaire} et des macromolécules qui y sont dissoutes (protéines, lipoprotéines). La membrane s'invagine en une petite dépression qui se referme en une \cle{vésicule de pinocytose}, de taille bien inférieure à celle du phagosome.
\end{definition}

\begin{experience}[Observation de la phagocytose par un leucocyte]
\textbf{Matériel :} lame portant un frottis de sang infecté, microscope optique à contraste de phase, électronographies de macrophages en action.

\textbf{Protocole et observations :}
\begin{enumerate}
\item Un leucocyte (macrophage) au contact d'une bactérie émet des pseudopodes qui s'allongent de part et d'autre de la bactérie.
\item Les pseudopodes fusionnent au-delà de la bactérie : celle-ci se retrouve enfermée dans un phagosome, à l'intérieur de la cellule.
\item Des lysosomes migrent vers le phagosome et fusionnent avec lui ; la bactérie est digérée en quelques dizaines de minutes.
\item Les résidus non digestibles sont évacués hors de la cellule.
\end{enumerate}

\textbf{Interprétation :} la bactérie n'a à aucun moment traversé la membrane : elle est entrée \emph{enveloppée} dans un repli de celle-ci. La déformation membranaire et les déplacements de vésicules nécessitent de l'énergie fournie par l'hydrolyse de l'ATP.

\textbf{Conclusion :} la phagocytose est une endocytose de particules solides, réalisée grâce à la fluidité membranaire et consommatrice d'ATP.
\end{experience}

\begin{popfigure}
\begin{center}
\begin{tikzpicture}[scale=0.92, every node/.style={font=\small}]
% Étape 1 : contact
\begin{scope}[shift={(0,0)}]
\draw[fill=popblueL, draw=popblue, thick] (0,0) circle (1.15);
\draw[fill=poppurpleL, draw=poppurple] (0.2,0.25) circle (0.35);
\node[font=\scriptsize] at (0.2,0.25) {noyau};
\draw[fill=poporange, draw=popdark] (1.55,0) circle (0.22);
\node[font=\scriptsize, below=2pt] at (1.55,-0.28) {bactérie};
\node[font=\small\bfseries, popdark] at (0,-1.6) {1. Contact};
\end{scope}
% Étape 2 : pseudopodes
\begin{scope}[shift={(4.4,0)}]
\draw[fill=popblueL, draw=popblue, thick] (0,0) circle (1.15);
\draw[fill=poppurpleL, draw=poppurple] (-0.25,0.1) circle (0.35);
\draw[fill=popblueL, draw=popblue, thick] (1.0,0.42) .. controls (1.5,0.55) and (1.85,0.45) .. (1.95,0.12);
\draw[fill=popblueL, draw=popblue, thick] (1.0,-0.42) .. controls (1.5,-0.55) and (1.85,-0.45) .. (1.95,-0.12);
\draw[fill=poporange, draw=popdark] (1.75,0) circle (0.22);
\node[font=\scriptsize, popblue, above=1pt] at (1.6,0.62) {pseudopodes};
\node[font=\small\bfseries, popdark] at (0.4,-1.6) {2. Émission de pseudopodes};
\end{scope}
% Étape 3 : phagosome
\begin{scope}[shift={(8.8,0)}]
\draw[fill=popblueL, draw=popblue, thick] (0,0) circle (1.15);
\draw[fill=poppurpleL, draw=poppurple] (-0.35,0.3) circle (0.35);
\draw[fill=popcream, draw=popgold, thick] (0.45,-0.25) circle (0.42);
\draw[fill=poporange, draw=popdark] (0.45,-0.25) circle (0.2);
\node[font=\scriptsize, popgold!60!black] at (0.45,-0.95) {phagosome};
\node[font=\small\bfseries, popdark] at (0,-1.6) {3. Formation du phagosome};
\end{scope}
% Étape 4 : digestion lysosomiale
\begin{scope}[shift={(13.2,0)}]
\draw[fill=popblueL, draw=popblue, thick] (0,0) circle (1.15);
\draw[fill=poppurpleL, draw=poppurple] (-0.4,0.35) circle (0.32);
\draw[fill=popgreenL, draw=popgreen!70!black, thick] (0.45,-0.2) circle (0.5);
\draw[fill=popgreen, draw=popgreen!70!black] (0.1,0.55) circle (0.14);
\draw[fill=popgreen, draw=popgreen!70!black] (0.85,0.5) circle (0.14);
\node[font=\scriptsize, popgreen!50!black] at (0.45,-1.0) {phagolysosome};
\node[font=\scriptsize, popgreen!50!black, above=1pt] at (0.85,0.66) {lysosomes};
\node[font=\small\bfseries, popdark] at (0,-1.6) {4. Digestion lysosomiale};
\end{scope}
% flèches entre étapes
\draw[-{Stealth[length=3mm]}, thick, popmuted] (2.6,0) -- (3.1,0);
\draw[-{Stealth[length=3mm]}, thick, popmuted] (7.0,0) -- (7.5,0);
\draw[-{Stealth[length=3mm]}, thick, popmuted] (11.4,0) -- (11.9,0);
\end{tikzpicture}
\end{center}
\legende{Schéma des quatre étapes de la phagocytose d'une bactérie par un leucocyte. 1 : contact de la bactérie avec la membrane ; 2 : émission de pseudopodes qui entourent la particule ; 3 : fermeture de la vésicule formant le phagosome ; 4 : fusion avec un lysosome formant un phagolysosome et digestion enzymatique de la bactérie. La particule reste en permanence entourée d'une membrane.}
\end{popfigure}

\begin{exemplebox}[La phagocytose, base de la défense immunitaire]
Un macrophage peut phagocyter plusieurs dizaines de bactéries au cours de sa vie. Lors d'une angine bactérienne ou d'une plaie infectée, des millions de leucocytes affluent vers le foyer infectieux et englobent les microbes : le \emph{pus} est en grande partie constitué de leucocytes morts après avoir phagocyté. La phagocytose est donc l'un des piliers de la défense de l'organisme contre les infections — un enjeu majeur de santé publique au Cameroun, où les infections bactériennes restent fréquentes.
\end{exemplebox}

\courssub{L'exocytose : sortie de particules hors de la cellule}

\begin{definition}[Exocytose]
L'\cle{exocytose} est le mécanisme inverse de l'endocytose : une \cle{vésicule intracellulaire} (issue du réticulum endoplasmique et de l'appareil de Golgi) migre vers la membrane plasmique, \cle{fusionne} avec elle et \cle{libère son contenu} dans le milieu extracellulaire. C'est également un processus actif, consommateur d'ATP.
\end{definition}

L'exocytose est le mode de \cle{sécrétion} de toutes les grosses molécules produites par la cellule :
\begin{itemize}
\item les \cle{enzymes digestives} sécrétées par le pancréas et les glandes salivaires dans le tube digestif ;
\item les \cle{hormones} protéiques (insuline sécrétée par le pancréas endocrine, par exemple) libérées dans le sang ;
\item les \cle{neuromédiateurs} (comme l'acétylcholine) libérés par les neurones au niveau des synapses ;
\item les \cle{anticorps} sécrétés par les plasmocytes.
\end{itemize}

\begin{popfigure}
\begin{center}
\begin{tikzpicture}[scale=0.95, every node/.style={font=\small}]
% ===== Pinocytose (gauche) =====
\begin{scope}[shift={(0,0)}]
\node[font=\small\bfseries, popblue] at (0,2.6) {Pinocytose (entrée)};
% milieu extracellulaire
\fill[poporangeL] (-2.6,1.1) rectangle (2.6,2.3);
\node[font=\scriptsize, popdark] at (0,2.05) {milieu extracellulaire (liquide + macromolécules)};
\foreach \x in {-1.8,-0.9,0.3,1.5,2.1} \draw[fill=poporange, draw=popdark] (\x,1.55) circle (0.09);
% membrane avec invagination
\draw[very thick, popblue] (-2.6,1.1) -- (-0.7,1.1) .. controls (-0.55,0.35) and (0.15,0.35) .. (0.3,1.1) -- (2.6,1.1);
% vésicule qui se détache
\draw[fill=popcream, draw=popblue, thick] (-0.2,-0.35) circle (0.42);
\draw[fill=poporange, draw=popdark] (-0.2,-0.35) circle (0.09);
\node[font=\scriptsize, popblue] at (1.5,-0.35) {vésicule de pinocytose};
\draw[-{Stealth[length=2.5mm]}, thick, popblue] (-0.2,0.25) -- (-0.2,-0.75);
% cytoplasme
\fill[popblueL, opacity=0.45] (-2.6,-1.6) rectangle (2.6,0.28);
\node[font=\scriptsize, popdark] at (0,-1.35) {cytoplasme};
\end{scope}
% ===== Exocytose (droite) =====
\begin{scope}[shift={(8.2,0)}]
\node[font=\small\bfseries, popgreen!60!black] at (0,2.6) {Exocytose (sortie)};
\fill[poporangeL] (-2.6,1.1) rectangle (2.6,2.3);
\node[font=\scriptsize, popdark] at (0,2.05) {milieu extracellulaire};
\foreach \x in {-1.6,-0.4,0.8,1.9} \draw[fill=popgreen, draw=popdark] (\x,1.6) circle (0.09);
\draw[very thick, popgreen!60!black] (-2.6,1.1) -- (2.6,1.1);
% vésicule en fusion
\draw[fill=popgreenL, draw=popgreen!60!black, thick] (-0.9,0.62) circle (0.42);
\draw[fill=popgreen, draw=popdark] (-0.9,0.62) circle (0.09);
\draw[-{Stealth[length=2.5mm]}, thick, popgreen!60!black] (-0.9,0.1) -- (-0.9,1.35);
% vésicule en approche
\draw[fill=popgreenL, draw=popgreen!60!black, thick] (1.1,-0.5) circle (0.42);
\draw[fill=popgreen, draw=popdark] (1.1,-0.5) circle (0.09);
\node[font=\scriptsize, popgreen!50!black] at (1.1,-1.15) {vésicule de sécrétion};
\draw[-{Stealth[length=2.5mm]}, thick, popgreen!60!black] (1.1,-0.05) -- (1.1,0.85);
\fill[popblueL, opacity=0.45] (-2.6,-1.6) rectangle (2.6,0.2);
\node[font=\scriptsize, popdark] at (-1.5,-1.35) {cytoplasme};
\end{scope}
\end{tikzpicture}
\end{center}
\legende{Comparaison pinocytose / exocytose. À gauche, la membrane s'invagine et emprisonne une gouttelette de liquide extracellulaire dans une vésicule entrante. À droite, une vésicule de sécrétion migre vers la membrane, fusionne avec elle et libère son contenu (points verts) à l'extérieur. Dans les deux cas, le contenu transporté ne traverse jamais la membrane : il reste entouré d'une membrane.}
\end{popfigure}

\courssub{Reconnaître endocytose et exocytose sur une électronographie}

\begin{methode}[Identifier le sens du mouvement d'une vésicule]
Face à une électronographie montrant une cellule et des vésicules, raisonne en quatre étapes :
\begin{enumerate}
\item \textbf{Repère la membrane plasmique} et oriente-toi : où est l'extérieur, où est le cytoplasme ?
\item \textbf{Cherche les pseudopodes} ou une invagination contenant une particule solide volumineuse (bactérie, débris) : il s'agit d'une \cle{phagocytose} (endocytose). Une petite invagination contenant du liquide ou des macromolécules indique une \cle{pinocytose}.
\item \textbf{Cherche une vésicule en fusion avec la membrane}, s'ouvrant vers l'extérieur, avec un contenu (souvent dense aux électrons) en train d'être libéré : il s'agit d'une \cle{exocytose}.
\item \textbf{Conclus sur le sens du mouvement} : vésicule qui se forme à partir de la membrane et entre = endocytose ; vésicule qui fusionne avec la membrane et libère son contenu dehors = exocytose.
\end{enumerate}
\end{methode}

\begin{attention}[Les particules ne « traversent » jamais la membrane]
Erreur fréquente au Baccalauréat : écrire que la bactérie ou l'hormone « traverse » la membrane plasmique. C'est faux. Dans l'endocytose comme dans l'exocytose, la particule reste \textbf{toujours entourée d'une membrane} : celle de la vésicule, qui dérive de la membrane plasmique (endocytose) ou qui fusionne avec elle (exocytose). Il n'y a jamais de contact direct entre la particule transportée et le cytoplasme. Autre piège : ne pas oublier de préciser que ces deux processus \textbf{consomment de l'ATP} — ce sont des transports \emph{actifs}, qui exploitent la \cle{fluidité membranaire}.
\end{attention}

\begin{exoresolu}[Analyse d'une électronographie]
\textbf{Énoncé.} Sur une électronographie de cellule pancréatique, on observe de nombreuses vésicules à contenu dense, dont l'une est en train de fusionner avec la membrane plasmique, son contenu étant rejeté vers l'extérieur de la cellule. (a) Nommer le phénomène observé. (b) Expliquer son déroulement. (c) Citer un produit ainsi libéré par le pancréas.
\tcblower
\textbf{Solution.}
(a) Il s'agit de l'\cle{exocytose} : une vésicule intracellulaire fusionne avec la membrane plasmique et libère son contenu dans le milieu extracellulaire.

(b) Les protéines de sécrétion sont synthétisées dans le réticulum endoplasmique, conditionnées en vésicules par l'appareil de Golgi ; ces vésicules migrent vers la membrane plasmique, fusionnent avec elle (grâce à la fluidité membranaire) et s'ouvrent vers l'extérieur, libérant leur contenu. Le processus consomme de l'ATP.

(c) Le pancréas exocrine libère ainsi des \cle{enzymes digestives} (amylase, lipase, protéases) dans le duodénum ; le pancréas endocrine libère l'\cle{insuline} dans le sang par le même mécanisme.
\end{exoresolu}

\begin{savaistu}[Choléra et réhydratation orale : quand les échanges cellulaires déraillent]
Le choléra, qui sévit encore par épidémies dans certaines régions du Cameroun (Extrême-Nord, Littoral), est dû à une toxine produite par la bactérie \emph{Vibrio cholerae}. Cette toxine dérègle le transport d'ions (notamment \ce{Cl-} et \ce{Na+}) à travers la membrane des cellules intestinales : un flux massif d'ions vers la lumière de l'intestin entraîne par osmose une perte d'eau considérable — jusqu'à \SI{1}{L} par heure — provoquant une déshydratation mortelle. Le traitement repose sur les \cle{solutés de réhydratation orale} (SRO) : une solution d'eau, de sel et de glucose. Le glucose permet l'absorption du sodium par cotransport, et l'eau suit par osmose : une application directe, qui sauve des vies, des mécanismes étudiés dans cette leçon.
\end{savaistu}

\begin{aretenir}[Bilan de la section]
\begin{itemize}
\item Les particules volumineuses ne peuvent traverser la membrane : la cellule utilise sa \cle{fluidité membranaire} pour les transporter dans des \cle{vésicules}.
\item \cle{Endocytose} = entrée : \cle{phagocytose} (particules solides, pseudopodes, phagosome, digestion lysosomiale) et \cle{pinocytose} (gouttelettes de liquide, macromolécules dissoutes).
\item \cle{Exocytose} = sortie : fusion d'une vésicule de sécrétion avec la membrane et libération du contenu (enzymes, hormones, neuromédiateurs).
\item Endocytose et exocytose sont des processus \cle{actifs} : ils consomment de l'\cle{ATP}.
\item La particule transportée reste \textbf{toujours séparée du cytoplasme par une membrane}.
\end{itemize}
\end{aretenir}

\coursec{Importance et applications des échanges cellulaires}

Dans les sections précédentes, nous avons vu que la cellule échange avec son milieu non seulement de l'eau (osmose) et des substances dissoutes (diffusion simple, diffusion facilitée, transport actif), mais aussi des particules volumineuses grâce à l'\cle{endocytose} et à l'\cle{exocytose}. Ces mécanismes ne sont pas de simples curiosités de laboratoire : ils conditionnent chaque instant de la vie d'un organisme. Pourquoi une tige de salade flétrie redevient-elle croquante après un bain d'eau fraîche ? Pourquoi le poisson salé se conserve-t-il des mois sans réfrigérateur au marché de Mokolo ? Pourquoi le perfuse-t-on avec un sérum à 9 ‰ et non avec de l'eau pure ? C'est à toutes ces questions, qui relèvent de la \cle{limitation des conséquences des échanges cellulaires}, que cette dernière section répond.

\courssub{Importance des échanges cellulaires chez les végétaux}

\textbf{Absorption de l'eau et des sels minéraux par les racines.} Le sol, au voisinage des racines, est une solution très diluée : c'est un milieu \cle{hypotonique} par rapport au suc cellulaire des poils absorbants. L'eau pénètre donc dans les cellules racinaires par \cle{osmose}, à travers la membrane plasmique et les aquaporines. En revanche, certains ions minéraux (nitrates \ce{NO3-}, phosphates \ce{PO4^3-}, potassium \ce{K+}) sont souvent plus concentrés à l'intérieur de la racine que dans le sol : leur absorption ne peut alors pas se faire par diffusion. Elle s'effectue par \cle{transport actif}, grâce à des pompes membranaires consommant l'ATP fourni par la respiration cellulaire des racines. C'est pourquoi une racine asphyxiée (sol gorgé d'eau, sans \ce{O2}) cesse d'absorber les sels minéraux même si l'eau continue d'entrer par osmose : observation classique qui distingue nettement les deux mécanismes.

\textbf{La turgescence, squelette des plantes herbacées.} Une plante non lignifiée (salade, jeune tige de manioc, feuille de gombo) ne possède ni os ni tronc rigide : c'est la \cle{turgescence} de ses cellules qui la maintient dressée. Chaque cellule gorgée d'eau exerce sur sa paroi une pression de l'ordre de quelques bars (0,5 à 2 MPa) ; l'ensemble de ces cellules tendues confère à l'organe sa rigidité. Lorsque l'eau manque, la pression de turgescence chute : les cellules deviennent flasques (état de \cle{flétrissure}) et la plante se couche. Un arrosage restaure l'osmose entrante et la plante se redresse en quelques heures.

\begin{propriete}[Rôle physiologique de la turgescence]
La turgescence des cellules assure le maintien et la rigidité des organes végétaux non lignifiés (tiges herbacées, feuilles, pétales) et conditionne le fonctionnement des stomates, donc les échanges gazeux foliaires (entrée de \ce{CO2} pour la photosynthèse, sortie de \ce{O2} et de vapeur d'eau).
\end{propriete}

\textbf{Ouverture et fermeture des stomates.} Le stomate est un orifice de l'épiderme foliaire bordé de deux \cle{cellules stomatiques} (cellules de garde) en forme de haricot. Le mécanisme est un chef-d'œuvre d'hydraulique cellulaire :
\begin{itemize}
\item En présence de lumière, les cellules stomatiques accumulent activement des ions \ce{K+} (transport actif \ce{H+}/\ce{K+} couplé à une \ce{H+}-ATPase). Leur potentiel osmotique diminue (devient plus négatif), l'eau entre par osmose : les cellules se \cle{turgescent}. Leur paroi étant plus mince du côté de l'ostiole, elles se bombent et s'écartent : le stomate \textbf{s'ouvre}.
\item À l'obscurité ou en cas de sécheresse (signal \cle{ABA} -- acide abscissique), les ions \ce{K+} ressortent (canaux sortants), l'eau suit par osmose, les cellules se \cle{flétrissent} et se rapprochent : le stomate \textbf{se ferme}, limitant les pertes d'eau par transpiration.
\end{itemize}

\begin{popfigure}
\begin{center}
\begin{tikzpicture}[scale=1.0]
% ---- Stomate ouvert ----
% Cellule supérieure (rotée autour de son centre)
\draw[fill=popgreenL, draw=popgreen!70!black, thick, rotate around={25:(-2.6,0.55)}] (-2.6,0.55) ellipse (1.15cm and 0.42cm);
% Cellule inférieure
\draw[fill=popgreenL, draw=popgreen!70!black, thick, rotate around={-25:(-2.6,-0.55)}] (-2.6,-0.55) ellipse (1.15cm and 0.42cm);
% Ostiole
\draw[fill=white, draw=popdark, thick] (-2.6,0) ellipse (0.28cm and 0.5cm);
\node[popdark, font=\small\bfseries] at (-2.6,-1.6) {Stomate ouvert};
\node[popmuted, font=\footnotesize, align=center] at (-2.6,-2.15) {Cellules turgescentes\\(entrée d'eau par osmose)};
\draw[->, popblue, thick] (-4.6,0.9) -- (-3.6,0.75) node[midway, above, font=\scriptsize]{Entrée d'eau};
% ---- Stomate fermé ----
\draw[fill=popgreenL!60, draw=popgreen!70!black, thick] (2.35,0) ellipse (0.42cm and 1.05cm);
\draw[fill=popgreenL!60, draw=popgreen!70!black, thick] (3.05,0) ellipse (0.42cm and 1.05cm);
\draw[draw=popdark, thick] (2.7,-0.85) -- (2.7,0.85);
\node[popdark, font=\small\bfseries] at (2.7,-1.6) {Stomate fermé};
\node[popmuted, font=\footnotesize, align=center] at (2.7,-2.15) {Cellules flétries\\(sortie d'eau par osmose)};
\draw[->, poporange, thick] (4.6,0.9) -- (3.7,0.75) node[midway, above, font=\scriptsize]{Sortie d'eau};
\end{tikzpicture}
\end{center}
\legende{Schéma fonctionnel d'un stomate vu de face. À gauche : les deux cellules stomatiques gorgées d'eau se bombent et libèrent l'ostiole (ouverture). À droite : déshydratées, elles s'affaissent et obturent l'ostiole (fermeture). Le sens des flux d'eau est contrôlé par le transport actif des ions \ce{K+}.}
\end{popfigure}

\courssub{Importance des échanges cellulaires chez les animaux}

\textbf{Maintien du milieu intérieur.} Chez l'animal, chaque cellule baigne dans le \cle{milieu intérieur} (lymphe et plasma), dont la composition doit rester stable : c'est l'\cle{homéostasie}. Les globules rouges (hématies), par exemple, ne conservent leur intégrité que dans un plasma \cle{isotonique} : placés dans un milieu hypotonique, ils éclatent (\cle{hémolyse}) ; dans un milieu hypertonique, ils se \cle{crénèlent}. L'organisme limite en permanence les conséquences des échanges hydriques en régulant la concentration du sang (osmorégulation rénale, soif, ADH).

\textbf{Absorption intestinale et réabsorption rénale.} Au niveau de l'intestin grêle, le glucose et les acides aminés sont absorbés par les cellules des villosités grâce à un \cle{transport actif secondaire} (co-transport \ce{Na+}/glucose, symporteur SGLT1), couplé au gradient de sodium maintenu par la \ce{Na+/K+}-ATPase basolatérale. Cela permet l'absorption même lorsque la concentration luminale devient plus faible que dans le sang : sans ce mécanisme, une partie des nutriments serait perdue dans les selles. Au niveau du rein, les néphrons \cle{réabsorbent} l'eau (osmose, sous contrôle de l'ADH/vasopressine au niveau du tube collecteur) et les solutés utiles (glucose, ions \ce{Na+}, \ce{Cl-}, \ce{HCO3-}) par transport actif, tout en éliminant les déchets (urée, créatinine) : c'est la clé de la constance du milieu intérieur.

\textbf{Défense immunitaire et sécrétions.} La \cle{phagocytose}, forme d'endocytose, est assurée par les leucocytes phagocytaires (polynucléaires neutrophiles, macrophages) qui englobent et digèrent les microbes dans des phagolysosomes : c'est un rempart essentiel contre les infections, par exemple lors de la réaction inflammatoire d'une plaie. Inversement, l'\cle{exocytose} permet aux cellules glandulaires de libérer leurs produits de sécrétion : enzymes digestives (trypsinogène, lipase) par le pancréas exocrine, \cle{insuline} par les cellules \cle{$\beta$ des îlots de Langerhans}, anticorps par les plasmocytes, neurotransmetteurs (acétylcholine, dopamine) par les neurones. Sans exocytose, ni digestion, ni régulation de la glycémie, ni communication nerveuse ne seraient possibles.

\courssub{Applications médicales}

\begin{exemplebox}[Le sérum physiologique, une solution isotonique]
Le \cle{sérum physiologique} (NaCl 0,9 \%) contient \SI{9}{g} de NaCl par litre d'eau, soit une concentration en particules équivalente à celle du plasma sanguin (environ 300 mosmol/L) : il est \cle{isotonique}. Perfusé à un patient, il n'entraîne ni hémolyse (l'eau n'entre pas en excès dans les hématies) ni crénelure (l'eau n'en sort pas). De même, le \cle{sérum glucosé à 5 \%} (\SI{50}{g} de glucose par litre, soit \SI{278}{mmol/L}) est isotonique au plasma au moment de la perfusion (le glucose est ensuite métabolisé). À l'inverse, une perfusion d'eau pure provoquerait l'éclatement massif des hématies par osmose entrante : c'est une erreur mortelle que la connaissance de l'osmose permet d'éviter.
\end{exemplebox}

\textbf{Les solutés de réhydratation orale (SRO).} Lors d'une diarrhée aiguë (choléra, gastro-entérite virale), l'organisme perd massivement de l'eau et des électrolytes (\ce{Na+}, \ce{K+}, \ce{HCO3-}, \ce{Cl-}). Le SRO (formulation OMS à faible osmolarité : \SI{75}{mmol/L} NaCl, \SI{75}{mmol/L} glucose, \SI{20}{mmol/L} KCl, \SI{10}{mmol/L} citrate) exploite directement le cours : le glucose est absorbé par la cellule intestinale (entérocye) grâce à un \cle{co-transport actif \ce{Na+}/glucose} (SGLT1) à la membrane apicale ; le sodium entre dans la cellule, l'eau suit ce mouvement de solutés par \cle{osmose} (via aquaporines et jonctions serrées). On réhydrate ainsi le malade par voie orale, sans perfusion veineuse, sauvant des vies là où le matériel de perfusion fait défaut.

\begin{savaistu}[Les SRO au Cameroun]
Dans les centres de santé camerounais, des hôpitaux de district de Yaoundé et Douala aux aires de santé de l'Extrême-Nord, les sachets de solutés de réhydratation orale, distribués gratuitement ou à très bas coût dans le cadre des programmes de lutte contre les maladies diarrhéiques (PNLP, appui UNICEF/OMS), ont considérablement réduit la mortalité infantile due aux diarrhées, autrefois l'une des premières causes de décès des enfants de moins de cinq ans. Une simple solution d'eau, de sel et de sucre, correctement dosée, sauve chaque année des milliers de vies : belle illustration de la puissance des connaissances sur l'osmose et le transport actif.
\end{savaistu}

\textbf{Le rein artificiel (hémodialyse).} Chez les insuffisants rénaux terminaux, le sang n'est plus épuré. La \cle{dialyse} (rein artificiel) fait circuler le sang du patient d'un côté d'une membrane semi-perméable (fibres creuses en polysulfone), et un liquide de dialyse de composition contrôlée (électrolytes, bicarbonate, glucose) de l'autre, en contre-courant. Les déchets (urée, créatinine, urate, ions \ce{K+} en excès, excès d'eau par ultrafiltration) diffusent du sang vers le liquide, du milieu le plus concentré vers le moins concentré (loi de Fick), tandis que les grosses molécules utiles (protéines plasmatiques > 20 kDa) et les cellules sanguines sont retenues. C'est une application directe du principe de la \cle{dialyse} expérimentale (membrane de Dutrochet/Pfeffer) : le patient doit cependant subir plusieurs séances par semaine (généralement 3 séances de 4 h), faute de quoi les toxines s'accumulent (urée > 1,5 g/L risque encéphalopathie).

\courssub{Applications alimentaires et domestiques}

\textbf{Conservation des aliments par le sel et le sucre.} Le poisson salé (hareng, morue), la viande fumée-salée (kpem, viande de brousse séchée), les confitures (\SI{60}{--}{65}{\%} sucre) et les fruits confits se conservent longtemps sans chaîne du froid. Le mécanisme est osmotique : le sel (saumure \SI{>}{15}{\%} NaCl) ou le sucre en forte concentration crée à la surface et à l'intérieur de l'aliment (par diffusion) un milieu \cle{hypertonique} (activité de l'eau $a_w < 0,7$). Les cellules microbiennes (bactéries de putréfaction, moisissures, levures) qui s'y trouvent perdent leur eau par osmose et se \cle{plasmolysent} : leur métabolisme est bloqué (inhibition enzymatique), elles ne peuvent plus se multiplier. L'aliment est ainsi protégé de la putréfaction.

\begin{attention}[Le sel ne « tue » pas directement les microbes]
Erreur fréquente au Bac : affirmer que le sel ou le sucre « tue » les microbes comme le ferait un poison (bactéricide). En réalité, le sel provoque surtout leur \cle{déshydratation par osmose} (plasmolyse), ce qui \textbf{bloque leur développement et leur multiplication} (effet bactériostatique/fongistatique). Beaucoup de microbes (spores, bactéries halotolérantes comme \emph{Staphylococcus aureus}) survivent à l'état dormant et peuvent reprendre leur activité si le milieu redevient dilué (réhydratation). Rédige donc : « le milieu hypertonique provoque la plasmolyse des cellules microbiennes, ce qui inhibe leur prolifération ».
\end{attention}

\textbf{Redressement des légumes flétris.} Une salade ou des feuilles de ngòn flétries plongées dans l'eau fraîche (milieu hypotonique) redeviennent croquantes en 30 à 60 min : l'eau entre dans les cellules par osmose et restaure leur \cle{turgescence}. C'est la même raison pour laquelle les fleuristes conservent les fleurs coupées dans l'eau (maintien de la turgescence des cellules du pétale et du pédoncule).

\textbf{Le danger de l'excès d'engrais (brûlure racinaire).} Un apport massif d'engrais minéraux (NPK, urée) au pied d'une plante rend la solution du sol \cle{hypertonique} par rapport aux cellules racinaires (potentiel osmotique $\Psi_s$ très négatif) : l'eau sort alors des racines par osmose (exosmose), les cellules se plasmolysent et la plante se dessèche même en terrain humide. C'est la « \cle{brûlure par les engrais} », bien connue des maraîchers qui fertilisent trop près des plants de tomates ou de piments. De même, il ne faut jamais arroser les plantes avec de l'eau salée (eau de mer, eau de forage saumâtre).

\begin{exoresolu}[Pourquoi un arrosage à l'eau salée tue-t-il une plante verte ?]
Énoncé. Un jardinier arrose par erreur ses plants de laitue avec de l'eau de mer (environ \SI{35}{g} de sels par litre, soit \SI{\approx 600}{mosmol/L}). En quelques jours, les plants flétrissent puis meurent, alors que le sol reste humide. Expliquer ce phénomène en précisant le mécanisme cellulaire mis en jeu.
\tcblower
Solution détaillée. L'eau de mer, très chargée en sels dissous (\ce{Na+}, \ce{Cl-}, \ce{Mg^2+}, \ce{SO4^2-}), constitue un milieu \cle{hypertonique} par rapport au suc cellulaire des cellules des racines et des feuilles (environ \SI{300}{mosmol/L}). Le gradient de potentiel hydrique ($\Psi_w = \Psi_s + \Psi_p$) est donc inversé : le potentiel hydrique du sol est très inférieur (plus négatif) à celui de la cellule. L'eau \textbf{sort} des cellules par \cle{osmose}, à travers la membrane plasmique semi-perméable, vers la solution du sol. Les cellules végétales perdent leur pression de turgescence ($\Psi_p \to 0$) : la vacuole rétrécit, la membrane plasmique se décolle de la paroi pecto-cellulosique rigide, c'est la \cle{plasmolyse}. Les feuilles, privées du soutien mécanique que leur confère la turgescence, flétrissent ; les stomates se ferment par manque de turgescence des cellules de garde, stoppant les échanges gazeux (photosynthèse nulle). La plasmolyse prolongée devenant irréversible (dénaturation des protéines, rupture de la membrane), les cellules meurent et la plante entière périt, bien que le sol soit humide : l'eau présente est \emph{indisponible} pour la plante car elle ne peut pas l'absorber contre le gradient osmotique (il faudrait un transport actif d'eau, qui n'existe pas).
\end{exoresolu}

\courssub{Synthèse générale : comment la cellule limite les conséquences des échanges}

Tout au long de cette leçon, un même fil conducteur apparaît : livrée à elle-même, la cellule subirait passivement les lois des gradients physico-chimiques (éclatement en milieu hypotonique, plasmolyse en milieu hypertonique, épuisement ou accumulation toxique de solutés). Elle \cle{limite les conséquences} de ces échanges grâce à trois dispositifs complémentaires :
\begin{enumerate}
\item la \cle{sélectivité de la membrane plasmique} (bicouche phospholipidique + protéines intégrales), barrière semi-perméable qui autorise certains passages (eau via aquaporines, gaz \ce{O2}/\ce{CO2}, petites molécules apolaires) et en interdit d'autres (ions, molécules polaires, macromolécules) sans transporteurs spécifiques ;
\item les \cle{transports actifs} (primaires \ce{Na+/K+}-ATPase, \ce{H+}-ATPase, \ce{Ca^2+}-ATPase ; secondaires symport/antiport), qui dépensent de l'ATP (énergie chimique) pour déplacer des solutés contre leur gradient électrochimique et maintenir des concentrations internes différentes du milieu extérieur (homéostasie ionique, charge nutritionnelle) ;
\item chez les végétaux, la \cle{paroi pecto-cellulosique rigide} (cellulose, hémicelluloses, pectines, protéines structurales), qui oppose une résistance mécanique ($\Psi_p > 0$) à l'entrée d'eau en milieu hypotonique et transforme un danger potentiel (l'osmose entrante illimitée) en atout physiologique : la turgescence (squelette hydraulique, croissance par allongement cellulaire).
\end{enumerate}

\begin{popfigure}
\begin{center}
\begin{tikzpicture}[mindmap, grow cyclic, every node/.style={concept, font=\small, text=white},
root concept/.append style={concept color=popblue, font=\small\bfseries, minimum size=3.1cm},
level 1/.append style={level distance=3.6cm, sibling angle=72},
level 1 concept/.append style={font=\scriptsize\bfseries, minimum size=2.2cm},
every annotation/.style={font=\scriptsize}]
\node[root concept]{Échanges cellulaires}
  child[concept color=popgreen] { node[align=center]{Osmose\\ (eau, passif,\\ aquaporines)} }
  child[concept color=poporange] { node[align=center]{Diffusion simple\\ (\ce{O2}, \ce{CO2}, lipophiles)} }
  child[concept color=poppurple] { node[align=center]{Diffusion facilitée\\ (canaux, transporteurs\\ passifs, pas d'ATP)} }
  child[concept color=poppink] { node[align=center]{Transport actif\\ (pompes, co-transport\\ ATP requis)} }
  child[concept color=popgold] { node[align=center]{Endocytose / Exocytose\\ (vésicules, particules\\ ATP requis)} };
\end{tikzpicture}
\end{center}
\legende{Carte mentale de synthèse des modes d'échanges entre la cellule et le milieu ambiant. Les transports passifs (osmose, diffusions) suivent le gradient électrochimique et ne consomment pas d'ATP directement ; le transport actif et les mouvements membranaires vésiculaires (endocytose, exocytose) consomment de l'ATP.}
\end{popfigure}

\begin{methode}[Rédiger une réponse de synthèse au Bac]
Face à une question du type « Expliquer le phénomène observé » ou « Interpréter les résultats de l'expérience de... », procède toujours en trois étapes rigoureuses :
\begin{enumerate}
\item \textbf{Partir du fait observé} : décris précisément le phénomène (cellule gonflée/plasmolysée, plante flétrie/dressée, globule rouge crénelé/lyse, courbe de saturation d'un transporteur...) et les conditions expérimentales (nature des milieux, concentrations, température, présence/absence d'inhibiteur métabolique comme le cyanure ou le DNP).
\item \textbf{Identifier le mécanisme d'échange impliqué} : osmose, diffusion simple ou facilitée, transport actif (préciser primaire ou secondaire), endocytose (phagocytose/pinocytose) ou exocytose. Nomme-le explicitement avec le vocabulaire du programme.
\item \textbf{Justifier par les gradients et l'énergie} : précise le sens du gradient de concentration (ou électrochimique) pour le soluté ou l'eau (milieu hypotonique/isotonique/hypertonique), le sens du mouvement net qui en résulte, et indique si le mécanisme consomme de l'ATP (transport actif, endo/exocytose, maintien du gradient \ce{Na+}) ou non (transports passifs). Conclus en reliant le mécanisme cellulaire au phénomène observé à l'échelle de l'organe ou de l'organisme (turgescence, absorption, sécrétion, défense).
\end{enumerate}
\end{methode}

\begin{aretenir}[Bilan de la leçon et ouverture]
Les échanges entre la cellule et le milieu ambiant — eau par osmose, solutés par diffusions (simple, facilitée) et transport actif (primaire, secondaire), particules par endocytose (phagocytose, pinocytose) et exocytose — sont contrôlés et limités par la membrane plasmique (sélectivité, protéines de transport), les pompes à ATP et, chez le végétal, la paroi rigide. Ils conditionnent le fonctionnement de l'organisme entier : la \textbf{nutrition} (absorption racinaire minérale, absorption intestinale glucidique/protidique), l'\textbf{excrétion/osmorégulation} (rein, transpiration stomatique), la \textbf{défense} (phagocytose, présentation d'antigènes) et la \textbf{communication intercellulaire} (exocytose des hormones, neurotransmetteurs, cytokines). Leurs applications irriguent la médecine (perfusions isotoniques, SRO, hémodialyse, conservation d'organes) comme la vie quotidienne (conservation des aliments par sel/sucre/séchage, jardinage, aquariophilie). La suite du programme (Module II : Nutrition, Module III : Relation) montrera comment ces échanges cellulaires s'intègrent dans les grandes fonctions intégrées de l'organisme.
\end{aretenir}

\newpage

\coursec{Activité d'intégration}

\coursec{Limitation des conséquences liées aux échanges d'eau, de substances dissoutes et de particules entre la cellule et le milieu ambiant}

\courssub{Objectifs de l'activité}

À l'issue de cette activité d'intégration, tu seras capable de :

\begin{itemize}
\item \cle{identifier} sur des microphotographies les différents états d'une cellule végétale (turgescence, plasmolyse) placée dans des solutions de concentrations différentes ;
\item \cle{interpréter} les résultats d'une expérience de type Dutrochet et en dégager la notion d'osmose ;
\item \cle{exploiter} une courbe d'absorption d'ions minéraux pour distinguer diffusion simple, diffusion facilitée et transport actif ;
\item \cle{mobiliser} les notions d'osmose, de transport passif et de transport actif pour expliquer deux situations concrètes de la vie courante (flétrissement de plants de tomates, conservation du poisson salé) ;
\item \cle{rédiger} un conseil argumenté, fondé sur des principes scientifiques, à destination d'un usager (maraîcher).
\end{itemize}

\courssub{Situation}

\textbf{Le mystère des plants de tomates flétris et du poisson salé.}

Monsieur Kamdem est maraîcher à Bafou, dans la région de l'Ouest du Cameroun. Le lundi matin, il épand autour de ses plants de tomates une forte dose d'engrais minéral NPK (environ \SI{500}{g} par m\textsuperscript{2}, soit 5 fois la dose recommandée), puis il arrose. Le mardi matin, stupeur : les plants sont flétris, les feuilles pendantes, alors que le sol est encore humide. Aucun insecte ravageur n'est observé.

Intrigué, son neveu, élève en Terminale D, prélève des racines et réalise au laboratoire du lycée les observations et expériences suivantes.

\textbf{Document 1 — Microphotographies de cellules de l'épiderme racinaire} (grossissement $\times$ 400) :

\begin{itemize}
\item \textbf{Photo A} (racines prélevées le lundi, avant l'épandage) : les cellules sont gonflées, la paroi est tendue, le cytoplasme et la membrane plasmique sont plaqués contre la paroi, la vacuole occupe presque tout le volume cellulaire.
\item \textbf{Photo B} (racines prélevées le mardi, après l'épandage) : la membrane plasmique s'est décollée de la paroi pectocellulosique, le protoplaste est rétracté au centre de la cellule, la vacuole a perdu une grande partie de son eau.
\end{itemize}

\textbf{Document 2 — Expérience de type Dutrochet.} On remplit une poche de membrane de collodion (membrane perméable à l'eau et aux petites molécules, imperméable aux grosses molécules) avec de l'eau du sol prélevée autour des racines après épandage (solution très concentrée en sels minéraux). On plonge cette poche, surmontée d'un tube capillaire, dans un bécher d'eau distillée. Résultat : le niveau du liquide \textbf{monte} dans le tube capillaire et se stabilise à $h = \SI{12}{cm}$ au bout de 2 heures. Témoin : une poche remplie d'eau distillée plongée dans l'eau distillée ne montre aucune variation de niveau.

\textbf{Document 3 — Absorption du potassium par les racines.} On mesure la vitesse d'absorption des ions $K^{+}$ par des racines de tomate placées dans des solutions de concentrations externes croissantes, en l'absence puis en présence de cyanure (inhibiteur de la respiration cellulaire, donc de la production d'ATP) :

\begin{center}
\begin{tabular}{|l|c|c|c|c|}
\hline
\rowcolor{popblueL}
Concentration externe en $K^{+}$ (mmol.L$^{-1}$) & 0,5 & 1 & 2 & 4 \\
\hline
Vitesse d'absorption sans cyanure (µmol.h$^{-1}$.g$^{-1}$) & 2,0 & 3,5 & 5,0 & 5,2 \\
\hline
Vitesse d'absorption avec cyanure (µmol.h$^{-1}$.g$^{-1}$) & 0,5 & 1,0 & 2,0 & 4,0 \\
\hline
\end{tabular}
\end{center}

\textbf{Document 4 — La recette du poisson salé.} Au marché de Bafou, la mère du jeune élève vend du poisson salé séché. Recette : couvrir les poissons frais d'une épaisse couche de gros sel pendant plusieurs jours, puis les faire sécher au soleil. Le poisson ainsi traité se conserve plusieurs semaines sans réfrigérateur, alors que le poisson frais se gâte en moins de 24 heures à la température ambiante.

\courssub{Consignes}

\begin{enumerate}
\item Identifier, en justifiant, l'état des cellules observées sur la photo A et sur la photo B (document 1). Nommer précisément chaque phénomène.
\item À l'aide du document 2, interpréter l'expérience de type Dutrochet : quel phénomène met-elle en évidence ? Définir ce phénomène et préciser le rôle de la membrane de collodion.
\item Expliquer, à l'échelle cellulaire, pourquoi les plants de tomates ont flétri alors que le sol était humide. Tu feras le lien entre la concentration de la solution du sol et le sens du mouvement de l'eau.
\item Tracer sur ton cahier le graphe de la vitesse d'absorption du $K^{+}$ en fonction de la concentration externe, sans et avec cyanure (document 3). Pour chaque courbe, décrire son allure et identifier le ou les mécanismes de transport mis en jeu (diffusion simple, diffusion facilitée, transport actif). Justifier en t'appuyant sur l'effet du cyanure et sur la saturation éventuelle.
\item Expliquer, en utilisant les mêmes principes osmotiques, pourquoi le salage permet de conserver le poisson (document 4). Que deviennent les micro-organismes de décomposition à la surface du poisson salé ?
\item Rédiger un conseil argumenté (10 à 15 lignes) à l'attention de M. Kamdem, lui expliquant l'origine du problème et lui proposant des solutions concrètes (dose d'engrais, arrosage abondant, délai avant récupération des plants).
\end{enumerate}

\courssub{Ressources à mobiliser}

\begin{aretenir}[Boîte à outils de la leçon]
\begin{itemize}
\item \textbf{Osmose} : diffusion de l'eau (solvant) à travers une membrane semi-perméable, du milieu le moins concentré en solutés (hypotonique) vers le milieu le plus concentré (hypertonique).
\item \textbf{Dialyse} : diffusion des solutés (substances dissoutes) à travers une membrane semi-perméable, du milieu le plus concentré vers le milieu le moins concentré (dans le sens du gradient de concentration). La membrane de collodion, perméable aux petites molécules, permet la dialyse.
\item \textbf{Cellule végétale} : milieu hypotonique $\Rightarrow$ \cle{turgescence} (entrée d'eau, protoplaste plaqué contre la paroi) ; milieu hypertonique $\Rightarrow$ \cle{plasmolyse} (sortie d'eau, décollement de la membrane plasmique de la paroi, rétraction du protoplaste).
\item \textbf{Cellule animale (ex. hématie)} : milieu hypotonique $\Rightarrow$ turgescence puis \cle{hémolyse} (éclatement par absence de paroi) ; milieu hypertonique $\Rightarrow$ \cle{crénelure} (sortie d'eau, rétrécissement, aspect ébouriffé).
\item \textbf{Transport passif} (sans ATP, dans le sens du gradient de concentration) :
      \begin{itemize}
      \item \cle{Diffusion simple} : proportionnelle à la concentration, non saturable, pas de protéine transporteur (ex. $O_2$, $CO_2$, lipophiles).
      \item \cle{Diffusion facilitée} : via protéines canaux ou transporteurs (glucose, ions), \textbf{saturable} (plateau), non sensible au cyanure.
      \end{itemize}
\item \textbf{Transport actif} (avec ATP, souvent contre le gradient) : \textbf{saturable} (plateau), \textbf{inhibé par le cyanure} (blocage respiration $\Rightarrow$ pas d'ATP). Ex. : pompe $Na^+/K^+$, absorption $K^+$ racinaire.
\item \textbf{Échanges de particules} : \cle{endocytose} (phagocytose : particules solides ; pinocytose : fluides) et \cle{exocytose} (sécrétion, rejet). Mécanismes vésiculaires coûtant de l'ATP.
\item \textbf{Expérience de Dutrochet} : la montée du liquide dans le tube traduit l'entrée d'eau dans la poche par osmose (milieu extérieur hypotonique $\rightarrow$ milieu intérieur hypertonique).
\end{itemize}
\end{aretenir}

\courssub{Démarche et corrigé commenté}

\begin{exoresolu}[Corrigé de l'activité d'intégration]
\textbf{Énoncé.} Voir situation et consignes 1 à 6 ci-dessus.

\tcblower

\textbf{1. Identification des états cellulaires.}

Photo A : les cellules sont gonflées, la vacuole est volumineuse et plaque le cytoplasme contre la paroi pectocellulosique : c'est l'état de \cle{turgescence}. Avant l'épandage, la solution du sol est hypotonique par rapport au suc vacuolaire : l'eau entre dans les cellules par osmose, assurant la rigidité des tissus.

Photo B : la membrane plasmique s'est décollée de la paroi et le protoplaste est rétracté au centre de la cellule : c'est la \cle{plasmolyse}. Après l'épandage massif d'engrais, la solution du sol est devenue hypertonique par rapport au suc vacuolaire : l'eau est sortie des cellules racinaires par osmose.

\textbf{2. Interprétation de l'expérience de type Dutrochet.}

La montée du liquide de $h = \SI{12}{cm}$ dans le tube capillaire montre qu'un volume d'eau est entré dans la poche, augmentant la pression interne (pression osmotique), alors qu'aucun mouvement n'a lieu dans le témoin (isotonie). L'eau a diffusé du bécher (eau distillée, milieu \textbf{hypotonique}) vers la poche (solution du sol concentrée, milieu \textbf{hypertonique}) à travers la membrane de collodion. Cette membrane joue le rôle de \cle{membrane semi-perméable} : elle est perméable aux molécules d'eau (et aux petits solutés) mais imperméable aux gros solutés (macromolécules, ions en grande quantité). Le phénomène mis en évidence est l'\cle{osmose}.

\textit{Précision :} La membrane de collodion laissant passer les petites molécules, une \cle{dialyse} (diffusion des solutés du milieu hypertonique vers l'hypotonique) se produit simultanément, mais l'expérience de Dutrochet est conçue pour visualiser spécifiquement le flux d'eau (osmose) par la variation de niveau dans le capillaire.

\textbf{3. Explication du flétrissement des plants.}

L'épandage de 5 fois la dose recommandée d'engrais NPK a rendu la solution du sol fortement \cle{hypertonique} par rapport au suc vacuolaire des cellules de l'épiderme racinaire. Le gradient de concentration en eau s'est inversé : le potentiel hydrique du sol est devenu plus bas (plus négatif) que celui de la cellule. Par osmose, l'eau est \textbf{sortie} massivement des cellules racinaires vers le sol. Les cellules se sont plasmolysées (photo B), les tissus ont perdu leur rigidité mécanique (perte de turgescence), d'où le flétrissement des parties aériennes, bien que le sol soit humide. L'eau du sol est présente mais \emph{indisponible} pour la plante car retenue par les sels : c'est une \textbf{sécheresse physiologique} (ou sécheresse saline).

\textbf{4. Exploitation du document 3.}

\begin{popfigure}
\begin{tikzpicture}
\begin{axis}[width=0.9\linewidth, height=7cm,
xlabel={Concentration externe en $K^{+}$ (mmol.L$^{-1}$)},
ylabel={Vitesse d'absorption (µmol.h$^{-1}$.g$^{-1}$)},
xmin=0, xmax=4.5, ymin=0, ymax=6,
legend pos=north west, grid=major, grid style={popline!40},
tick label style={/pgf/number format/use comma}]
\addplot[popblue, thick, mark=*, mark size=1.5pt] coordinates {(0.5,2.0) (1,3.5) (2,5.0) (4,5.2)};
\addlegendentry{Sans cyanure (Total)}
\addplot[poporange, thick, mark=square*, mark size=1.5pt] coordinates {(0.5,0.5) (1,1.0) (2,2.0) (4,4.0)};
\addlegendentry{Avec cyanure (Passif)}
\end{axis}
\end{tikzpicture}
\legende{Vitesse d'absorption du $K^{+}$ par des racines de tomate en fonction de la concentration externe, sans et avec cyanure. La courbe « Sans cyanure » représente le transport total (passif + actif). La courbe « Avec cyanure » isole la composante passive (diffusion simple).}
\end{popfigure}

\begin{itemize}
\item \textbf{Courbe « Avec cyanure » (orange)} : la vitesse est proportionnelle à la concentration externe (alignement quasi parfait sur une droite passant par l'origine : $v \approx 1 \times [K^+]_{ext}$), sans plateau de saturation jusqu'à 4 mmol.L$^{-1}$. Ce mécanisme ne nécessite pas d'ATP (insensible au cyanure) et ne fait pas intervenir de protéines saturables : c'est une \cle{diffusion simple} (transport passif) des ions $K^+$ dans le sens de leur gradient électrochimique.
\item \textbf{Courbe « Sans cyanure » (bleue)} : la vitesse augmente rapidement puis atteint un \textbf{plateau de saturation} vers 5,2 µmol.h$^{-1}$.g$^{-1}$ dès 2 mmol.L$^{-1}$. Cela traduit l'intervention de \cle{protéines transporteurs} en nombre limité (saturation). Ce mécanisme est \textbf{fortement inhibé par le cyanure} (la différence entre les deux courbes disparaît à saturation) : il consomme de l'ATP fourni par la respiration mitochondriale. Il s'agit donc d'un \cle{transport actif}, permettant l'absorption du $K^+$ \emph{contre} son gradient de concentration (accumulation dans la cellule).
\end{itemize}

La composante active (différence des deux courbes) est saturée dès 2 mmol.L$^{-1}$ ; la composante passive (diffusion simple), non saturable, devient majoritaire à forte concentration externe.

\textbf{5. Conservation du poisson salé.}

Le sel déposé en couche épaisse crée à la surface du poisson un milieu fortement \cle{hypertonique} (saumure saturée $\approx$ 5 mol.L$^{-1}$ NaCl). Deux conséquences osmotiques simultanées :
\begin{itemize}
\item \textbf{Déshydratation de la chair} : l'eau contenue dans les cellules musculaires du poisson (milieu hypotonique par rapport à la saumure) diffuse vers l'extérieur par osmose. Le poisson « rend » son eau, perd de sa masse, sa texture se raffermit.
\item \textbf{Inhibition microbienne} : les micro-organismes de décomposition (bactéries, levures, moisissures) déposés à la surface se retrouvent dans un milieu hypertonique extrême. L'eau quitte leur cytoplasme par osmose : leurs cellules subissent une \textbf{plasmolyse} (bactéries à paroi) ou une \textbf{crénelure} (levures), leur métabolisme s'arrête, elles meurent ou entrent en dormance.
\end{itemize}
Privé d'eau libre ($a_w$ très bas) et de flore microbienne active, le poisson se conserve plusieurs semaines. C'est exactement le même principe biophysique que le flétrissement des tomates : \textbf{l'eau quitte toujours le milieu hypotonique vers le milieu hypertonique}.

\textbf{6. Conseil argumenté à M. Kamdem.}

« Monsieur Kamdem, vos plants ont flétri parce que la dose d'engrais épandue (5 fois la dose normale, soit \SI{500}{g.m^{-2}} au lieu de \SI{100}{g.m^{-2}}) a rendu la solution du sol plus concentrée en sels minéraux que le suc vacuolaire de vos racines. Par osmose, l'eau est sortie des cellules racinaires au lieu d'y entrer : vos plants subissent une \textbf{sécheresse physiologique} (plasmolyse généralisée) alors que le sol est humide.

Pour sauver votre culture : (1) \textbf{Arrosez immédiatement et abondamment} (plusieurs arrosages espacés de quelques heures) avec de l'eau non salée pour \textbf{lessiver} l'excès d'engrais en profondeur, hors de la zone racinaire. La solution du sol redeviendra hypotonique, l'eau rentrera dans les racines (déplasmolyse) et la turgescence reviendra en 24 à 48 h si les membranes ne sont pas irréversiblement endommagées. (2) À l'avenir, \textbf{respectez scrupuleusement la dose recommandée} (\SI{100}{g.m^{-2}} pour le NPK 20-10-10) et \textbf{fractionnez les apports} (ex. 2 fois \SI{50}{g.m^{-2}} à 15 jours d'intervalle). (3) \textbf{N'épandez jamais l'engrais au contact direct des tiges ou des racines} : enfouissez-le légèrement à 5-10 cm du collet. Un sol bien drainé et riche en matière organique (compost) tamponne aussi les variations de salinité. »
\end{exoresolu}

\courssub{Bilan de l'activité}

Cette activité a permis de mettre en évidence que :

\begin{itemize}
\item un même principe biophysique — l'\cle{osmose} (mouvement de l'eau du milieu hypotonique vers le milieu hypertonique à travers une membrane semi-perméable) — explique des phénomènes en apparence très différents : le flétrissement de plants de tomates surfertilisés (agriculture) et la conservation du poisson salé (technique alimentaire traditionnelle) ;
\item les observations microscopiques (turgescence, plasmolyse chez le végétal ; hémolyse, crénelure chez l'animal) sont la traduction cellulaire directe des mouvements d'eau entre la cellule et son milieu ;
\item l'analyse de courbes cinétiques (proportionnalité $\Rightarrow$ diffusion simple ; plateau de saturation $\Rightarrow$ protéines transporteurs ; inhibition par cyanure $\Rightarrow$ dépendance à l'ATP $\Rightarrow$ transport actif) permet de distinguer rigoureusement \cle{diffusion simple}, \cle{diffusion facilitée} et \cle{transport actif} ;
\item la \cle{dialyse} (diffusion des solutés) complète l'osmose pour comprendre les échanges à travers les membranes semi-perméables (rein artificiel, expérience de Dutrochet) ;
\item les échanges de particules (\cle{endocytose}, \cle{exocytose}) permettent le franchissement membranaire de grosses molécules ou particules, au coût d'ATP ;
\item les savoirs scientifiques de la leçon ont une portée pratique immédiate : raisonner les apports d'engrais en agriculture camerounaise, comprendre les techniques traditionnelles de conservation (salage, séchage, fumage), et plus largement toute situation où une cellule est placée dans un milieu de concentration différente (perfusions médicales isotoniques, dessalement, cuisine).
\end{itemize}

\begin{attention}[Pièges à éviter dans ce type d'activité]
\begin{itemize}
\item \textbf{Vocabulaire osmose} : Ne jamais dire « le sel attire l'eau » ou « l'eau va vers le sel ». L'eau diffuse \emph{spontanément} du milieu où son potentiel chimique est élevé (hypotonique, eau « libre ») vers le milieu où il est bas (hypertonique, eau « liée » aux solutés). C'est une diffusion passive du solvant.
\item \textbf{Plasmolyse vs Crénelure} : Ne pas confondre. \cle{Plasmolyse} = cellule végétale (paroi rigide, décollement membrane/paroi). \cle{Crénelure} = cellule animale (pas de paroi, rétrécissement global, aspect ébouriffé).
\item \textbf{Interprétation des courbes de transport} : Un plateau (saturation) exclut la diffusion simple. Il est compatible avec la diffusion facilitée \emph{et} le transport actif. \textbf{Seul l'effet d'un inhibiteur métabolique (cyanure, DNP, ou absence d'O$_2$) ou la mise en évidence d'un transport contre gradient permet de conclure au transport actif.}
\item \textbf{Conseil agricole} : La solution n'est pas d'arrêter l'arrosage (ce qui aggraverait la concentration saline) mais au contraire d'arroser \textbf{abondamment} pour diluer et lessiver les sels. Ne pas confondre « excès d'eau » (asphyxie racinaire) et « excès de sels » (sécheresse physiologique).
\item \textbf{Unités et puissances de 10} : Vigilance sur les conversions : \SI{500}{g.m^{-2}} = \SI{5000}{kg.ha^{-1}} (dose massive). Vitesse en µmol.h$^{-1}$.g$^{-1}$ (masse fraîche ou sèche ? Préciser si possible, ici masse fraîche implicite).
\end{itemize}
\end{attention}

\newpage

\coursec{Méthodes et exercices résolus}

\courssub{Identifier au microscope les états des cellules animale et végétale en solutions de concentrations différentes}

\begin{methode}[Reconnaître turgescence, plasmolyse, hémolyse et crénelure]
\begin{enumerate}
\item \cle{Identifier le type cellulaire} : présence d'une paroi pecto-cellulosique et d'une grande vacuole centrale $\Rightarrow$ cellule \textbf{végétale} ; absence de paroi, membrane plasmique seule $\Rightarrow$ cellule \textbf{animale}.
\item \cle{Comparer la concentration} du milieu extérieur à celle du milieu intracellulaire : milieu \textbf{hypotonique} (moins concentré en solutés), \textbf{isotonique} (même concentration) ou \textbf{hypertonique} (plus concentré en solutés).
\item \cle{Prévoir le sens du flux d'eau} : l'eau se déplace toujours du milieu le \emph{moins} concentré en solutés (le plus riche en eau, hypotonique) vers le milieu le \emph{plus} concentré en solutés (le plus pauvre en eau, hypertonique).
\item \cle{Conclure sur l'état observé} :
\begin{itemize}
\item Cellule végétale en milieu hypotonique $\rightarrow$ \textbf{turgescence} (vacuole gonflée, cytoplasme plaqué contre la paroi) ;
\item Cellule végétale en milieu hypertonique $\rightarrow$ \textbf{plasmolyse} (membrane plasmique décollée de la paroi, vacuole rétractée) ;
\item Cellule animale (hématie) en milieu hypotonique $\rightarrow$ \textbf{hémolyse} (gonflement puis éclatement) ;
\item Cellule animale en milieu hypertonique $\rightarrow$ \textbf{crénelure} (rétraction, aspect hérissé).
\end{itemize}
\item \cle{Rédiger la conclusion} en citant le phénomène responsable : l'\textbf{osmose}, diffusion de l'eau (solvant) à travers une membrane semi-perméable.
\end{enumerate}
\end{methode}

\begin{exoresolu}[États d'épiderme d'oignon et d'hématies]
Des lamelles d'épiderme d'oignon sont placées dans trois milieux : eau distillée (A), solution de NaCl à \SI{9}{g/L} (B), solution de NaCl à \SI{100}{g/L} (C). Par ailleurs, des hématies sont placées dans de l'eau distillée (D) et dans du sérum à \SI{20}{g/L} de NaCl (E). Décrire et interpréter l'état des cellules dans chaque cas.
\tcblower
\textbf{Raisonnement préalable.} La vacuole des cellules d'oignon et le cytoplasme des hématies contiennent des substances dissoutes : leur milieu interne est plus concentré que l'eau distillée, voisin d'une solution physiologique ($\approx$ \SI{9}{g/L} de NaCl pour les hématies), et moins concentré qu'une solution à \SI{100}{g/L}.

\textbf{Cellules végétales (oignon).}
\begin{itemize}
\item \textbf{Milieu A (eau distillée, hypotonique)} : l'eau entre dans la vacuole par osmose ; la vacuole se gonfle et plaque le cytoplasme contre la paroi : c'est la \textbf{turgescence}. La paroi pecto-cellulosique rigide s'oppose à l'éclatement.
\item \textbf{Milieu B (NaCl \SI{9}{g/L}, environ isotonique)} : les flux d'eau entrants et sortants s'équilibrent ; la cellule garde son aspect normal.
\item \textbf{Milieu C (NaCl \SI{100}{g/L}, hypertonique)} : l'eau sort de la vacuole ; le volume vacuolaire diminue et la membrane plasmique se décolle de la paroi : c'est la \textbf{plasmolyse}.
\end{itemize}

\textbf{Cellules animales (hématies).}
\begin{itemize}
\item \textbf{Milieu D (eau distillée, hypotonique)} : l'eau entre massivement ; dépourvue de paroi protectrice, l'hématie gonfle puis éclate : c'est l'\textbf{hémolyse} (l'hémoglobine libérée rend le liquide rouge translucide).
\item \textbf{Milieu E (NaCl \SI{20}{g/L}, hypertonique)} : l'eau sort de l'hématie qui se rétracte et se déforme : c'est la \textbf{crénelure}.
\end{itemize}

\textbf{Conclusion.} Le sens des échanges d'eau dépend de la concentration relative des milieux séparés par la membrane : l'eau diffuse du milieu hypotonique vers le milieu hypertonique. C'est pourquoi les perfusions hospitalières utilisent du sérum physiologique à \SI{9}{g/L}, isotonique au plasma sanguin, afin de préserver l'intégrité des hématies.
\end{exoresolu}

\courssub{Interpréter les résultats de l'expérience de Dutrochet}

\begin{methode}[Analyser l'osmomètre de Dutrochet]
\begin{enumerate}
\item \cle{Décrire le dispositif} : une vessie de porc (membrane animale semi-perméable) remplie d'une solution concentrée (eau sucrée), fermée et prolongée par un tube capillaire, le tout plongé dans de l'eau pure.
\item \cle{Observer} : le niveau du liquide \textbf{monte} dans le tube capillaire.
\item \cle{Formuler l'hypothèse} : de l'eau est entrée dans la vessie.
\item \cle{Interpréter} : la membrane est perméable à l'eau mais pas (ou très peu) au saccharose ; l'eau diffuse du milieu le moins concentré en solutés (eau pure) vers le milieu le plus concentré (solution sucrée) : c'est l'\textbf{osmose}.
\item \cle{Conclure} : le volume interne augmente, d'où l'ascension du liquide ; elle s'arrête lorsque la pression hydrostatique exercée par la colonne de liquide compense la pression osmotique.
\item \cle{Vérifier le contrôle} : avec de l'eau pure des deux côtés de la membrane, aucune variation de niveau n'est observée (dispositif témoin).
\end{enumerate}
\end{methode}

\begin{exoresolu}[L'osmomètre de Dutrochet]
Une vessie de porc remplie d'une solution de saccharose à \SI{100}{g/L} est reliée à un tube capillaire et plongée dans un bécher d'eau distillée. Après 2 heures, le niveau du liquide dans le capillaire s'est élevé de \SI{15}{cm}. Un dispositif témoin contenant de l'eau distillée dans la vessie ne montre aucune variation. Interpréter ces résultats.
\tcblower
\textbf{Observations.} (1) Avec la solution sucrée, le niveau monte de \SI{15}{cm} : le volume de liquide contenu dans la vessie a augmenté. (2) Avec l'eau distillée des deux côtés, aucune variation n'est enregistrée.

\textbf{Hypothèse.} De l'eau a traversé la membrane de la vessie, de l'extérieur vers l'intérieur.

\textbf{Interprétation.} La membrane de la vessie est \textbf{semi-perméable} : elle laisse passer les molécules d'eau mais retient les molécules de saccharose. La concentration en eau libre est plus grande dans le bécher (eau pure) que dans la solution sucrée ; l'eau diffuse donc du bécher vers la vessie, c'est-à-dire du milieu \textbf{hypotonique} vers le milieu \textbf{hypertonique}. Ce mouvement du solvant à travers une membrane semi-perméable est l'\textbf{osmose}. L'augmentation du volume interne provoque l'ascension du liquide dans le capillaire. Le témoin confirme que, sans différence de concentration, il n'y a pas de flux net d'eau.

\textbf{Conclusion.} L'expérience de Dutrochet met en évidence l'osmose : diffusion de l'eau du milieu le moins concentré en solutés vers le milieu le plus concentré, à travers une membrane semi-perméable.
\end{exoresolu}

\courssub{Interpréter les résultats de l'expérience de Pfeffer}

\begin{methode}[Exploiter l'osmomètre de Pfeffer et la pression osmotique]
\begin{enumerate}
\item \cle{Décrire le dispositif} : vase en terre cuite poreuse dont les parois sont imprégnées de ferrocyanure de cuivre (membrane \emph{artificielle} rigide et semi-perméable), rempli de solution sucrée, relié à un manomètre à mercure, et plongé dans l'eau pure.
\item \cle{Observer} : l'eau entre dans le vase, mais la paroi rigide empêche toute variation de volume ; c'est la \textbf{pression} interne qui augmente, comme l'indique le manomètre.
\item \cle{Définir la pression osmotique} $\pi$ : pression qu'il faut exercer sur la solution pour arrêter l'entrée d'eau à travers la membrane semi-perméable.
\item \cle{Exploiter les mesures} : vérifier que $\pi$ est proportionnelle à la concentration molaire $C$ du soluté et à la température absolue $T$ (loi de Van't Hoff) :
\[ \pi = R\,T\,C \quad \text{avec } R = \SI{8,31}{J.K^{-1}.mol^{-1}},\ T \text{ en K},\ C \text{ en } \si{mol.m^{-3}},\ \pi \text{ en Pa}. \]
\item \cle{Convertir} systématiquement avant tout calcul : $T(\mathrm{K}) = \theta(^{\circ}\mathrm{C}) + 273$ ; $C$ en \si{mol.m^{-3}} ($\SI{1}{mol/L} = \SI{1000}{mol.m^{-3}}$).
\end{enumerate}
\end{methode}

\begin{exoresolu}[Mesure de la pression osmotique]
Avec l'osmomètre de Pfeffer maintenu à \SI{27}{\degreeCelsius}, on obtient les résultats suivants pour des solutions de saccharose :

\begin{center}
\begin{tabular}{|c|c|c|c|}
\hline
\rowcolor{popblueL} Concentration $C$ (\si{mol/L}) & 0,1 & 0,2 & 0,4 \\
\hline Pression osmotique $\pi$ (\si{atm}) & 2,46 & 4,93 & 9,86 \\
\hline
\end{tabular}
\end{center}

Montrer que $\pi$ est proportionnelle à $C$, puis calculer la pression osmotique d'une solution de saccharose à \SI{0,5}{mol/L} à \SI{27}{\degreeCelsius}, en pascals (on donne $R = \SI{8,31}{J.K^{-1}.mol^{-1}}$).
\tcblower
\textbf{1. Proportionnalité.} Calculons le rapport $\pi/C$ pour chaque solution :
\[ \frac{2,46}{0,1} = 24,6 \ ;\quad \frac{4,93}{0,2} = 24,65 \ ;\quad \frac{9,86}{0,4} = 24,65 \quad (\si{atm.L/mol}). \]
Le rapport $\pi/C$ est constant ($\approx 24,65$ \si{atm.L/mol}) : la pression osmotique est donc \textbf{proportionnelle à la concentration molaire} du soluté, conformément à la loi de Van't Hoff $\pi = RTC$.

\textbf{2. Calcul de $\pi$ pour $C = \SI{0,5}{mol/L}$.}
\begin{align*}
T &= 27 + 273 = \SI{300}{K} \\
C &= \SI{0,5}{mol/L} = \SI{500}{mol.m^{-3}} \\
\pi &= R\,T\,C = 8,31 \times 300 \times 500 \\
\pi &= \SI{1246500}{Pa} \approx 1,25 \times 10^{6}\ \si{Pa} \approx \SI{12,5}{bar}.
\end{align*}

\textbf{Conclusion.} La pression osmotique de la solution de saccharose à \SI{0,5}{mol/L} à \SI{27}{\degreeCelsius} vaut environ $1,25 \times 10^{6}$ \si{Pa}, soit plus de 12 fois la pression atmosphérique : les forces osmotiques mises en jeu au niveau des cellules sont considérables, ce qui explique l'éclatement des hématies en eau pure et le gonflement des tissus végétaux en milieu hypotonique.
\end{exoresolu}

\courssub{Réaliser et interpréter le graphe dégageant la notion de diffusion facilitée}

\begin{methode}[Distinguer diffusion simple et diffusion facilitée sur un graphe]
\begin{enumerate}
\item \cle{Tracer} la vitesse de transport $v$ (ordonnée, en \si{mol.s^{-1}} ou en unité arbitraire) en fonction de la concentration extracellulaire du soluté $[S]$ (abscisse).
\item \cle{Comparer les allures} :
\begin{itemize}
\item \textbf{Diffusion simple} : droite passant par l'origine ; $v$ est proportionnelle à $[S]$ (aucune limite, aucune protéine impliquée) ;
\item \textbf{Diffusion facilitée} : courbe qui croît puis atteint un \textbf{plateau} ($v_{max}$) : les transporteurs membranaires (protéines canaux ou perméases) sont en nombre limité, donc \textbf{saturables}.
\end{itemize}
\item \cle{Interpréter le plateau} : lorsque tous les transporteurs sont occupés, augmenter $[S]$ n'augmente plus la vitesse $\Rightarrow$ notion de \textbf{saturation}.
\item \cle{Conclure} : la diffusion facilitée est un transport \textbf{passif} (dans le sens du gradient de concentration, sans consommation d'ATP) mais \textbf{médié par des protéines membranaires}, donc spécifique et saturable.
\end{enumerate}
\end{methode}

\begin{exoresolu}[Entrée du glucose dans l'hématie]
On mesure la vitesse d'entrée du glucose dans des hématies en fonction de la concentration extracellulaire en glucose, en l'absence (courbe 1) puis en présence (courbe 2) de cytochalasine B, inhibiteur des transporteurs de glucose :

\begin{center}
\begin{popfigure}
\begin{tikzpicture}
\begin{axis}[width=0.85\linewidth,height=6.5cm,
xlabel={$[S]$ glucose extracellulaire (\si{mmol/L})},
ylabel={Vitesse d'entrée $v$ (u.a.)},
xmin=0,xmax=10,ymin=0,ymax=12,
legend style={at={(0.97,0.35)},anchor=east,font=\small}]
\addplot[popcurve,domain=0:10,samples=100,color=popblue]{12*x/(2+x)};
\addlegendentry{Courbe 1 : sans inhibiteur}
\addplot[popcurve,domain=0:10,samples=60,color=poporange]{0.4*x};
\addlegendentry{Courbe 2 : avec cytochalasine B}
\end{axis}
\end{tikzpicture}
\legende{Vitesse d'entrée du glucose dans l'hématie en fonction de sa concentration extracellulaire, en l'absence (courbe 1) ou en présence (courbe 2) de cytochalasine B.}
\end{popfigure}
\end{center}

Attribuer chaque courbe à un type de transport et dégager la notion de diffusion facilitée.
\tcblower
\textbf{Analyse de la courbe 1.} La vitesse croît avec $[S]$ puis plafonne vers $v_{max} \approx 12$ u.a. : il existe un \textbf{plateau de saturation}. Cette allure hyperbolique traduit l'intervention d'un nombre limité de \textbf{protéines transporteurs} (les perméases GLUT de la membrane des hématies) : lorsque tous les transporteurs sont occupés, la vitesse ne peut plus augmenter. C'est la \textbf{diffusion facilitée}.

\textbf{Analyse de la courbe 2.} En présence de cytochalasine B, qui bloque les transporteurs, la vitesse devient faible et \textbf{proportionnelle} à $[S]$ (droite passant par l'origine) : il ne subsiste que la \textbf{diffusion simple} à travers la bicouche phospholipidique, très lente car le glucose est une molécule hydrophile.

\textbf{Interprétation.} L'écart entre les deux courbes représente le transport assuré par les protéines membranaires. La diffusion facilitée : (1) s'effectue dans le sens du gradient de concentration, donc \textbf{sans dépense d'ATP} (transport passif) ; (2) exige une \textbf{protéine membranaire spécifique} ; (3) est \textbf{saturable} (existence d'une vitesse maximale $v_{max}$).

\textbf{Conclusion.} L'existence d'un plateau de saturation et l'effet d'un inhibiteur spécifique du transporteur sont les deux critères expérimentaux qui établissent la diffusion facilitée, mode d'entrée du glucose dans la plupart des cellules de l'organisme.
\end{exoresolu}

\courssub{Concevoir un protocole expérimental mettant en évidence l'osmose}

\begin{methode}[Construire un protocole expérimental rigoureux]
\begin{enumerate}
\item \cle{Formuler le problème et l'hypothèse} : « l'eau diffuse à travers une membrane semi-perméable, du milieu le moins concentré en solutés vers le milieu le plus concentré ».
\item \cle{Choisir le matériel} : membrane semi-perméable (vessie, cellophane, tube à dialyse, ou tubercule de pomme de terre évidé), solution concentrée (eau sucrée ou salée), eau distillée, tube capillaire ou balance de précision.
\item \cle{Définir la variable mesurée} : variation de niveau de liquide, de masse ou de longueur.
\item \cle{Prévoir un \textbf{dispositif témoin}} : même montage avec de l'eau distillée des deux côtés de la membrane, afin de prouver que l'effet observé est bien dû à la différence de concentration.
\item \cle{Décrire le protocole} en étapes numérotées, reproductibles, avec des quantités précises (volumes, masses, durées, température).
\item \cle{Prévoir les résultats attendus} et le critère de validation ou de rejet de l'hypothèse.
\end{enumerate}
\end{methode}

\begin{exoresolu}[Mise en évidence de l'osmose avec des cylindres de pomme de terre]
Un élève veut démontrer que l'eau entre dans les cellules d'un végétal placé en milieu hypotonique. Il dispose de pommes de terre, d'eau distillée, d'une solution de NaCl à \SI{100}{g/L}, d'une balance au centigramme et de couteaux. Concevoir le protocole, prévoir les résultats et conclure.
\tcblower
\textbf{Hypothèse.} Les cellules de la pomme de terre, dont le milieu interne (jus vacuolaire) est concentré, absorbent l'eau en milieu hypotonique par osmose et la perdent en milieu hypertonique.

\textbf{Protocole.}
\begin{enumerate}
\item Découper 6 cylindres de pomme de terre de même diamètre et de même longueur (\SI{5}{cm}) à l'aide d'un emporte-pièce.
\item Peser chaque cylindre (masse initiale $m_0$, par exemple \SI{5,00}{g}) et les répartir en deux lots de 3.
\item Plonger le lot 1 dans l'eau distillée (milieu hypotonique) et le lot 2 dans la solution de NaCl à \SI{100}{g/L} (milieu hypertonique), pendant \SI{60}{min} à température constante.
\item Essuyer délicatement chaque cylindre, le peser (masse finale $m_f$) et calculer la variation relative de masse : $\dfrac{m_f - m_0}{m_0} \times 100$ (en \%).
\end{enumerate}

\textbf{Résultats attendus.} Lot 1 (eau distillée) : $m_f > m_0$, gain de masse de l'ordre de $+10$ à $+20\,\%$, cylindres fermes et rigides (turgescence des cellules). Lot 2 (NaCl concentré) : $m_f < m_0$, perte de masse, cylindres mous et flétris (plasmolyse des cellules).

\textbf{Interprétation.} En milieu hypotonique, l'eau entre dans les cellules par osmose (la membrane plasmique est semi-perméable) : la masse des cylindres augmente. En milieu hypertonique, l'eau sort des cellules : la masse diminue. Les trois répétitions par lot permettent de calculer une moyenne et de limiter les erreurs de mesure.

\textbf{Conclusion.} L'hypothèse est validée : le tubercule de pomme de terre se comporte comme un osmomètre ; le sens du flux d'eau dépend de la concentration du milieu extérieur, ce qui met en évidence le phénomène d'osmose.
\end{exoresolu}

\courssub{Distinguer transport passif, transport actif et transport de particules}

\begin{methode}[Identifier un mode de transport à partir de données expérimentales]
\begin{enumerate}
\item \cle{Comparer les concentrations} de part et d'autre de la membrane : le transport se fait-il dans le sens du gradient (du milieu le plus concentré vers le moins concentré) ou \textbf{contre} le gradient ?
\item \cle{Tester la dépendance énergétique} : ajouter un inhibiteur de la respiration cellulaire (cyanure, azoture) ou supprimer l'apport d'ATP. Si le transport s'arrête $\Rightarrow$ \textbf{transport actif} ; s'il se poursuit $\Rightarrow$ transport passif.
\item \cle{Tester la saturation} : si la vitesse plafonne à forte concentration $\Rightarrow$ intervention de protéines transporteurs (diffusion facilitée ou transport actif) ; si elle reste proportionnelle à la concentration $\Rightarrow$ diffusion simple.
\item \cle{Observer la taille des éléments transportés} : macromolécules ou particules (bactéries, gouttelettes liquides) $\Rightarrow$ \textbf{endocytose} (phagocytose si la particule est solide, pinocytose si elle est liquide) ou \textbf{exocytose} ; ces transports s'accompagnent toujours de déformations membranaires et d'une consommation d'ATP.
\item \cle{Rédiger la réponse} en justifiant chaque critère par la donnée expérimentale correspondante.
\end{enumerate}
\end{methode}

\begin{exoresolu}[Absorption du potassium par des racines]
Des racines de maïs sont placées dans une solution contenant des ions $K^{+}$ à \SI{0,5}{mmol/L}. On mesure la concentration intracellulaire en $K^{+}$ : \SI{80}{mmol/L}. L'expérience est répétée en présence de cyanure (bloqueur de la respiration cellulaire) : l'absorption de $K^{+}$ chute de 95\,\%. Enfin, on observe que la vitesse d'absorption plafonne lorsque la concentration externe augmente. Identifier le mode de transport du $K^{+}$ et justifier.
\tcblower
\textbf{Critère 1 : sens du transport.} La concentration interne (\SI{80}{mmol/L}) est 160 fois supérieure à la concentration externe (\SI{0,5}{mmol/L}) : le $K^{+}$ est absorbé \textbf{contre son gradient de concentration}. Aucun transport passif ne peut réaliser une telle accumulation.

\textbf{Critère 2 : dépendance énergétique.} Le cyanure bloque la respiration cellulaire, donc la production d'ATP ; l'absorption chute alors de 95\,\% : le transport \textbf{consomme de l'énergie}. C'est la signature du \textbf{transport actif}.

\textbf{Critère 3 : saturation.} Le plafonnement de la vitesse à forte concentration externe montre l'intervention d'un nombre limité de \textbf{protéines transporteurs} (pompes à $K^{+}$ membranaires), saturables.

\textbf{Conclusion.} L'absorption du potassium par les racines est un \textbf{transport actif} : déplacement contre le gradient de concentration, médié par des pompes membranaires saturables et alimenté par l'ATP issu de la respiration cellulaire. Ce mécanisme explique l'accumulation des sels minéraux dans les racines, indispensable à la nutrition minérale des plantes cultivées comme le maïs.
\end{exoresolu}

\courssub{Exploiter des données sur l'endocytose et l'exocytose}

\begin{methode}[Analyser un document sur les transports de particules]
\begin{enumerate}
\item \cle{Repérer la nature de la particule} : solide (bactérie, débris cellulaire) $\rightarrow$ \textbf{phagocytose} ; liquide ou macromolécules dissoutes $\rightarrow$ \textbf{pinocytose} ; expulsion de substances vers le milieu extracellulaire $\rightarrow$ \textbf{exocytose}.
\item \cle{Décrire les étapes} sur le document : pour l'endocytose, invagination de la membrane plasmique $\rightarrow$ formation d'une vésicule intracellulaire (phagosome ou vésicule de pinocytose) $\rightarrow$ fusion avec un lysosome (digestion enzymatique) ; pour l'exocytose, migration d'une vésicule de sécrétion $\rightarrow$ fusion avec la membrane plasmique $\rightarrow$ libération du contenu à l'extérieur.
\item \cle{Citer un exemple biologique} : leucocyte phagocytant une bactérie ; cellule du pancréas libérant l'insuline par exocytose.
\item \cle{Rappeler le coût énergétique} : ces transports exigent de l'ATP et mettent en jeu des déformations de la membrane plasmique.
\end{enumerate}
\end{methode}

\begin{exoresolu}[Le macrophage et la bactérie]
Sur une électronographie, on observe un macrophage (leucocyte) émettant des prolongements cytoplasmiques qui entourent une bactérie ; à l'intérieur de la cellule, une vésicule contenant la bactérie fusionne avec un organite riche en enzymes. Décrire et nommer les phénomènes observés, puis préciser leur importance.
\tcblower
\textbf{Description.} (1) Les prolongements cytoplasmiques (pseudopodes) du macrophage entourent la bactérie : c'est l'\textbf{invagination} de la membrane plasmique. (2) La bactérie se retrouve enfermée dans une vésicule intracellulaire, le \textbf{phagosome}. (3) Le phagosome fusionne avec un \textbf{lysosome} (organite riche en enzymes digestives), formant une vacuole digestive au sein de laquelle la bactérie est détruite.

\textbf{Identification.} L'ensemble constitue la \textbf{phagocytose}, forme d'\textbf{endocytose} concernant les particules solides. Elle nécessite des déformations membranaires et une dépense d'ATP.

\textbf{Importance.} La phagocytose est un mécanisme de \textbf{défense de l'organisme} : les macrophages et les polynucléaires éliminent les agents pathogènes (bactéries, hématies parasitées par \emph{Plasmodium} au cours du paludisme) ainsi que les débris cellulaires. Symétriquement, l'\textbf{exocytose} permet la libération de substances volumineuses (hormones protéiques comme l'insuline, enzymes digestives, neurotransmetteurs) sans traversée directe de la membrane.

\textbf{Conclusion.} Endocytose et exocytose assurent le transport des particules et des macromolécules incapables de traverser la membrane plasmique par les voies classiques, au prix d'une consommation d'énergie.
\end{exoresolu}

\begin{attention}[Les 5 erreurs qui coûtent des points au Bac]
\begin{enumerate}
\item \textbf{Inverser le sens du flux d'eau.} L'eau se déplace du milieu le \emph{moins} concentré en solutés (hypotonique) vers le milieu le \emph{plus} concentré (hypertonique). Astuce : l'eau va « diluer » le milieu le plus concentré. Ne jamais écrire que « le sel diffuse vers l'eau » dans l'expérience de Dutrochet : la membrane semi-perméable retient le soluté.
\item \textbf{Confondre turgescence et plasmolyse, hémolyse et crénelure.} Milieu hypotonique : turgescence (végétal) ou hémolyse (animal) ; milieu hypertonique : plasmolyse (végétal) ou crénelure (animal). Préciser que seule la cellule végétale résiste à l'éclatement grâce à sa \textbf{paroi pecto-cellulosique}.
\item \textbf{Oublier les conversions dans la loi de Van't Hoff.} Dans $\pi = RTC$, la température doit être exprimée en \textbf{kelvins} ($T = \theta + 273$) et la concentration en \si{mol.m^{-3}} ($\SI{1}{mol/L} = \SI{1000}{mol.m^{-3}}$). Une température laissée en degrés Celsius fausse tout le calcul.
\item \textbf{Affirmer que la diffusion facilitée consomme de l'ATP.} La diffusion facilitée est un transport \textbf{passif} (dans le sens du gradient, sans dépense d'énergie) ; ses critères sont la \textbf{saturation} (plateau de $v_{max}$) et la \textbf{spécificité} (protéine transporteur, effet d'inhibiteurs). Seul le transport actif consomme de l'ATP et s'effectue contre le gradient.
\item \textbf{Décrire sans interpréter ni conclure.} Dans l'exploitation d'une expérience (Dutrochet, Pfeffer, graphe de transport), le correcteur attend la chaîne complète : observation $\rightarrow$ hypothèse $\rightarrow$ interprétation (mécanisme : osmose, diffusion, transport actif) $\rightarrow$ conclusion qui répond au problème posé. Citer toujours le \textbf{dispositif témoin} et préciser son rôle.
\end{enumerate}
\end{attention}

\newpage

\coursec{Exercices}

\courssub{Je vérifie mes connaissances}

\begin{exercice}[QCM : les échanges cellulaires]
Pour chaque question, une seule réponse est exacte. Recopie la lettre correspondante.

\textbf{1.} Une cellule végétale placée dans une solution hypotonique :
\begin{itemize}
\item[a)] se plasmolyse ;
\item[b)] devient turgescente ;
\item[c)] éclate ;
\item[d)] subit une crénelure.
\end{itemize}

\textbf{2.} L'osmose est le mouvement :
\begin{itemize}
\item[a)] des molécules de soluté à travers une membrane perméable ;
\item[b)] de l'eau à travers une membrane semi-perméable, du milieu le moins concentré vers le milieu le plus concentré ;
\item[c)] de l'eau à travers une membrane perméable, du milieu concentré vers le milieu dilué ;
\item[d)] des particules à travers la membrane plasmique.
\end{itemize}

\textbf{3.} La diffusion facilitée se distingue de la diffusion simple par :
\begin{itemize}
\item[a)] la consommation d'ATP ;
\item[b)] l'intervention de protéines membranaires spécifiques ;
\item[c)] le transport contre le gradient de concentration ;
\item[d)] la formation de vésicules.
\end{itemize}

\textbf{4.} La phagocytose est :
\begin{itemize}
\item[a)] l'entrée de gouttelettes liquides dans la cellule ;
\item[b)] la sortie de particules par fusion d'une vésicule avec la membrane ;
\item[c)] l'ingestion de particules solides par émission de pseudopodes ;
\item[d)] le passage d'ions par un canal.
\end{itemize}
\end{exercice}

\begin{exercice}[Vrai ou faux, à justifier]
Réponds par vrai ou faux, puis justifie ta réponse en une ou deux phrases.
\begin{enumerate}
\item Un globule rouge placé dans de l'eau distillée subit une crénelure.
\item La dialyse concerne le passage des molécules de soluté à travers une membrane perméable.
\item Le transport actif s'effectue toujours dans le sens du gradient de concentration.
\item L'exocytose et l'endocytose sont des mécanismes qui consomment de l'énergie.
\item Dans l'expérience de Dutrochet, la colonne d'eau monte dans le tube parce que l'eau pénètre dans la vessie remplie de solution concentrée.
\end{enumerate}
\end{exercice}

\begin{exercice}[Définitions express]
Définis en deux lignes maximum chacun des termes suivants : \cle{turgescence}, \cle{plasmolyse}, \cle{hémolyse}, \cle{solution isotonique}, \cle{pinocytose}, \cle{pression osmotique}.
\end{exercice}

\begin{exercice}[Calculs directs : tonicité d'une solution]
On dissout \num{4,5}~g de chlorure de sodium (NaCl) dans \num{500}~mL d'eau distillée.
\begin{enumerate}
\item Calcule la concentration massique de la solution obtenue, en g/L puis en pour mille (\textperthousand).
\item Le sérum physiologique a une concentration de 9~\textperthousand. La solution préparée est-elle hypotonique, isotonique ou hypertonique par rapport au sang ? Justifie.
\item Que subira un globule rouge placé dans cette solution ?
\end{enumerate}
\end{exercice}

\courssub{Je m'entraîne}

\begin{exercice}[Épiderme d'oignon en solutions salines]
On dépose trois lamelles d'épiderme d'oignon dans trois solutions : S$_1$ (eau distillée), S$_2$ (NaCl à 9~\textperthousand), S$_3$ (NaCl à 100~\textperthousand). Après 10 minutes, on observe au microscope.
\begin{enumerate}
\item Décris l'état des cellules végétales attendu dans chaque solution, en précisant la position du protoplaste par rapport à la paroi.
\item Nomme les phénomènes observés dans S$_1$ et S$_3$.
\item Explique, en termes de mouvements d'eau, ce qui se passe dans S$_3$.
\item Pourquoi les cellules végétales n'éclatent-elles pas dans S$_1$, contrairement aux globules rouges ?
\end{enumerate}
\end{exercice}

\begin{exercice}[L'expérience de Dutrochet]
Une vessie de porc remplie d'une solution de saccharose à 100~g/L est fermée par un bouchon traversé par un tube vertical, puis plongée dans un bécher d'eau distillée. Après quelques heures, le niveau liquide monte dans le tube de 25~cm.
\begin{enumerate}
\item Schématise et légende le montage.
\item Interprète la montée du liquide dans le tube.
\item Prédis le résultat si l'on inverse les liquides (eau distillée dans la vessie, solution de saccharose dans le bécher).
\item Précise le rôle de la membrane de la vessie et le nom de la propriété qu'elle présente.
\end{enumerate}
\end{exercice}

\begin{exercice}[Diffusion simple ou diffusion facilitée ?]
On mesure la vitesse d'entrée de deux solutés dans des cellules en fonction de leur concentration externe :

\begin{center}
\begin{tabularx}{\linewidth}{|l|*{5}{>{\centering\arraybackslash}X|}}
\hline
\rowcolor{popblueL} Concentration externe (mmol/L) & 0,5 & 1 & 2 & 4 & 8 \\
\hline
Vitesse d'entrée du soluté A (unités arbitraires) & 2 & 4 & 8 & 16 & 32 \\
\hline
Vitesse d'entrée du soluté B (unités arbitraires) & 5 & 9 & 14 & 17 & 18 \\
\hline
\end{tabularx}
\end{center}

\begin{enumerate}
\item Trace sur un même graphique les courbes vitesse d'entrée = f(concentration externe) pour les deux solutés.
\item Compare l'allure des deux courbes.
\item Attribue à chaque soluté son mode de transport (diffusion simple ou diffusion facilitée) en justifiant.
\item Explique l'existence du plateau observé pour le soluté B.
\end{enumerate}
\end{exercice}

\begin{exercice}[Un macrophage au travail]
Un macrophage rencontre une bactérie dans un tissu infecté. On observe au microscope électronique l'émission de prolongements cytoplasmiques qui entourent la bactérie, puis la formation d'une vésicule intracellulaire contenant la bactérie.
\begin{enumerate}
\item Nomme les prolongements émis par le macrophage et le mécanisme mis en jeu.
\item Schématise les étapes de ce mécanisme.
\item Cite une autre forme d'endocytose et précise ce qui la distingue de celle observée ici.
\item Quel organite cellulaire interviendra ensuite pour digérer la bactérie ?
\end{enumerate}
\end{exercice}

\courssub{Je me prépare au Bac}

\begin{exercice}[Type Bac : trois solutions, trois états des globules rouges]
Au laboratoire du lycée, on place des gouttes de sang dans trois solutions A, B et C de concentrations inconnues. L'observation microscopique donne :
\begin{itemize}
\item dans la solution A : des globules rouges biconcaves, de diamètre normal (\SI{7,5}{\micro m}) ;
\item dans la solution B : des globules rouges éclatés, seuls des « fantômes » membranaires sont visibles ;
\item dans la solution C : des globules rouges ratatinés, à contours irréguliers, de diamètre réduit (\SI{5}{\micro m}).
\end{itemize}

\begin{enumerate}
\item Identifie l'état des globules rouges dans chaque solution et nomme précisément le phénomène observé dans B et dans C.
\item Classe les trois solutions par tonicité croissante en justifiant ton classement par le sens des mouvements d'eau dans chaque cas.
\item Laquelle de ces solutions pourrait être du sérum physiologique à 9~\textperthousand ? Justifie.
\item À l'hôpital de district de Fundong, un infirmier s'apprête à perfuser un patient avec de l'eau distillée par erreur. Explique les conséquences de cette erreur sur les globules rouges du patient et le danger encouru.
\item Schématise un globule rouge dans chacune des trois solutions, en représentant par des flèches le sens du mouvement d'eau.
\end{enumerate}
\end{exercice}

\begin{exercice}[Type Bac : l'expérience de Pfeiffer et la pression osmotique]
On réalise un osmomètre de Pfeiffer : une cellule à paroi de ferrocyanure de cuivre (membrane semi-perméable) remplie d'une solution de saccharose est reliée à un manomètre et plongée dans l'eau pure à \SI{20}{\celsius}. On mesure la pression osmotique $\pi$ pour différentes concentrations $C$ en saccharose :

\begin{center}
\begin{tabularx}{\linewidth}{|l|*{5}{>{\centering\arraybackslash}X|}}
\hline
\rowcolor{popblueL} $C$ (mol/L) & 0,05 & 0,10 & 0,15 & 0,20 & 0,25 \\
\hline
$\pi$ (atm) & 1,20 & 2,40 & 3,61 & 4,81 & 6,01 \\
\hline
\end{tabularx}
\end{center}

\begin{enumerate}
\item Rappelle le principe de l'expérience de Pfeiffer et précise le rôle de la membrane de ferrocyanure de cuivre.
\item Trace le graphe $\pi = f(C)$. Quelle relation mathématique simple lie $\pi$ et $C$ ?
\item Calcule le coefficient de proportionnalité $\pi/C$ et montre qu'il est voisin de $RT$, avec $R = \SI{0,082}{L.atm.mol^{-1}.K^{-1}}$ et $T = 293$~K. Détaille ton calcul.
\item En déduire la pression osmotique d'une solution de saccharose à 0,30~mol/L à la même température.
\item Deux solutions de saccharose à 0,10~mol/L et 0,25~mol/L sont séparées par la même membrane semi-perméable. Prédis le sens du mouvement d'eau et justifie en utilisant les résultats du tableau.
\end{enumerate}
\end{exercice}

\begin{exercice}[Type Bac : modes de transport et mise en évidence expérimentale de l'osmose]
\textbf{Partie A.} On étudie l'absorption de deux solutés X et Y par des cellules racinaires d'une plante cultivée dans la région de l'Ouest camerounais. Le graphe vitesse d'absorption = f(concentration externe) donne : pour X, une droite passant par l'origine ; pour Y, une courbe qui croît puis plafonne à partir de 4~mmol/L.

\begin{enumerate}
\item Interprète l'allure de chaque courbe et déduis-en le mode de transport de chaque soluté.
\item Pour vérifier si l'un de ces transports est actif, on ajoute au milieu du cyanure, inhibiteur de la respiration cellulaire. Quel résultat attend-on pour chaque soluté si Y est transporté activement ? Justifie.
\end{enumerate}

\textbf{Partie B.} Tu disposes au laboratoire du matériel suivant : un tube à dialyse (cellophane), une solution de glucose concentrée, de l'eau distillée, deux béchers, un capillaire en verre, un bouchon, un support, du fil, une solution de Lugol et de l'eau iodée.

\begin{enumerate}
\item[3.] Conçois un protocole expérimental complet (montage, dispositif témoin) permettant de mettre en évidence le phénomène d'osmose. Précise le matériel utilisé à chaque étape.
\item[4.] Décris les résultats attendus pour le dispositif expérimental et pour le témoin, et donne l'interprétation.
\item[5.] Propose une utilisation de l'eau iodée ou du Lugol pour mettre en évidence, dans le même montage, le phénomène de dialyse.
\end{enumerate}
\end{exercice}

\newpage

\coursec{Corrigés des exercices}

\coursec{Exercices résolus et corrigés — Limitation des conséquences liées aux échanges cellulaires}

\begin{corrige}[Exercice 1 — QCM : les échanges cellulaires]
\textbf{1. Réponse b).} Dans une solution \cle{hypotonique} (moins concentrée que le cytosol), l'eau entre dans la cellule végétale par osmose : la vacuole gonfle, le protoplaste plaque contre la paroi pectocellulosique qui résiste : la cellule devient \cle{turgescente}. Elle n'éclate pas grâce à la paroi rigide (réponse c fausse) ; la plasmolyse (a) et la crénelure (d) correspondent à un milieu hypertonique.

\textbf{2. Réponse b).} L'\cle{osmose} est le mouvement du solvant (l'eau) à travers une membrane \cle{semi-perméable}, du milieu le moins concentré en soluté (hypotonique) vers le milieu le plus concentré (hypertonique). Les réponses a et d décrivent la dialyse ou l'endocytose ; la réponse c inverse le sens du mouvement.

\textbf{3. Réponse b).} La \cle{diffusion facilitée} fait intervenir des \cle{protéines membranaires spécifiques} (transporteurs ou canaux), ce qui explique sa saturation à forte concentration. Elle reste passive : pas d'ATP (a faux) et elle suit le gradient de concentration (c faux). La formation de vésicules (d) caractérise l'endo/exocytose.

\textbf{4. Réponse c).} La \cle{phagocytose} est l'ingestion de particules \emph{solides} (bactéries, débris cellulaires) grâce à l'émission de \cle{pseudopodes} qui les enferment dans une vésicule, le phagosome. L'entrée de gouttelettes liquides (a) est la pinocytose ; la sortie par fusion de vésicule (b) est l'exocytose.
\end{corrige}

\begin{corrige}[Exercice 2 — Vrai ou faux, à justifier]
\begin{enumerate}
\item \textbf{Faux.} L'eau distillée est très hypotonique par rapport au cytosol du globule rouge : l'eau entre massivement par osmose, la cellule gonfle puis éclate : c'est l'\cle{hémolyse}. La crénelure se produit au contraire en milieu hypertonique (sortie d'eau).
\item \textbf{Vrai.} La \cle{dialyse} est le passage des molécules et ions de soluté à travers une membrane perméable (ou dialysante), du milieu le plus concentré vers le milieu le moins concentré, jusqu'à l'équilibre des concentrations.
\item \textbf{Faux.} Le \cle{transport actif} s'effectue \emph{contre} le gradient de concentration (du milieu le moins concentré vers le plus concentré), grâce à des pompes membranaires qui hydrolysent l'ATP. C'est le transport passif qui suit le gradient.
\item \textbf{Vrai.} L'\cle{endocytose} et l'\cle{exocytose} impliquent des déformations de la membrane plasmique et le déplacement de vésicules, processus qui consomment de l'énergie fournie par l'hydrolyse de l'ATP.
\item \textbf{Vrai.} La membrane de la vessie est semi-perméable : l'eau distillée du bécher pénètre par osmose dans la vessie où la solution est plus concentrée ; le volume interne augmente et refoule le liquide dans le tube vertical : la colonne monte.
\end{enumerate}
\end{corrige}

\begin{corrige}[Exercice 3 — Définitions express]
\begin{itemize}
\item \textbf{Turgescence :} état d'une cellule végétale gorgée d'eau en milieu hypotonique ; la vacuole distendue plaque le cytoplasme contre la paroi, ce qui rigidifie les tissus.
\item \textbf{Plasmolyse :} rétractation du protoplaste qui se décolle de la paroi pectocellulosique lorsqu'une cellule végétale perd de l'eau en milieu hypertonique.
\item \textbf{Hémolyse :} éclatement d'un globule rouge placé en milieu hypotonique, avec libération de l'hémoglobine ; il ne subsiste que des « fantômes » membranaires.
\item \textbf{Solution isotonique :} solution dont la concentration en solutés est égale à celle du milieu intracellulaire ; les flux d'eau entrants et sortants se compensent, la cellule garde sa forme normale (ex. sérum physiologique à 9~\textperthousand).
\item \textbf{Pinocytose :} endocytose de gouttelettes de liquide extracellulaire et de molécules dissoutes, piégées dans de petites vésicules formées par invagination de la membrane plasmique.
\item \textbf{Pression osmotique :} pression qu'il faut exercer sur une solution pour empêcher l'entrée d'eau par osmose à travers une membrane semi-perméable ; elle est proportionnelle à la concentration en soluté ($\pi = RTC$).
\end{itemize}
\end{corrige}

\begin{corrige}[Exercice 4 — Calculs directs : tonicité d'une solution]
\begin{enumerate}
\item \textbf{Concentration massique.} On applique $C_m = \dfrac{m}{V}$ avec $m = \num{4,5}$~g et $V = \num{500}$~mL $= \num{0,500}$~L :
\[ C_m = \frac{\num{4,5}}{\num{0,500}} = \num{9,0}~\text{g/L}. \]
En pour mille : $9,0$~g pour \num{1000}~g de solution, soit $C_m = 9$~\textperthousand.
\item \textbf{Tonicité.} La solution préparée (9~\textperthousand) a exactement la même concentration que le sérum physiologique (9~\textperthousand), lui-même isotonique au plasma sanguin. La solution est donc \cle{isotonique} par rapport au sang.
\item \textbf{Comportement du globule rouge.} En milieu isotonique, les entrées et sorties d'eau par osmose se compensent exactement : le globule rouge \textbf{conserve sa forme normale} de disque biconcave (diamètre $\approx \SI{7,5}{\micro m}$). Ni hémolyse, ni crénelure.
\end{enumerate}
\end{corrige}

\begin{attention}[Point de vigilance — Exercice 4]
Ne pas confondre g/L et \textperthousand : pour les solutions aqueuses diluées, 1~g/L $\approx$ 1~\textperthousand, mais il faut justifier l'équivalence (masse volumique de la solution $\approx 1$~kg/L).
\end{attention}

\begin{corrige}[Exercice 5 — Épiderme d'oignon en solutions salines]
\begin{enumerate}
\item \textbf{États attendus.}
\begin{itemize}
\item \textbf{S$_1$ (eau distillée, hypotonique) :} l'eau entre dans les cellules ; la vacuole gonfle et le protoplaste est \textbf{plaqué contre la paroi} : cellules \cle{turgescentes}.
\item \textbf{S$_2$ (NaCl à 9~\textperthousand, proche de l'isotonie) :} les flux d'eau sont équilibrés ; le protoplaste occupe normalement la cellule, sans décollement ni tension excessive : cellules en état « normal » (turgescence physiologique).
\item \textbf{S$_3$ (NaCl à 100~\textperthousand, hypertonique) :} l'eau sort des cellules ; la vacuole se vide et le protoplaste \textbf{se rétracte et se décolle de la paroi} : cellules \cle{plasmolysées}.
\end{itemize}
\item Dans S$_1$ : la \cle{turgescence}. Dans S$_3$ : la \cle{plasmolyse}.
\item \textbf{Dans S$_3$ :} le milieu externe (100~\textperthousand) est beaucoup plus concentré en solutés que le suc vacuolaire. Par \cle{osmose}, l'eau quitte la vacuole puis le cytoplasme pour gagner le milieu extérieur, à travers la membrane plasmique semi-perméable. La perte d'eau fait diminuer le volume du protoplaste, qui se détache de la paroi rigide, elle, ne se rétracte pas.
\item \textbf{Pourquoi pas d'éclatement dans S$_1$ ?} La cellule végétale est entourée d'une \textbf{paroi pectocellulosique rigide et résistante} qui s'oppose à la distension : la pression de turgescence s'établit et finit par équilibrer le flux osmotique entrant. Le globule rouge, dépourvu de paroi, ne possède qu'une membrane plasmique fragile : il gonfle jusqu'à éclater (hémolyse).
\end{enumerate}
\end{corrige}

\begin{corrige}[Exercice 6 — L'expérience de Dutrochet]
\begin{enumerate}
\item \textbf{Schéma du montage.}
\begin{popfigure}
\begin{center}
\begin{tikzpicture}[scale=0.9, font=\small]
% bécher
\draw[thick, popblue] (-3,0) -- (-3,2.6) (3,0) -- (3,2.6);
\draw[thick, popblue] (-3,0) -- (3,0);
\fill[popblue!15] (-2.95,0.05) rectangle (2.95,1.9);
\node[popblue] at (0,0.35) {eau distillée};
% vessie
\draw[thick, poporange, fill=poporange!20] (0,1.0) ellipse (1.5 and 0.85);
\node at (0,1.0) {solution de saccharose};
\node at (0,0.55) {\footnotesize 100 g/L};
% bouchon + tube
\draw[thick, fill=popgold!40] (-0.35,1.75) rectangle (0.35,2.05);
\draw[thick, popdark] (-0.12,2.05) -- (-0.12,4.6) (0.12,2.05) -- (0.12,4.6);
\fill[poporange!40] (-0.12,2.05) rectangle (0.12,4.0);
\draw[popdark, dashed] (-0.5,2.6) -- (0.5,2.6) node[right, popdark]{niveau initial};
\draw[popdark, dashed] (-0.5,4.0) -- (0.5,4.0) node[right, popdark]{niveau final (+25 cm)};
\draw[->, thick, popgreen] (-2.4,1.5) -- (-1.2,1.2) node[midway, above left, popgreen]{eau};
\draw[->, thick, popgreen] (2.4,1.5) -- (1.2,1.2);
\end{tikzpicture}
\end{center}
\legende{Expérience de Dutrochet : la vessie (membrane semi-perméable) contient la solution concentrée de saccharose ; l'eau du bécher y pénètre par osmose et fait monter la colonne liquide dans le tube.}
\end{popfigure}
\item \textbf{Interprétation.} La membrane de la vessie est semi-perméable : elle laisse passer l'eau mais pas le saccharose. Le milieu intérieur (100~g/L) étant hypertonique par rapport à l'eau distillée, l'eau pénètre dans la vessie par \cle{osmose}. Le volume interne augmente, la pression hydrostatique s'élève et refoule le liquide dans le tube : la colonne monte de 25~cm, jusqu'à ce que la pression hydrostatique équilibre la \cle{pression osmotique}.
\item \textbf{Liquides inversés.} L'eau distillée est alors dans la vessie et la solution concentrée à l'extérieur : l'eau \textbf{sort} de la vessie par osmose. Le niveau dans le tube \textbf{descend} et la vessie s'affaisse (se vide partiellement).
\item \textbf{Rôle de la membrane.} La membrane de la vessie joue le rôle de \textbf{barrière sélective} : elle possède la propriété de \cle{semi-perméabilité} (perméable à l'eau, imperméable au saccharose). C'est cette propriété qui rend le phénomène d'osmose possible.
\end{enumerate}
\end{corrige}

\begin{corrige}[Exercice 7 — Diffusion simple ou diffusion facilitée ?]
\begin{enumerate}
\item \textbf{Graphique.}
\begin{popfigure}
\begin{center}
\begin{tikzpicture}
\begin{axis}[
width=0.85\linewidth, height=6.5cm,
xlabel={Concentration externe (mmol/L)},
ylabel={Vitesse d'entrée (u. a.)},
xmin=0, xmax=9, ymin=0, ymax=35,
xtick={0,1,2,3,4,5,6,7,8},
legend pos=north west, legend style={font=\footnotesize},
grid=both, grid style={popline!40}
]
\addplot[popcurve, popblue, mark=*, samples=50, domain=0:8.5] {4*x};
\addlegendentry{Soluté A}
\addplot[popcurve, poporange, mark=square*, samples=100, domain=0:8.5] {18*x/(x+0.9)};
\addlegendentry{Soluté B}
\end{axis}
\end{tikzpicture}
\end{center}
\legende{Vitesse d'entrée de deux solutés en fonction de la concentration externe : droite passant par l'origine pour A (diffusion simple), courbe avec plateau pour B (diffusion facilitée).}
\end{popfigure}
\item \textbf{Comparaison.} Pour le soluté A, la vitesse est une \textbf{droite passant par l'origine} : elle est proportionnelle à la concentration externe (elle double quand la concentration double : 2, 4, 8, 16, 32). Pour le soluté B, la vitesse croît d'abord rapidement (5, 9, 14) puis \textbf{ralentit et plafonne} vers 17--18 u. a. : il existe une vitesse maximale $V_{\max} \approx 18$~u. a.
\item \textbf{Attribution.}
\begin{itemize}
\item \textbf{Soluté A : diffusion simple.} La proportionnalité vitesse/concentration, sans saturation, traduit un passage direct à travers la bicouche phospholipidique, uniquement dépendant du gradient de concentration.
\item \textbf{Soluté B : diffusion facilitée.} La saturation (plateau) est la signature de l'intervention de \cle{protéines transporteuses} membranaires en nombre limité.
\end{itemize}
\item \textbf{Origine du plateau.} Les protéines de transport spécifiques du soluté B sont en \textbf{nombre fini} dans la membrane. Lorsque la concentration externe devient élevée, tous les transporteurs sont occupés en permanence (\emph{saturation}) : la vitesse ne peut plus augmenter, même si la concentration externe continue de croître. C'est la notion de \cle{vitesse maximale} $V_{\max}$.
\end{enumerate}
\end{corrige}

\begin{corrige}[Exercice 8 — Un macrophage au travail]
\begin{enumerate}
\item Les prolongements cytoplasmiques sont des \cle{pseudopodes} ; le mécanisme mis en jeu est la \cle{phagocytose} (une forme d'endocytose de particules solides).
\item \textbf{Schéma des étapes.}
\begin{popfigure}
\begin{center}
\begin{tikzpicture}[scale=0.8, font=\footnotesize]
% étape 1
\draw[thick, popblue, fill=popblue!10] (0,0) circle (1);
\fill[poporange] (1.35,0) circle (0.18);
\node at (0,-1.5) {1. Contact};
\node[poporange] at (1.9,0.5) {bactérie};
% étape 2
\begin{scope}[xshift=4.2cm]
\draw[thick, popblue, fill=popblue!10] (0,0) circle (1);
\draw[thick, popblue, fill=popblue!10] (1.1,0.45) .. controls (1.7,0.6) and (1.9,0.3) .. (1.75,0) .. controls (1.9,-0.3) and (1.7,-0.6) .. (1.1,-0.45);
\fill[poporange] (1.35,0) circle (0.18);
\node at (0.6,-1.5) {2. Pseudopodes};
\end{scope}
% étape 3
\begin{scope}[xshift=9cm]
\draw[thick, popblue, fill=popblue!10] (0,0) circle (1);
\draw[thick, popgreen, fill=popgreen!20] (0.3,0) circle (0.35);
\fill[poporange] (0.3,0) circle (0.18);
\node at (0,-1.5) {3. Phagosome};
\end{scope}
\draw[->, very thick, popdark] (2.2,0) -- (3.0,0);
\draw[->, very thick, popdark] (7.0,0) -- (7.8,0);
\end{tikzpicture}
\end{center}
\legende{Phagocytose d'une bactérie par un macrophage : (1) adhésion de la bactérie à la membrane ; (2) émission de pseudopodes qui l'entourent ; (3) fermeture de la vésicule intracellulaire, le phagosome.}
\end{popfigure}
\item L'autre forme d'endocytose est la \cle{pinocytose} : elle concerne l'absorption de \textbf{gouttelettes de liquide} extracellulaire et de molécules dissoutes dans de petites vésicules, alors que la phagocytose concerne des \textbf{particules solides} et met en jeu des pseudopodes et de grandes vésicules.
\item Le \cle{lysosome} interviendra ensuite : il fusionne avec le phagosome (formation du phagolysosome) et ses \textbf{enzymes digestives} (hydrolases acides) dégradent la bactérie.
\end{enumerate}
\end{corrige}

\begin{corrige}[Exercice 9 — Type Bac : trois solutions, trois états des globules rouges]
\begin{enumerate}
\item \textbf{États observés.}
\begin{itemize}
\item \textbf{Solution A :} globules rouges biconcaves de diamètre normal (\SI{7,5}{\micro m}) : état \textbf{normal}, milieu isotonique.
\item \textbf{Solution B :} globules éclatés, « fantômes » membranaires : c'est l'\cle{hémolyse}, caractéristique d'un milieu \textbf{hypotonique}.
\item \textbf{Solution C :} globules ratatinés, contours irréguliers, diamètre réduit (\SI{5}{\micro m}) : c'est la \cle{crénelure}, caractéristique d'un milieu \textbf{hypertonique}.
\end{itemize}
\item \textbf{Classement par tonicité croissante : B $<$ A $<$ C.}
\begin{itemize}
\item Dans B, l'eau \emph{entre} dans les globules (milieu externe moins concentré que le cytosol) jusqu'à l'éclatement : B est hypotonique.
\item Dans A, les flux d'eau entrants et sortants \emph{se compensent} : A est isotonique.
\item Dans C, l'eau \emph{sort} des globules (milieu externe plus concentré) : C est hypertonique.
\end{itemize}
\item \textbf{La solution A} pourrait être le sérum physiologique à 9~\textperthousand : ce sérum est précisément \textbf{isotonique} au plasma sanguin, ce qui explique que les globules rouges y conservent leur forme et leur diamètre normaux.
\item \textbf{Conséquences de la perfusion d'eau distillée.} L'eau distillée, très hypotonique, dilue le plasma qui devient hypotonique par rapport au cytosol des hématies. L'eau pénètre massivement dans les globules rouges par osmose : ils gonflent puis \textbf{éclatent (hémolyse)}. L'hémoglobine libérée dans le plasma est toxique pour les reins (elle peut obstruer les tubules rénaux) et la destruction massive des hématies provoque une \textbf{anémie aiguë} et une hypoxie des tissus : le pronostic vital du patient est engagé. C'est pourquoi on perfuse toujours avec du sérum physiologique isotonique à 9~\textperthousand.
\item \textbf{Schémas.}
\begin{popfigure}
\begin{center}
\begin{tikzpicture}[scale=0.85, font=\footnotesize]
% A : normal
\draw[thick, popblue, fill=poppink!25] (0,0) ellipse (1 and 0.45);
\draw[popblue!60, fill=white] (0,0) ellipse (0.45 and 0.18);
\draw[<->, popgreen, thick] (-1.6,0.7) -- (-1.0,0.45);
\draw[<->, popgreen, thick] (1.6,0.7) -- (1.0,0.45);
\node at (0,-1.0) {A : isotonique};
\node at (0,-1.5) {flux équilibrés};
% B : hémolyse
\begin{scope}[xshift=5cm]
\draw[thick, popblue, dashed, fill=poppink!10] (0,0) circle (1.05);
\draw[poppink, thick] plot[smooth] coordinates {(-0.6,0.2) (-0.3,0.5) (0.1,0.3) (0.5,0.6)};
\foreach \a in {30,90,150,210,270,330} \draw[->, popgreen, thick] (\a:1.7) -- (\a:1.15);
\node at (0,-1.5) {B : hypotonique};
\node at (0,-2.0) {entrée d'eau, hémolyse};
\end{scope}
% C : crénelure
\begin{scope}[xshift=10cm]
\draw[thick, popblue, fill=poppink!25] plot[smooth cycle] coordinates {(0.7,0) (0.5,0.35) (0.15,0.25) (0,0.65) (-0.2,0.3) (-0.6,0.4) (-0.45,0) (-0.6,-0.4) (-0.2,-0.3) (0,-0.65) (0.15,-0.25) (0.5,-0.35)};
\foreach \a in {30,90,150,210,270,330} \draw[->, popgreen, thick] (\a:0.95) -- (\a:1.55);
\node at (0,-1.5) {C : hypertonique};
\node at (0,-2.0) {sortie d'eau, crénelure};
\end{scope}
\end{tikzpicture}
\end{center}
\legende{Globule rouge dans les trois solutions ; les flèches vertes indiquent le sens du mouvement d'eau par osmose.}
\end{popfigure}
\end{enumerate}
\textbf{Barème indicatif (10 points) :} Q1 : 2 pts (états + noms précis des phénomènes) ; Q2 : 2 pts (classement + sens des flux d'eau) ; Q3 : 1 pt ; Q4 : 2,5 pts (hémolyse + danger expliqué) ; Q5 : 2,5 pts (trois schémas légendés avec flèches correctement orientées).
\end{corrige}

\begin{attention}[Points de vigilance — Exercice 9]
Le jury attend les termes exacts \cle{hémolyse} et \cle{crénelure}, et la justification du classement par le \emph{sens des mouvements d'eau}, pas seulement par l'aspect des cellules. Pour la question 4, il faut enchaîner logiquement : hypotonie $\rightarrow$ entrée d'eau $\rightarrow$ hémolyse $\rightarrow$ conséquences (anémie, toxicité rénale).
\end{attention}

\begin{corrige}[Exercice 10 — Type Bac : l'expérience de Pfeffer et la pression osmotique]
\begin{enumerate}
\item \textbf{Principe.} Une cellule dont la paroi poreuse est imprégnée de ferrocyanure de cuivre est remplie de solution de saccharose et plongée dans l'eau pure. L'eau entre par osmose, la pression interne s'élève et est mesurée par le manomètre : à l'équilibre, cette pression est la \cle{pression osmotique} $\pi$. La membrane de ferrocyanure de cuivre joue le rôle de membrane \textbf{semi-perméable} : perméable à l'eau, imperméable au saccharose.
\item \textbf{Graphe $\pi = f(C)$.}
\begin{popfigure}
\begin{center}
\begin{tikzpicture}
\begin{axis}[
width=0.8\linewidth, height=6cm,
xlabel={$C$ (mol/L)}, ylabel={$\pi$ (atm)},
xmin=0, xmax=0.28, ymin=0, ymax=7,
xtick={0,0.05,0.10,0.15,0.20,0.25},
grid=both, grid style={popline!40}
]
\addplot[popcurve, popblue, domain=0:0.27, samples=20] {24.04*x};
\addplot[only marks, mark=*, poporange, mark size=2pt] coordinates {(0.05,1.20) (0.10,2.40) (0.15,3.61) (0.20,4.81) (0.25,6.01)};
\end{axis}
\end{tikzpicture}
\end{center}
\legende{Pression osmotique en fonction de la concentration en saccharose : les points sont alignés sur une droite passant par l'origine.}
\end{popfigure}
Les points sont alignés sur une \textbf{droite passant par l'origine} : $\pi$ est proportionnelle à $C$, soit $\pi = k\,C$.
\item \textbf{Calcul du coefficient.} Avec le point $C = \num{0,10}$~mol/L, $\pi = \num{2,40}$~atm :
\[ \frac{\pi}{C} = \frac{\num{2,40}}{\num{0,10}} = \num{24,0}~\text{L.atm.mol}^{-1}. \]
(Vérification avec les autres points : $\num{3,61}/\num{0,15} \approx \num{24,1}$ ; $\num{4,81}/\num{0,20} \approx \num{24,1}$ ; $\num{6,01}/\num{0,25} \approx \num{24,0}$ : le rapport est constant.)
Or :
\[ RT = \num{0,082} \times 293 = \num{24,0}~\text{L.atm.mol}^{-1}. \]
On a bien $\dfrac{\pi}{C} \approx RT$, d'où la \textbf{loi de Van't Hoff} : $\pi = RTC$.
\item \textbf{Pression osmotique à 0,30~mol/L :}
\[ \pi = RTC = \num{0,082} \times 293 \times \num{0,30} \approx \num{7,2}~\text{atm}. \]
\item \textbf{Sens du mouvement d'eau.} D'après le tableau, la solution à 0,25~mol/L a une pression osmotique ($\num{6,01}$~atm) supérieure à celle de la solution à 0,10~mol/L ($\num{2,40}$~atm) : elle « attire » davantage l'eau. L'eau passe donc de la solution à \textbf{0,10~mol/L vers la solution à 0,25~mol/L}, c'est-à-dire du milieu le moins concentré vers le milieu le plus concentré, jusqu'à l'égalisation des concentrations.
\end{enumerate}
\textbf{Barème indicatif (10 points) :} Q1 : 2 pts ; Q2 : 2 pts (graphe + relation de proportionnalité) ; Q3 : 2,5 pts (calcul détaillé de $\pi/C$ et de $RT$, comparaison) ; Q4 : 1,5 pt ; Q5 : 2 pts.
\end{corrige}

\begin{attention}[Points de vigilance — Exercice 10]
Exigences du correcteur : unités portées dans le calcul de $RT$ ; conversion de la température en kelvins ($T = 20 + 273 = 293$~K) ; justification du sens du flux d'eau par la \emph{comparaison des pressions osmotiques} et non par une simple affirmation.
\end{attention}

\begin{corrige}[Exercice 11 — Type Bac : modes de transport et mise en évidence expérimentale de l'osmose]
\textbf{Partie A.}
\begin{enumerate}
\item \textbf{Interprétation des courbes.}
\begin{itemize}
\item \textbf{Soluté X :} droite passant par l'origine $\Rightarrow$ la vitesse d'absorption est proportionnelle à la concentration externe, sans saturation : il s'agit de la \cle{diffusion simple} (transport passif, passage direct à travers la membrane).
\item \textbf{Soluté Y :} courbe qui croît puis plafonne à partir de 4~mmol/L $\Rightarrow$ existence d'une vitesse maximale due à la saturation de transporteurs membranaires en nombre limité : il s'agit d'un transport \textbf{via des protéines} (diffusion facilitée ou transport actif).
\end{itemize}
\item \textbf{Test au cyanure.} Le cyanure bloque la respiration cellulaire, donc la production d'ATP.
\begin{itemize}
\item \textbf{Pour X (diffusion simple, passive) :} la vitesse d'absorption reste \textbf{inchangée}, car ce transport ne consomme pas d'énergie.
\item \textbf{Pour Y, si son transport est actif :} la vitesse d'absorption \textbf{chute fortement (voire s'annule)}, car le transport actif contre le gradient de concentration exige l'hydrolyse de l'ATP. Si au contraire la vitesse de Y est inchangée, c'est que Y entre par diffusion facilitée (passive malgré les transporteurs).
\end{itemize}
\end{enumerate}

\textbf{Partie B.}
\begin{enumerate}
\item[3.] \textbf{Protocole expérimental.}
\textbf{Méthode : Mise en évidence de l'osmose}
\begin{enumerate}
\item Remplir le tube à dialyse (cellophane, membrane semi-perméable) de solution concentrée de glucose ; le fermer à une extrémité avec du fil.
\item Obturer l'autre extrémité avec le bouchon traversé par le capillaire en verre ; repérer le niveau initial du liquide dans le capillaire.
\item Fixer l'ensemble sur le support et le plonger dans un bécher d'eau distillée.
\item \textbf{Dispositif témoin :} monter un second tube à dialyse rempli d'\emph{eau distillée} (au lieu de la solution de glucose), plongé dans un bécher d'eau distillée, dans les mêmes conditions.
\item Suivre le niveau du liquide dans les deux capillaires pendant 30 à 60 minutes.
\end{enumerate}
\item[4.] \textbf{Résultats attendus et interprétation.}
\begin{itemize}
\item \textbf{Dispositif expérimental :} le niveau du liquide \textbf{monte} dans le capillaire. Interprétation : la cellophane, semi-perméable, laisse passer l'eau mais pas le glucose ; l'eau distillée (hypotonique) pénètre par \cle{osmose} dans le tube où la solution est concentrée, augmentant le volume interne.
\item \textbf{Témoin :} le niveau reste \textbf{stable}, car il n'existe aucun gradient de concentration de part et d'autre de la membrane. La montée observée dans le dispositif expérimental est donc bien due à l'osmose et non à un artefact du montage.
\end{itemize}
\item[5.] \textbf{Mise en évidence de la dialyse.} Remplir le tube à dialyse d'une solution contenant de l'\textbf{amidon} et le plonger dans un bécher d'eau distillée additionnée d'\textbf{eau iodée} (ou placer du glucose dans le tube et tester le bécher). L'eau iodée (petites molécules) traverse la cellophane et entre dans le tube : l'amidon vire au \textbf{bleu-noir} en présence de Lugol/d'iode, alors que l'amidon (grosses molécules) ne sort pas. Le passage sélectif des petites molécules de soluté à travers la membrane, du milieu concentré vers le milieu dilué, met en évidence la \cle{dialyse}.
\end{enumerate}
\textbf{Barème indicatif (10 points) :} Partie A : Q1 : 2 pts, Q2 : 2 pts. Partie B : Q3 : 3 pts (montage + témoin indispensable), Q4 : 2 pts, Q5 : 1 pt.
\end{corrige}

\begin{attention}[Points de vigilance — Exercice 11]
Le \textbf{dispositif témoin} est exigé dans toute démarche expérimentale : son absence coûte des points. Pour le test au cyanure, il faut explicitement écrire que le cyanure inhibe la respiration cellulaire, donc la production d'ATP, et que seul le transport \emph{actif} en dépend. Dans l'interprétation de l'osmose, toujours préciser la nature \emph{semi-perméable} de la cellophane.
\end{attention}

\newpage

\coursec{Fiche bilan}

\begin{popfigure}
\begin{tikzpicture}[mindmap, grow cyclic, every node/.style={concept, font=\small}, text width=2.6cm,
  root concept/.append style={concept color=popblue, font=\small\bfseries, text width=3.2cm},
  level 1/.append style={level distance=3.6cm, sibling angle=60, font=\footnotesize},
  level 2/.append style={level distance=2.4cm, sibling angle=45, font=\scriptsize, text width=2.2cm}]
\node[root concept]{Échanges cellule--milieu}
  child[concept color=popgreen]{ node {Échanges d'eau}
    child[concept color=popgreen!60]{ node {Turgescence / plasmolyse (végétal)} }
    child[concept color=popgreen!60]{ node {Hémolyse / crénelure (animal)} }
  }
  child[concept color=poporange]{ node {Osmose et dialyse}
    child[concept color=poporange!60]{ node {Dutrochet, Pfeiffer} }
    child[concept color=poporange!60]{ node {Membrane hémiperméable} }
  }
  child[concept color=poppurple]{ node {Transport passif}
    child[concept color=poppurple!60]{ node {Diffusion simple} }
    child[concept color=poppurple!60]{ node {Diffusion facilitée (saturation)} }
  }
  child[concept color=poppink]{ node {Transport actif}
    child[concept color=poppink!60]{ node {ATP + transporteur} }
    child[concept color=poppink!60]{ node {Contre le gradient} }
  }
  child[concept color=popgold]{ node {Échanges de particules}
    child[concept color=popgold!60]{ node {Endocytose : phago-, pinocytose} }
    child[concept color=popgold!60]{ node {Exocytose} }
  }
  child[concept color=popmuted]{ node {Importance et applications}
    child[concept color=popmuted!70]{ node {Nutrition, défense, dialyse rénale, conservation} }
  };
\end{tikzpicture}
\legende{Carte mentale de la leçon : les six grandes familles de savoirs sur les échanges entre la cellule et le milieu ambiant.}
\end{popfigure}

\begin{bilan}
\textbf{Définitions clés}
\begin{itemize}
  \item \cle{Osmose} : diffusion du solvant (eau) à travers une membrane hémiperméable, du milieu \cle{hypotonique} (moins concentré) vers le milieu \cle{hypertonique} (plus concentré).
  \item \cle{Dialyse} : diffusion des petites molécules dissoutes (cristalloïdes) à travers une membrane perméable, du milieu concentré vers le milieu dilué ; les colloïdes ne traversent pas.
  \item \cle{Turgescence} : gonflement de la cellule végétale en milieu hypotonique (vacuole pleine, membrane plaquée contre la paroi) ; \cle{plasmolyse} : décollement du protoplasme en milieu hypertonique.
  \item \cle{Hémolyse} : éclatement du globule rouge en milieu hypotonique ; \cle{crénelure} : rétraction et aspect festonné en milieu hypertonique.
  \item \cle{Diffusion simple} : passage direct de molécules (gaz, liposolubles) à travers la bicouche lipidique, selon le gradient, sans énergie.
  \item \cle{Diffusion facilitée} : passage via une protéine transporteur (canal ou permease), selon le gradient, sans ATP, avec \textbf{saturation} des transporteurs.
  \item \cle{Transport actif} : passage \textbf{contre le gradient} de concentration, nécessitant un transporteur et de l'\cle{ATP}.
  \item \cle{Endocytose} : incorporation de particules par invagination membranaire (\cle{phagocytose} : solides ; \cle{pinocytose} : liquides) ; \cle{exocytose} : expulsion par fusion d'une vésicule avec la membrane.
\end{itemize}

\textbf{Expériences et résultats à connaître par c\oe ur}
\begin{center}
\begin{tabularx}{\linewidth}{|l|X|X|}
\hline
\rowcolor{popblueL} \textbf{Expérience} & \textbf{Protocole} & \textbf{Résultat / conclusion} \\
\hline
Dutrochet & Vessie de porc séparant eau distillée et solution concentrée, reliée à un tube vertical & Montée du liquide dans le tube : l'eau passe vers le milieu concentré $\Rightarrow$ osmose, pression osmotique \\
\hline
Pfeiffer & Cellule osmométrique à membrane de ferrocyanure de cuivre + manomètre & Mesure de la pression osmotique : elle augmente avec la concentration du soluté \\
\hline
Diffusion facilitée & Vitesse d'entrée d'un soluté en fonction de sa concentration extracellulaire & Courbe en \textbf{hyperbole avec plateau} (saturation des transporteurs), contrairement à la droite de la diffusion simple \\
\hline
GR en solutions de NaCl & \SI{0,9}{\percent} (isotonique), \SI{0,3}{\percent} (hypotonique), \SI{3}{\percent} (hypertonique) & GR normaux / hémolyse / crénelure : le sang est isotonique à \SI{0,9}{\percent} NaCl \\
\hline
\end{tabularx}
\end{center}

\textbf{Méthodes express}
\begin{itemize}
  \item Sens du flux d'eau : toujours du milieu \emph{le moins concentré} vers \emph{le plus concentré} en solutés.
  \item Cellule végétale : paroi = protection contre l'éclatement ; cellule animale : fragile (hémolyse).
  \item Distinguer passif/actif : présence d'ATP et sens par rapport au gradient (contre le gradient $\Rightarrow$ actif).
  \item Distinguer diffusion simple/facilitée : cinétique linéaire (simple) ou saturable en plateau (facilitée).
  \item Protocole osmose : membrane hémiperméable + deux milieux de concentrations différentes + repérage du niveau de liquide.
\end{itemize}
\end{bilan}

\begin{aretenir}[Checklist avant le Bac]
\begin{itemize}
  \item[$\square$] Je sais définir l'osmose et la dialyse et les distinguer.
  \item[$\square$] Je prédis le sens du mouvement de l'eau selon la tonicité des milieux.
  \item[$\square$] Je décris et j'identifie turgescence, plasmolyse, hémolyse et crénelure sur une préparation ou une photographie.
  \item[$\square$] J'interprète l'expérience de Dutrochet (montée du liquide, pression osmotique).
  \item[$\square$] J'interprète l'expérience de Pfeiffer (mesure de la pression osmotique).
  \item[$\square$] Je sais proposer un protocole expérimental mettant en évidence l'osmose.
  \item[$\square$] Je distingue diffusion simple et diffusion facilitée grâce à l'allure du graphe (droite vs plateau de saturation).
  \item[$\square$] Je justifie le caractère actif d'un transport (contre le gradient, besoin d'ATP).
  \item[$\square$] Je décris les étapes de la phagocytose, de la pinocytose et de l'exocytose.
  \item[$\square$] Je cite des applications : perfusions isotoniques, dialyse rénale, conservation des aliments (salage, sucrage), absorption racinaire.
  \item[$\square$] Je rédige mes interprétations avec le vocabulaire exact : hypotonique, isotonique, hypertonique, membrane hémiperméable.
  \item[$\square$] Je sais que \SI{0,9}{\percent} NaCl est isotonique au sang (sérum physiologique).
\end{itemize}
\end{aretenir}

\findecours

\end{document}
